% Môn: Vật lý
% Lớp: 12
% Chương: Chương IV. Vật lí hạt nhân
% Bài: L12C4 Bài 22. Phản ứng hạt nhân và năng lượng liên kết
% Số câu: 204
% App đọc file này trực tiếp — sửa rồi Commit trên GitHub, bấm Đồng bộ GitHub .tex

% L12C4 Bài 22. Phản ứng hạt nhân và năng lượng liên kết

% ===== Câu 1 =====
% ID: L12C4B22-01-DS
% Mức: TH
\dangbt{Bài toán tổng hợp về phản ứng hạt nhân, phân hạch và nhiệt hạch}
\begin{ex}
% ID: L12C4B22-01-DS
% Mức: TH
Xét bài toán tổng hợp về phản ứng hạt nhân, phân hạch và nhiệt hạch (tình huống 1). Hãy xác định tính đúng, sai của bốn phát biểu sau.
\choiceTF
{\True Phân hạch hạt nhân nặng và nhiệt hạch hạt nhân nhẹ đều có thể tỏa năng lượng do sản phẩm bền hơn về năng lượng liên kết riêng.}
{Mọi phản ứng hạt nhân đều tự phát và tỏa cùng một năng lượng.}
{\True Khi giải cần đọc đúng dữ kiện, dùng hệ đơn vị thống nhất và kiểm tra điều kiện áp dụng.}
{Có thể bỏ qua dữ kiện, đơn vị và ý nghĩa vật lí của kết quả.}
\loigiai{
\begin{itemchoice}
\item (Đúng) Phân hạch hạt nhân nặng và nhiệt hạch hạt nhân nhẹ đều có thể tỏa năng lượng do sản phẩm bền hơn về năng lượng liên kết riêng. (Vì): Phân hạch hạt nhân nặng và nhiệt hạch hạt nhân nhẹ đều có thể tỏa năng lượng do sản phẩm bền hơn về năng lượng liên kết riêng. b) Sai: Mọi phản ứng hạt nhân đều tự phát và tỏa cùng một năng lượng. c) Đúng: Khi giải cần đọc đúng dữ kiện, dùng hệ đơn vị thống nhất và kiểm tra điều kiện áp dụng. d) Sai: Có thể bỏ qua dữ kiện, đơn vị và ý nghĩa vật lí của kết quả. Kiến thức cốt lõi: Phân hạch hạt nhân nặng và nhiệt hạch hạt nhân nhẹ đều có thể tỏa năng lượng do sản phẩm bền hơn về năng lượng liên kết riêng
\item (Sai) Mọi phản ứng hạt nhân đều tự phát và tỏa cùng một năng lượng.
\item (Đúng) Khi giải cần đọc đúng dữ kiện, dùng hệ đơn vị thống nhất và kiểm tra điều kiện áp dụng.
\item (Sai) Có thể bỏ qua dữ kiện, đơn vị và ý nghĩa vật lí của kết quả.
\end{itemchoice}
}
\end{ex}

% ===== Câu 2 =====
% ID: L12C4B22-01-TL
% Mức: TH
\begin{bt}
% ID: L12C4B22-01-TL
% Mức: TH
Trình bày kiến thức cốt lõi của bài toán tổng hợp về phản ứng hạt nhân, phân hạch và nhiệt hạch và nêu điều kiện áp dụng.
\loigiai{
Cần nêu được: Phân hạch hạt nhân nặng và nhiệt hạch hạt nhân nhẹ đều có thể tỏa năng lượng do sản phẩm bền hơn về năng lượng liên kết riêng. Khi vận dụng phải xác định đúng dữ kiện, điều kiện áp dụng, hệ đơn vị và kiểm tra tính hợp lí của kết quả. Nhận định “Mọi phản ứng hạt nhân đều tự phát và tỏa cùng một năng lượng.” sai vì trái với kiến thức trên.
}
\end{bt}

% ===== Câu 3 =====
% ID: L12C4B22-01-TLN
% Mức: TH
\begin{ex}
% ID: L12C4B22-01-TLN
% Mức: TH
Trong bốn phát biểu sau về bài toán tổng hợp về phản ứng hạt nhân, phân hạch và nhiệt hạch, có bao nhiêu phát biểu đúng? Chỉ ghi một số nguyên. (1) Phân hạch hạt nhân nặng và nhiệt hạch hạt nhân nhẹ đều có thể tỏa năng lượng do sản phẩm bền hơn về năng lượng liên kết riêng. (2) Mọi phản ứng hạt nhân đều tự phát và tỏa cùng một năng lượng. (3) Cần kiểm tra điều kiện áp dụng trước khi tính toán. (4) Có thể bỏ qua đơn vị của kết quả.
\shortans{2}
\loigiai{
Đối chiếu từng phát biểu với kiến thức cốt lõi: Phân hạch hạt nhân nặng và nhiệt hạch hạt nhân nhẹ đều có thể tỏa năng lượng do sản phẩm bền hơn về năng lượng liên kết riêng. Có 2 phát biểu đúng.
}
\end{ex}

% ===== Câu 4 =====
% ID: L12C4B22-01-TN
% Mức: NB
\begin{ex}
% ID: L12C4B22-01-TN
% Mức: NB
Phát biểu nào sau đây đúng về bài toán tổng hợp về phản ứng hạt nhân, phân hạch và nhiệt hạch?
\choice
{\True Phân hạch hạt nhân nặng và nhiệt hạch hạt nhân nhẹ đều có thể tỏa năng lượng do sản phẩm bền hơn về năng lượng liên kết riêng.}
{Mọi phản ứng hạt nhân đều tự phát và tỏa cùng một năng lượng.}
{Có thể bỏ qua điều kiện áp dụng và đơn vị trong mọi trường hợp.}
{Kết quả của bài toán không phụ thuộc dữ kiện đã cho.}
\loigiai{
Phân hạch hạt nhân nặng và nhiệt hạch hạt nhân nhẹ đều có thể tỏa năng lượng do sản phẩm bền hơn về năng lượng liên kết riêng. Vì vậy phương án $A$ đúng; các phương án còn lại trái với kiến thức hoặc điều kiện áp dụng.
}
\end{ex}

% ===== Câu 5 =====
% ID: L12C4B22-02-DS
% Mức: TH
\begin{ex}
% ID: L12C4B22-02-DS
% Mức: TH
Xét bài toán tổng hợp về phản ứng hạt nhân, phân hạch và nhiệt hạch (tình huống 2). Hãy xác định tính đúng, sai của bốn phát biểu sau.
\choiceTF
{Mọi phản ứng hạt nhân đều tự phát và tỏa cùng một năng lượng.}
{\True Phân hạch hạt nhân nặng và nhiệt hạch hạt nhân nhẹ đều có thể tỏa năng lượng do sản phẩm bền hơn về năng lượng liên kết riêng.}
{Có thể bỏ qua dữ kiện, đơn vị và ý nghĩa vật lí của kết quả.}
{\True Khi giải cần đọc đúng dữ kiện, dùng hệ đơn vị thống nhất và kiểm tra điều kiện áp dụng.}
\loigiai{
\begin{itemchoice}
\item (Sai) Mọi phản ứng hạt nhân đều tự phát và tỏa cùng một năng lượng. (Vì): Mọi phản ứng hạt nhân đều tự phát và tỏa cùng một năng lượng. b) Đúng: Phân hạch hạt nhân nặng và nhiệt hạch hạt nhân nhẹ đều có thể tỏa năng lượng do sản phẩm bền hơn về năng lượng liên kết riêng. c) Sai: Có thể bỏ qua dữ kiện, đơn vị và ý nghĩa vật lí của kết quả. d) Đúng: Khi giải cần đọc đúng dữ kiện, dùng hệ đơn vị thống nhất và kiểm tra điều kiện áp dụng. Kiến thức cốt lõi: Phân hạch hạt nhân nặng và nhiệt hạch hạt nhân nhẹ đều có thể tỏa năng lượng do sản phẩm bền hơn về năng lượng liên kết riêng
\item (Đúng) Phân hạch hạt nhân nặng và nhiệt hạch hạt nhân nhẹ đều có thể tỏa năng lượng do sản phẩm bền hơn về năng lượng liên kết riêng.
\item (Sai) Có thể bỏ qua dữ kiện, đơn vị và ý nghĩa vật lí của kết quả.
\item (Đúng) Khi giải cần đọc đúng dữ kiện, dùng hệ đơn vị thống nhất và kiểm tra điều kiện áp dụng.
\end{itemchoice}
}
\end{ex}

% ===== Câu 6 =====
% ID: L12C4B22-02-TL
% Mức: TH
\begin{bt}
% ID: L12C4B22-02-TL
% Mức: TH
Giải thích vì sao nhận định sau đúng: “Phân hạch hạt nhân nặng và nhiệt hạch hạt nhân nhẹ đều có thể tỏa năng lượng do sản phẩm bền hơn về năng lượng liên kết riêng.”
\loigiai{
Cần nêu được: Phân hạch hạt nhân nặng và nhiệt hạch hạt nhân nhẹ đều có thể tỏa năng lượng do sản phẩm bền hơn về năng lượng liên kết riêng. Khi vận dụng phải xác định đúng dữ kiện, điều kiện áp dụng, hệ đơn vị và kiểm tra tính hợp lí của kết quả. Nhận định “Mọi phản ứng hạt nhân đều tự phát và tỏa cùng một năng lượng.” sai vì trái với kiến thức trên.
}
\end{bt}

% ===== Câu 7 =====
% ID: L12C4B22-02-TLN
% Mức: TH
\begin{ex}
% ID: L12C4B22-02-TLN
% Mức: TH
Trong bốn phát biểu sau về bài toán tổng hợp về phản ứng hạt nhân, phân hạch và nhiệt hạch, có bao nhiêu phát biểu đúng? Chỉ ghi một số nguyên. (1) Mọi phản ứng hạt nhân đều tự phát và tỏa cùng một năng lượng. (2) Phân hạch hạt nhân nặng và nhiệt hạch hạt nhân nhẹ đều có thể tỏa năng lượng do sản phẩm bền hơn về năng lượng liên kết riêng. (3) Kết quả phải phù hợp ý nghĩa vật lí. (4) Mọi đại lượng đều có cùng bản chất.
\shortans{2}
\loigiai{
Đối chiếu từng phát biểu với kiến thức cốt lõi: Phân hạch hạt nhân nặng và nhiệt hạch hạt nhân nhẹ đều có thể tỏa năng lượng do sản phẩm bền hơn về năng lượng liên kết riêng. Có 2 phát biểu đúng.
}
\end{ex}

% ===== Câu 8 =====
% ID: L12C4B22-02-TN
% Mức: NB
\begin{ex}
% ID: L12C4B22-02-TN
% Mức: NB
Nhận định nào phù hợp với kiến thức Vật lí về bài toán tổng hợp về phản ứng hạt nhân, phân hạch và nhiệt hạch?
\choice
{Mọi phản ứng hạt nhân đều tự phát và tỏa cùng một năng lượng.}
{\True Phân hạch hạt nhân nặng và nhiệt hạch hạt nhân nhẹ đều có thể tỏa năng lượng do sản phẩm bền hơn về năng lượng liên kết riêng.}
{Có thể bỏ qua điều kiện áp dụng và đơn vị trong mọi trường hợp.}
{Kết quả của bài toán không phụ thuộc dữ kiện đã cho.}
\loigiai{
Phân hạch hạt nhân nặng và nhiệt hạch hạt nhân nhẹ đều có thể tỏa năng lượng do sản phẩm bền hơn về năng lượng liên kết riêng. Vì vậy phương án $B$ đúng; các phương án còn lại trái với kiến thức hoặc điều kiện áp dụng.
}
\end{ex}

% ===== Câu 9 =====
% ID: L12C4B22-03-DS
% Mức: VD
\begin{ex}
% ID: L12C4B22-03-DS
% Mức: VD
Xét bài toán tổng hợp về phản ứng hạt nhân, phân hạch và nhiệt hạch (tình huống 3). Hãy xác định tính đúng, sai của bốn phát biểu sau.
\choiceTF
{\True Phân hạch hạt nhân nặng và nhiệt hạch hạt nhân nhẹ đều có thể tỏa năng lượng do sản phẩm bền hơn về năng lượng liên kết riêng.}
{Mọi phản ứng hạt nhân đều tự phát và tỏa cùng một năng lượng.}
{\True Khi giải cần đọc đúng dữ kiện, dùng hệ đơn vị thống nhất và kiểm tra điều kiện áp dụng.}
{Có thể bỏ qua dữ kiện, đơn vị và ý nghĩa vật lí của kết quả.}
\loigiai{
\begin{itemchoice}
\item (Đúng) Phân hạch hạt nhân nặng và nhiệt hạch hạt nhân nhẹ đều có thể tỏa năng lượng do sản phẩm bền hơn về năng lượng liên kết riêng. (Vì): Phân hạch hạt nhân nặng và nhiệt hạch hạt nhân nhẹ đều có thể tỏa năng lượng do sản phẩm bền hơn về năng lượng liên kết riêng. b) Sai: Mọi phản ứng hạt nhân đều tự phát và tỏa cùng một năng lượng. c) Đúng: Khi giải cần đọc đúng dữ kiện, dùng hệ đơn vị thống nhất và kiểm tra điều kiện áp dụng. d) Sai: Có thể bỏ qua dữ kiện, đơn vị và ý nghĩa vật lí của kết quả. Kiến thức cốt lõi: Phân hạch hạt nhân nặng và nhiệt hạch hạt nhân nhẹ đều có thể tỏa năng lượng do sản phẩm bền hơn về năng lượng liên kết riêng
\item (Sai) Mọi phản ứng hạt nhân đều tự phát và tỏa cùng một năng lượng.
\item (Đúng) Khi giải cần đọc đúng dữ kiện, dùng hệ đơn vị thống nhất và kiểm tra điều kiện áp dụng.
\item (Sai) Có thể bỏ qua dữ kiện, đơn vị và ý nghĩa vật lí của kết quả.
\end{itemchoice}
}
\end{ex}

% ===== Câu 10 =====
% ID: L12C4B22-03-TL
% Mức: VD
\begin{bt}
% ID: L12C4B22-03-TL
% Mức: VD
Phân tích sai lầm trong nhận định: “Mọi phản ứng hạt nhân đều tự phát và tỏa cùng một năng lượng.”
\loigiai{
Cần nêu được: Phân hạch hạt nhân nặng và nhiệt hạch hạt nhân nhẹ đều có thể tỏa năng lượng do sản phẩm bền hơn về năng lượng liên kết riêng. Khi vận dụng phải xác định đúng dữ kiện, điều kiện áp dụng, hệ đơn vị và kiểm tra tính hợp lí của kết quả. Nhận định “Mọi phản ứng hạt nhân đều tự phát và tỏa cùng một năng lượng.” sai vì trái với kiến thức trên.
}
\end{bt}

% ===== Câu 11 =====
% ID: L12C4B22-03-TLN
% Mức: TH
\begin{ex}
% ID: L12C4B22-03-TLN
% Mức: TH
Trong bốn phát biểu sau về bài toán tổng hợp về phản ứng hạt nhân, phân hạch và nhiệt hạch, có bao nhiêu phát biểu đúng? Chỉ ghi một số nguyên. (1) Cần đổi các đại lượng về hệ đơn vị thống nhất. (2) Phân hạch hạt nhân nặng và nhiệt hạch hạt nhân nhẹ đều có thể tỏa năng lượng do sản phẩm bền hơn về năng lượng liên kết riêng. (3) Mọi phản ứng hạt nhân đều tự phát và tỏa cùng một năng lượng. (4) Không cần kiểm tra dữ kiện.
\shortans{2}
\loigiai{
Đối chiếu từng phát biểu với kiến thức cốt lõi: Phân hạch hạt nhân nặng và nhiệt hạch hạt nhân nhẹ đều có thể tỏa năng lượng do sản phẩm bền hơn về năng lượng liên kết riêng. Có 2 phát biểu đúng.
}
\end{ex}

% ===== Câu 12 =====
% ID: L12C4B22-03-TN
% Mức: NB
\begin{ex}
% ID: L12C4B22-03-TN
% Mức: NB
Chọn kết luận chính xác về bài toán tổng hợp về phản ứng hạt nhân, phân hạch và nhiệt hạch?
\choice
{Mọi phản ứng hạt nhân đều tự phát và tỏa cùng một năng lượng.}
{Có thể bỏ qua điều kiện áp dụng và đơn vị trong mọi trường hợp.}
{\True Phân hạch hạt nhân nặng và nhiệt hạch hạt nhân nhẹ đều có thể tỏa năng lượng do sản phẩm bền hơn về năng lượng liên kết riêng.}
{Kết quả của bài toán không phụ thuộc dữ kiện đã cho.}
\loigiai{
Phân hạch hạt nhân nặng và nhiệt hạch hạt nhân nhẹ đều có thể tỏa năng lượng do sản phẩm bền hơn về năng lượng liên kết riêng. Vì vậy phương án $C$ đúng; các phương án còn lại trái với kiến thức hoặc điều kiện áp dụng.
}
\end{ex}

% ===== Câu 13 =====
% ID: L12C4B22-04-DS
% Mức: VD
\begin{ex}
% ID: L12C4B22-04-DS
% Mức: VD
Xét bài toán tổng hợp về phản ứng hạt nhân, phân hạch và nhiệt hạch (tình huống 4). Hãy xác định tính đúng, sai của bốn phát biểu sau.
\choiceTF
{Mọi phản ứng hạt nhân đều tự phát và tỏa cùng một năng lượng.}
{\True Phân hạch hạt nhân nặng và nhiệt hạch hạt nhân nhẹ đều có thể tỏa năng lượng do sản phẩm bền hơn về năng lượng liên kết riêng.}
{Có thể bỏ qua dữ kiện, đơn vị và ý nghĩa vật lí của kết quả.}
{\True Khi giải cần đọc đúng dữ kiện, dùng hệ đơn vị thống nhất và kiểm tra điều kiện áp dụng.}
\loigiai{
\begin{itemchoice}
\item (Sai) Mọi phản ứng hạt nhân đều tự phát và tỏa cùng một năng lượng. (Vì): Mọi phản ứng hạt nhân đều tự phát và tỏa cùng một năng lượng. b) Đúng: Phân hạch hạt nhân nặng và nhiệt hạch hạt nhân nhẹ đều có thể tỏa năng lượng do sản phẩm bền hơn về năng lượng liên kết riêng. c) Sai: Có thể bỏ qua dữ kiện, đơn vị và ý nghĩa vật lí của kết quả. d) Đúng: Khi giải cần đọc đúng dữ kiện, dùng hệ đơn vị thống nhất và kiểm tra điều kiện áp dụng. Kiến thức cốt lõi: Phân hạch hạt nhân nặng và nhiệt hạch hạt nhân nhẹ đều có thể tỏa năng lượng do sản phẩm bền hơn về năng lượng liên kết riêng
\item (Đúng) Phân hạch hạt nhân nặng và nhiệt hạch hạt nhân nhẹ đều có thể tỏa năng lượng do sản phẩm bền hơn về năng lượng liên kết riêng.
\item (Sai) Có thể bỏ qua dữ kiện, đơn vị và ý nghĩa vật lí của kết quả.
\item (Đúng) Khi giải cần đọc đúng dữ kiện, dùng hệ đơn vị thống nhất và kiểm tra điều kiện áp dụng.
\end{itemchoice}
}
\end{ex}

% ===== Câu 14 =====
% ID: L12C4B22-04-TL
% Mức: VD
\begin{bt}
% ID: L12C4B22-04-TL
% Mức: VD
Đề xuất các bước giải một bài toán thuộc dạng bài toán tổng hợp về phản ứng hạt nhân, phân hạch và nhiệt hạch.
\loigiai{
Cần nêu được: Phân hạch hạt nhân nặng và nhiệt hạch hạt nhân nhẹ đều có thể tỏa năng lượng do sản phẩm bền hơn về năng lượng liên kết riêng. Khi vận dụng phải xác định đúng dữ kiện, điều kiện áp dụng, hệ đơn vị và kiểm tra tính hợp lí của kết quả. Nhận định “Mọi phản ứng hạt nhân đều tự phát và tỏa cùng một năng lượng.” sai vì trái với kiến thức trên.
}
\end{bt}

% ===== Câu 15 =====
% ID: L12C4B22-04-TLN
% Mức: VD
\begin{ex}
% ID: L12C4B22-04-TLN
% Mức: VD
Trong bốn phát biểu sau về bài toán tổng hợp về phản ứng hạt nhân, phân hạch và nhiệt hạch, có bao nhiêu phát biểu đúng? Chỉ ghi một số nguyên. (1) Phân hạch hạt nhân nặng và nhiệt hạch hạt nhân nhẹ đều có thể tỏa năng lượng do sản phẩm bền hơn về năng lượng liên kết riêng. (2) Cần đối chiếu kết quả với điều kiện bài toán. (3) Mọi phản ứng hạt nhân đều tự phát và tỏa cùng một năng lượng. (4) Mọi trường hợp đều cho cùng một kết quả.
\shortans{2}
\loigiai{
Đối chiếu từng phát biểu với kiến thức cốt lõi: Phân hạch hạt nhân nặng và nhiệt hạch hạt nhân nhẹ đều có thể tỏa năng lượng do sản phẩm bền hơn về năng lượng liên kết riêng. Có 2 phát biểu đúng.
}
\end{ex}

% ===== Câu 16 =====
% ID: L12C4B22-04-TN
% Mức: TH
\begin{ex}
% ID: L12C4B22-04-TN
% Mức: TH
Nội dung nào mô tả đúng nhất về bài toán tổng hợp về phản ứng hạt nhân, phân hạch và nhiệt hạch?
\choice
{Mọi phản ứng hạt nhân đều tự phát và tỏa cùng một năng lượng.}
{Có thể bỏ qua điều kiện áp dụng và đơn vị trong mọi trường hợp.}
{Kết quả của bài toán không phụ thuộc dữ kiện đã cho.}
{\True Phân hạch hạt nhân nặng và nhiệt hạch hạt nhân nhẹ đều có thể tỏa năng lượng do sản phẩm bền hơn về năng lượng liên kết riêng.}
\loigiai{
Phân hạch hạt nhân nặng và nhiệt hạch hạt nhân nhẹ đều có thể tỏa năng lượng do sản phẩm bền hơn về năng lượng liên kết riêng. Vì vậy phương án $D$ đúng; các phương án còn lại trái với kiến thức hoặc điều kiện áp dụng.
}
\end{ex}

% ===== Câu 17 =====
% ID: L12C4B22-05-DS
% Mức: VD
\begin{ex}
% ID: L12C4B22-05-DS
% Mức: VD
Xét bài toán tổng hợp về phản ứng hạt nhân, phân hạch và nhiệt hạch (tình huống 5). Hãy xác định tính đúng, sai của bốn phát biểu sau.
\choiceTF
{\True Phân hạch hạt nhân nặng và nhiệt hạch hạt nhân nhẹ đều có thể tỏa năng lượng do sản phẩm bền hơn về năng lượng liên kết riêng.}
{Có thể bỏ qua dữ kiện, đơn vị và ý nghĩa vật lí của kết quả.}
{Mọi phản ứng hạt nhân đều tự phát và tỏa cùng một năng lượng.}
{Có thể bỏ qua dữ kiện, đơn vị và ý nghĩa vật lí của kết quả.}
\loigiai{
\begin{itemchoice}
\item (Đúng) Phân hạch hạt nhân nặng và nhiệt hạch hạt nhân nhẹ đều có thể tỏa năng lượng do sản phẩm bền hơn về năng lượng liên kết riêng. (Vì): Phân hạch hạt nhân nặng và nhiệt hạch hạt nhân nhẹ đều có thể tỏa năng lượng do sản phẩm bền hơn về năng lượng liên kết riêng. b) Sai: Có thể bỏ qua dữ kiện, đơn vị và ý nghĩa vật lí của kết quả. c) Sai: Mọi phản ứng hạt nhân đều tự phát và tỏa cùng một năng lượng. d) Sai: Có thể bỏ qua dữ kiện, đơn vị và ý nghĩa vật lí của kết quả. Kiến thức cốt lõi: Phân hạch hạt nhân nặng và nhiệt hạch hạt nhân nhẹ đều có thể tỏa năng lượng do sản phẩm bền hơn về năng lượng liên kết riêng
\item (Sai) Có thể bỏ qua dữ kiện, đơn vị và ý nghĩa vật lí của kết quả.
\item (Sai) Mọi phản ứng hạt nhân đều tự phát và tỏa cùng một năng lượng.
\item (Sai) Có thể bỏ qua dữ kiện, đơn vị và ý nghĩa vật lí của kết quả.
\end{itemchoice}
}
\end{ex}

% ===== Câu 18 =====
% ID: L12C4B22-05-TL
% Mức: VD
\begin{bt}
% ID: L12C4B22-05-TL
% Mức: VD
Nêu một tình huống thực tiễn liên quan đến bài toán tổng hợp về phản ứng hạt nhân, phân hạch và nhiệt hạch và giải thích bằng kiến thức Vật lí.
\loigiai{
Cần nêu được: Phân hạch hạt nhân nặng và nhiệt hạch hạt nhân nhẹ đều có thể tỏa năng lượng do sản phẩm bền hơn về năng lượng liên kết riêng. Khi vận dụng phải xác định đúng dữ kiện, điều kiện áp dụng, hệ đơn vị và kiểm tra tính hợp lí của kết quả. Nhận định “Mọi phản ứng hạt nhân đều tự phát và tỏa cùng một năng lượng.” sai vì trái với kiến thức trên.
}
\end{bt}

% ===== Câu 19 =====
% ID: L12C4B22-05-TLN
% Mức: VD
\begin{ex}
% ID: L12C4B22-05-TLN
% Mức: VD
Trong bốn phát biểu sau về bài toán tổng hợp về phản ứng hạt nhân, phân hạch và nhiệt hạch, có bao nhiêu phát biểu đúng? Chỉ ghi một số nguyên. (1) Mọi phản ứng hạt nhân đều tự phát và tỏa cùng một năng lượng. (2) Cần đọc đúng dữ kiện và giả thiết. (3) Phân hạch hạt nhân nặng và nhiệt hạch hạt nhân nhẹ đều có thể tỏa năng lượng do sản phẩm bền hơn về năng lượng liên kết riêng. (4) Có thể thay số trước khi chọn công thức.
\shortans{2}
\loigiai{
Đối chiếu từng phát biểu với kiến thức cốt lõi: Phân hạch hạt nhân nặng và nhiệt hạch hạt nhân nhẹ đều có thể tỏa năng lượng do sản phẩm bền hơn về năng lượng liên kết riêng. Có 2 phát biểu đúng.
}
\end{ex}

% ===== Câu 20 =====
% ID: L12C4B22-05-TN
% Mức: TH
\begin{ex}
% ID: L12C4B22-05-TN
% Mức: TH
Khi phân tích, phát biểu nào đúng về bài toán tổng hợp về phản ứng hạt nhân, phân hạch và nhiệt hạch?
\choice
{\True Phân hạch hạt nhân nặng và nhiệt hạch hạt nhân nhẹ đều có thể tỏa năng lượng do sản phẩm bền hơn về năng lượng liên kết riêng.}
{Mọi phản ứng hạt nhân đều tự phát và tỏa cùng một năng lượng.}
{Có thể bỏ qua điều kiện áp dụng và đơn vị trong mọi trường hợp.}
{Kết quả của bài toán không phụ thuộc dữ kiện đã cho.}
\loigiai{
Phân hạch hạt nhân nặng và nhiệt hạch hạt nhân nhẹ đều có thể tỏa năng lượng do sản phẩm bền hơn về năng lượng liên kết riêng. Vì vậy phương án $A$ đúng; các phương án còn lại trái với kiến thức hoặc điều kiện áp dụng.
}
\end{ex}

% ===== Câu 21 =====
% ID: L12C4B22-06-DS
% Mức: TH
\dangbt{Hoàn thành phương trình phản ứng hạt nhân}
\begin{ex}
% ID: L12C4B22-06-DS
% Mức: TH
Xét hoàn thành phương trình phản ứng hạt nhân (tình huống 1). Hãy xác định tính đúng, sai của bốn phát biểu sau.
\choiceTF
{\True Muốn tìm hạt còn thiếu phải bảo toàn tổng số khối $A$ và tổng điện tích $Z$ ở hai vế.}
{Chỉ cần bảo toàn số khối, không cần bảo toàn điện tích.}
{\True Khi giải cần đọc đúng dữ kiện, dùng hệ đơn vị thống nhất và kiểm tra điều kiện áp dụng.}
{Có thể bỏ qua dữ kiện, đơn vị và ý nghĩa vật lí của kết quả.}
\loigiai{
\begin{itemchoice}
\item (Đúng) Muốn tìm hạt còn thiếu phải bảo toàn tổng số khối $A$ và tổng điện tích $Z$ ở hai vế. (Vì): Muốn tìm hạt còn thiếu phải bảo toàn tổng số khối $A$ và tổng điện tích $Z$ ở hai vế. b) Sai: Chỉ cần bảo toàn số khối, không cần bảo toàn điện tích. c) Đúng: Khi giải cần đọc đúng dữ kiện, dùng hệ đơn vị thống nhất và kiểm tra điều kiện áp dụng. d) Sai: Có thể bỏ qua dữ kiện, đơn vị và ý nghĩa vật lí của kết quả. Kiến thức cốt lõi: Muốn tìm hạt còn thiếu phải bảo toàn tổng số khối $A$ và tổng điện tích $Z$ ở hai vế
\item (Sai) Chỉ cần bảo toàn số khối, không cần bảo toàn điện tích.
\item (Đúng) Khi giải cần đọc đúng dữ kiện, dùng hệ đơn vị thống nhất và kiểm tra điều kiện áp dụng.
\item (Sai) Có thể bỏ qua dữ kiện, đơn vị và ý nghĩa vật lí của kết quả.
\end{itemchoice}
}
\end{ex}

% ===== Câu 22 =====
% ID: L12C4B22-06-TL
% Mức: VDC
\dangbt{Bài toán tổng hợp về phản ứng hạt nhân, phân hạch và nhiệt hạch}
\begin{bt}
% ID: L12C4B22-06-TL
% Mức: VDC
So sánh nhận định đúng “Phân hạch hạt nhân nặng và nhiệt hạch hạt nhân nhẹ đều có thể tỏa năng lượng do sản phẩm bền hơn về năng lượng liên kết riêng.” với nhận định sai “Mọi phản ứng hạt nhân đều tự phát và tỏa cùng một năng lượng.”; chỉ rõ căn cứ Vật lí.
\loigiai{
Cần nêu được: Phân hạch hạt nhân nặng và nhiệt hạch hạt nhân nhẹ đều có thể tỏa năng lượng do sản phẩm bền hơn về năng lượng liên kết riêng. Khi vận dụng phải xác định đúng dữ kiện, điều kiện áp dụng, hệ đơn vị và kiểm tra tính hợp lí của kết quả. Nhận định “Mọi phản ứng hạt nhân đều tự phát và tỏa cùng một năng lượng.” sai vì trái với kiến thức trên.
}
\end{bt}

% ===== Câu 23 =====
% ID: L12C4B22-06-TLN
% Mức: VDC
\begin{ex}
% ID: L12C4B22-06-TLN
% Mức: VDC
Trong bốn phát biểu sau về bài toán tổng hợp về phản ứng hạt nhân, phân hạch và nhiệt hạch, có bao nhiêu phát biểu đúng? Chỉ ghi một số nguyên. (1) Kết quả cần có ý nghĩa vật lí. (2) Mọi phản ứng hạt nhân đều tự phát và tỏa cùng một năng lượng. (3) Phải dùng đúng định luật cho tình huống. (4) Phân hạch hạt nhân nặng và nhiệt hạch hạt nhân nhẹ đều có thể tỏa năng lượng do sản phẩm bền hơn về năng lượng liên kết riêng.
\shortans{3}
\loigiai{
Đối chiếu từng phát biểu với kiến thức cốt lõi: Phân hạch hạt nhân nặng và nhiệt hạch hạt nhân nhẹ đều có thể tỏa năng lượng do sản phẩm bền hơn về năng lượng liên kết riêng. Có 3 phát biểu đúng.
}
\end{ex}

% ===== Câu 24 =====
% ID: L12C4B22-06-TN
% Mức: TH
\begin{ex}
% ID: L12C4B22-06-TN
% Mức: TH
Kết luận nào của học sinh là đúng về bài toán tổng hợp về phản ứng hạt nhân, phân hạch và nhiệt hạch?
\choice
{Mọi phản ứng hạt nhân đều tự phát và tỏa cùng một năng lượng.}
{\True Phân hạch hạt nhân nặng và nhiệt hạch hạt nhân nhẹ đều có thể tỏa năng lượng do sản phẩm bền hơn về năng lượng liên kết riêng.}
{Có thể bỏ qua điều kiện áp dụng và đơn vị trong mọi trường hợp.}
{Kết quả của bài toán không phụ thuộc dữ kiện đã cho.}
\loigiai{
Phân hạch hạt nhân nặng và nhiệt hạch hạt nhân nhẹ đều có thể tỏa năng lượng do sản phẩm bền hơn về năng lượng liên kết riêng. Vì vậy phương án $B$ đúng; các phương án còn lại trái với kiến thức hoặc điều kiện áp dụng.
}
\end{ex}

% ===== Câu 25 =====
% ID: L12C4B22-07-DS
% Mức: TH
\dangbt{Hoàn thành phương trình phản ứng hạt nhân}
\begin{ex}
% ID: L12C4B22-07-DS
% Mức: TH
Xét hoàn thành phương trình phản ứng hạt nhân (tình huống 2). Hãy xác định tính đúng, sai của bốn phát biểu sau.
\choiceTF
{Chỉ cần bảo toàn số khối, không cần bảo toàn điện tích.}
{\True Muốn tìm hạt còn thiếu phải bảo toàn tổng số khối $A$ và tổng điện tích $Z$ ở hai vế.}
{Có thể bỏ qua dữ kiện, đơn vị và ý nghĩa vật lí của kết quả.}
{\True Khi giải cần đọc đúng dữ kiện, dùng hệ đơn vị thống nhất và kiểm tra điều kiện áp dụng.}
\loigiai{
\begin{itemchoice}
\item (Sai) Chỉ cần bảo toàn số khối, không cần bảo toàn điện tích. (Vì): Chỉ cần bảo toàn số khối, không cần bảo toàn điện tích. b) Đúng: Muốn tìm hạt còn thiếu phải bảo toàn tổng số khối $A$ và tổng điện tích $Z$ ở hai vế. c) Sai: Có thể bỏ qua dữ kiện, đơn vị và ý nghĩa vật lí của kết quả. d) Đúng: Khi giải cần đọc đúng dữ kiện, dùng hệ đơn vị thống nhất và kiểm tra điều kiện áp dụng. Kiến thức cốt lõi: Muốn tìm hạt còn thiếu phải bảo toàn tổng số khối $A$ và tổng điện tích $Z$ ở hai vế
\item (Đúng) Muốn tìm hạt còn thiếu phải bảo toàn tổng số khối $A$ và tổng điện tích $Z$ ở hai vế.
\item (Sai) Có thể bỏ qua dữ kiện, đơn vị và ý nghĩa vật lí của kết quả.
\item (Đúng) Khi giải cần đọc đúng dữ kiện, dùng hệ đơn vị thống nhất và kiểm tra điều kiện áp dụng.
\end{itemchoice}
}
\end{ex}

% ===== Câu 26 =====
% ID: L12C4B22-07-TL
% Mức: TH
\dangbt{Nhận biết phản ứng hạt nhân và các định luật bảo toàn}
\begin{bt}
% ID: L12C4B22-07-TL
% Mức: TH
Trình bày kiến thức cốt lõi của nhận biết phản ứng hạt nhân và các định luật bảo toàn và nêu điều kiện áp dụng.
\loigiai{
Cần nêu được: Trong phản ứng hạt nhân bảo toàn điện tích, số nucleon, năng lượng toàn phần và động lượng. Khi vận dụng phải xác định đúng dữ kiện, điều kiện áp dụng, hệ đơn vị và kiểm tra tính hợp lí của kết quả. Nhận định “Trong phản ứng hạt nhân luôn bảo toàn khối lượng nghỉ.” sai vì trái với kiến thức trên.
}
\end{bt}

% ===== Câu 27 =====
% ID: L12C4B22-07-TLN
% Mức: TH
\begin{ex}
% ID: L12C4B22-07-TLN
% Mức: TH
Trong bốn phát biểu sau về nhận biết phản ứng hạt nhân và các định luật bảo toàn, có bao nhiêu phát biểu đúng? Chỉ ghi một số nguyên. (1) Trong phản ứng hạt nhân bảo toàn điện tích, số nucleon, năng lượng toàn phần và động lượng. (2) Trong phản ứng hạt nhân luôn bảo toàn khối lượng nghỉ. (3) Cần kiểm tra điều kiện áp dụng trước khi tính toán. (4) Có thể bỏ qua đơn vị của kết quả.
\shortans{2}
\loigiai{
Đối chiếu từng phát biểu với kiến thức cốt lõi: Trong phản ứng hạt nhân bảo toàn điện tích, số nucleon, năng lượng toàn phần và động lượng. Có 2 phát biểu đúng.
}
\end{ex}

% ===== Câu 28 =====
% ID: L12C4B22-07-TN
% Mức: TH
\dangbt{Bài toán tổng hợp về phản ứng hạt nhân, phân hạch và nhiệt hạch}
\begin{ex}
% ID: L12C4B22-07-TN
% Mức: TH
Chọn phát biểu không trái với kiến thức về bài toán tổng hợp về phản ứng hạt nhân, phân hạch và nhiệt hạch?
\choice
{Mọi phản ứng hạt nhân đều tự phát và tỏa cùng một năng lượng.}
{Có thể bỏ qua điều kiện áp dụng và đơn vị trong mọi trường hợp.}
{\True Phân hạch hạt nhân nặng và nhiệt hạch hạt nhân nhẹ đều có thể tỏa năng lượng do sản phẩm bền hơn về năng lượng liên kết riêng.}
{Kết quả của bài toán không phụ thuộc dữ kiện đã cho.}
\loigiai{
Phân hạch hạt nhân nặng và nhiệt hạch hạt nhân nhẹ đều có thể tỏa năng lượng do sản phẩm bền hơn về năng lượng liên kết riêng. Vì vậy phương án $C$ đúng; các phương án còn lại trái với kiến thức hoặc điều kiện áp dụng.
}
\end{ex}

% ===== Câu 29 =====
% ID: L12C4B22-08-DS
% Mức: VD
\dangbt{Hoàn thành phương trình phản ứng hạt nhân}
\begin{ex}
% ID: L12C4B22-08-DS
% Mức: VD
Xét hoàn thành phương trình phản ứng hạt nhân (tình huống 3). Hãy xác định tính đúng, sai của bốn phát biểu sau.
\choiceTF
{\True Muốn tìm hạt còn thiếu phải bảo toàn tổng số khối $A$ và tổng điện tích $Z$ ở hai vế.}
{Chỉ cần bảo toàn số khối, không cần bảo toàn điện tích.}
{\True Khi giải cần đọc đúng dữ kiện, dùng hệ đơn vị thống nhất và kiểm tra điều kiện áp dụng.}
{Có thể bỏ qua dữ kiện, đơn vị và ý nghĩa vật lí của kết quả.}
\loigiai{
\begin{itemchoice}
\item (Đúng) Muốn tìm hạt còn thiếu phải bảo toàn tổng số khối $A$ và tổng điện tích $Z$ ở hai vế. (Vì): Muốn tìm hạt còn thiếu phải bảo toàn tổng số khối $A$ và tổng điện tích $Z$ ở hai vế. b) Sai: Chỉ cần bảo toàn số khối, không cần bảo toàn điện tích. c) Đúng: Khi giải cần đọc đúng dữ kiện, dùng hệ đơn vị thống nhất và kiểm tra điều kiện áp dụng. d) Sai: Có thể bỏ qua dữ kiện, đơn vị và ý nghĩa vật lí của kết quả. Kiến thức cốt lõi: Muốn tìm hạt còn thiếu phải bảo toàn tổng số khối $A$ và tổng điện tích $Z$ ở hai vế
\item (Sai) Chỉ cần bảo toàn số khối, không cần bảo toàn điện tích.
\item (Đúng) Khi giải cần đọc đúng dữ kiện, dùng hệ đơn vị thống nhất và kiểm tra điều kiện áp dụng.
\item (Sai) Có thể bỏ qua dữ kiện, đơn vị và ý nghĩa vật lí của kết quả.
\end{itemchoice}
}
\end{ex}

% ===== Câu 30 =====
% ID: L12C4B22-08-TL
% Mức: TH
\dangbt{Nhận biết phản ứng hạt nhân và các định luật bảo toàn}
\begin{bt}
% ID: L12C4B22-08-TL
% Mức: TH
Giải thích vì sao nhận định sau đúng: “Trong phản ứng hạt nhân bảo toàn điện tích, số nucleon, năng lượng toàn phần và động lượng.”
\loigiai{
Cần nêu được: Trong phản ứng hạt nhân bảo toàn điện tích, số nucleon, năng lượng toàn phần và động lượng. Khi vận dụng phải xác định đúng dữ kiện, điều kiện áp dụng, hệ đơn vị và kiểm tra tính hợp lí của kết quả. Nhận định “Trong phản ứng hạt nhân luôn bảo toàn khối lượng nghỉ.” sai vì trái với kiến thức trên.
}
\end{bt}

% ===== Câu 31 =====
% ID: L12C4B22-08-TLN
% Mức: TH
\begin{ex}
% ID: L12C4B22-08-TLN
% Mức: TH
Trong bốn phát biểu sau về nhận biết phản ứng hạt nhân và các định luật bảo toàn, có bao nhiêu phát biểu đúng? Chỉ ghi một số nguyên. (1) Trong phản ứng hạt nhân luôn bảo toàn khối lượng nghỉ. (2) Trong phản ứng hạt nhân bảo toàn điện tích, số nucleon, năng lượng toàn phần và động lượng. (3) Kết quả phải phù hợp ý nghĩa vật lí. (4) Mọi đại lượng đều có cùng bản chất.
\shortans{2}
\loigiai{
Đối chiếu từng phát biểu với kiến thức cốt lõi: Trong phản ứng hạt nhân bảo toàn điện tích, số nucleon, năng lượng toàn phần và động lượng. Có 2 phát biểu đúng.
}
\end{ex}

% ===== Câu 32 =====
% ID: L12C4B22-08-TN
% Mức: VD
\dangbt{Bài toán tổng hợp về phản ứng hạt nhân, phân hạch và nhiệt hạch}
\begin{ex}
% ID: L12C4B22-08-TN
% Mức: VD
Cơ sở Vật lí đúng để giải bài là gì về bài toán tổng hợp về phản ứng hạt nhân, phân hạch và nhiệt hạch?
\choice
{Mọi phản ứng hạt nhân đều tự phát và tỏa cùng một năng lượng.}
{Có thể bỏ qua điều kiện áp dụng và đơn vị trong mọi trường hợp.}
{Kết quả của bài toán không phụ thuộc dữ kiện đã cho.}
{\True Phân hạch hạt nhân nặng và nhiệt hạch hạt nhân nhẹ đều có thể tỏa năng lượng do sản phẩm bền hơn về năng lượng liên kết riêng.}
\loigiai{
Phân hạch hạt nhân nặng và nhiệt hạch hạt nhân nhẹ đều có thể tỏa năng lượng do sản phẩm bền hơn về năng lượng liên kết riêng. Vì vậy phương án $D$ đúng; các phương án còn lại trái với kiến thức hoặc điều kiện áp dụng.
}
\end{ex}

% ===== Câu 33 =====
% ID: L12C4B22-09-DS
% Mức: VD
\dangbt{Hoàn thành phương trình phản ứng hạt nhân}
\begin{ex}
% ID: L12C4B22-09-DS
% Mức: VD
Xét hoàn thành phương trình phản ứng hạt nhân (tình huống 4). Hãy xác định tính đúng, sai của bốn phát biểu sau.
\choiceTF
{Chỉ cần bảo toàn số khối, không cần bảo toàn điện tích.}
{\True Muốn tìm hạt còn thiếu phải bảo toàn tổng số khối $A$ và tổng điện tích $Z$ ở hai vế.}
{Có thể bỏ qua dữ kiện, đơn vị và ý nghĩa vật lí của kết quả.}
{\True Khi giải cần đọc đúng dữ kiện, dùng hệ đơn vị thống nhất và kiểm tra điều kiện áp dụng.}
\loigiai{
\begin{itemchoice}
\item (Sai) Chỉ cần bảo toàn số khối, không cần bảo toàn điện tích. (Vì): Chỉ cần bảo toàn số khối, không cần bảo toàn điện tích. b) Đúng: Muốn tìm hạt còn thiếu phải bảo toàn tổng số khối $A$ và tổng điện tích $Z$ ở hai vế. c) Sai: Có thể bỏ qua dữ kiện, đơn vị và ý nghĩa vật lí của kết quả. d) Đúng: Khi giải cần đọc đúng dữ kiện, dùng hệ đơn vị thống nhất và kiểm tra điều kiện áp dụng. Kiến thức cốt lõi: Muốn tìm hạt còn thiếu phải bảo toàn tổng số khối $A$ và tổng điện tích $Z$ ở hai vế
\item (Đúng) Muốn tìm hạt còn thiếu phải bảo toàn tổng số khối $A$ và tổng điện tích $Z$ ở hai vế.
\item (Sai) Có thể bỏ qua dữ kiện, đơn vị và ý nghĩa vật lí của kết quả.
\item (Đúng) Khi giải cần đọc đúng dữ kiện, dùng hệ đơn vị thống nhất và kiểm tra điều kiện áp dụng.
\end{itemchoice}
}
\end{ex}

% ===== Câu 34 =====
% ID: L12C4B22-09-TL
% Mức: VD
\dangbt{Nhận biết phản ứng hạt nhân và các định luật bảo toàn}
\begin{bt}
% ID: L12C4B22-09-TL
% Mức: VD
Phân tích sai lầm trong nhận định: “Trong phản ứng hạt nhân luôn bảo toàn khối lượng nghỉ.”
\loigiai{
Cần nêu được: Trong phản ứng hạt nhân bảo toàn điện tích, số nucleon, năng lượng toàn phần và động lượng. Khi vận dụng phải xác định đúng dữ kiện, điều kiện áp dụng, hệ đơn vị và kiểm tra tính hợp lí của kết quả. Nhận định “Trong phản ứng hạt nhân luôn bảo toàn khối lượng nghỉ.” sai vì trái với kiến thức trên.
}
\end{bt}

% ===== Câu 35 =====
% ID: L12C4B22-09-TLN
% Mức: TH
\begin{ex}
% ID: L12C4B22-09-TLN
% Mức: TH
Trong bốn phát biểu sau về nhận biết phản ứng hạt nhân và các định luật bảo toàn, có bao nhiêu phát biểu đúng? Chỉ ghi một số nguyên. (1) Cần đổi các đại lượng về hệ đơn vị thống nhất. (2) Trong phản ứng hạt nhân bảo toàn điện tích, số nucleon, năng lượng toàn phần và động lượng. (3) Trong phản ứng hạt nhân luôn bảo toàn khối lượng nghỉ. (4) Không cần kiểm tra dữ kiện.
\shortans{2}
\loigiai{
Đối chiếu từng phát biểu với kiến thức cốt lõi: Trong phản ứng hạt nhân bảo toàn điện tích, số nucleon, năng lượng toàn phần và động lượng. Có 2 phát biểu đúng.
}
\end{ex}

% ===== Câu 36 =====
% ID: L12C4B22-09-TN
% Mức: VD
\dangbt{Bài toán tổng hợp về phản ứng hạt nhân, phân hạch và nhiệt hạch}
\begin{ex}
% ID: L12C4B22-09-TN
% Mức: VD
Trong một bài thực hành, nhận xét nào đúng về bài toán tổng hợp về phản ứng hạt nhân, phân hạch và nhiệt hạch?
\choice
{\True Phân hạch hạt nhân nặng và nhiệt hạch hạt nhân nhẹ đều có thể tỏa năng lượng do sản phẩm bền hơn về năng lượng liên kết riêng.}
{Mọi phản ứng hạt nhân đều tự phát và tỏa cùng một năng lượng.}
{Có thể bỏ qua điều kiện áp dụng và đơn vị trong mọi trường hợp.}
{Kết quả của bài toán không phụ thuộc dữ kiện đã cho.}
\loigiai{
Phân hạch hạt nhân nặng và nhiệt hạch hạt nhân nhẹ đều có thể tỏa năng lượng do sản phẩm bền hơn về năng lượng liên kết riêng. Vì vậy phương án $A$ đúng; các phương án còn lại trái với kiến thức hoặc điều kiện áp dụng.
}
\end{ex}

% ===== Câu 37 =====
% ID: L12C4B22-10-DS
% Mức: TH
\dangbt{Nhận biết phản ứng hạt nhân và các định luật bảo toàn}
\begin{ex}
% ID: L12C4B22-10-DS
% Mức: TH
Xét nhận biết phản ứng hạt nhân và các định luật bảo toàn (tình huống 1). Hãy xác định tính đúng, sai của bốn phát biểu sau.
\choiceTF
{\True Trong phản ứng hạt nhân bảo toàn điện tích, số nucleon, năng lượng toàn phần và động lượng.}
{Trong phản ứng hạt nhân luôn bảo toàn khối lượng nghỉ.}
{\True Khi giải cần đọc đúng dữ kiện, dùng hệ đơn vị thống nhất và kiểm tra điều kiện áp dụng.}
{Có thể bỏ qua dữ kiện, đơn vị và ý nghĩa vật lí của kết quả.}
\loigiai{
\begin{itemchoice}
\item (Đúng) Trong phản ứng hạt nhân bảo toàn điện tích, số nucleon, năng lượng toàn phần và động lượng. (Vì): Trong phản ứng hạt nhân bảo toàn điện tích, số nucleon, năng lượng toàn phần và động lượng. b) Sai: Trong phản ứng hạt nhân luôn bảo toàn khối lượng nghỉ. c) Đúng: Khi giải cần đọc đúng dữ kiện, dùng hệ đơn vị thống nhất và kiểm tra điều kiện áp dụng. d) Sai: Có thể bỏ qua dữ kiện, đơn vị và ý nghĩa vật lí của kết quả. Kiến thức cốt lõi: Trong phản ứng hạt nhân bảo toàn điện tích, số nucleon, năng lượng toàn phần và động lượng
\item (Sai) Trong phản ứng hạt nhân luôn bảo toàn khối lượng nghỉ.
\item (Đúng) Khi giải cần đọc đúng dữ kiện, dùng hệ đơn vị thống nhất và kiểm tra điều kiện áp dụng.
\item (Sai) Có thể bỏ qua dữ kiện, đơn vị và ý nghĩa vật lí của kết quả.
\end{itemchoice}
}
\end{ex}

% ===== Câu 38 =====
% ID: L12C4B22-10-TL
% Mức: VD
\begin{bt}
% ID: L12C4B22-10-TL
% Mức: VD
Đề xuất các bước giải một bài toán thuộc dạng nhận biết phản ứng hạt nhân và các định luật bảo toàn.
\loigiai{
Cần nêu được: Trong phản ứng hạt nhân bảo toàn điện tích, số nucleon, năng lượng toàn phần và động lượng. Khi vận dụng phải xác định đúng dữ kiện, điều kiện áp dụng, hệ đơn vị và kiểm tra tính hợp lí của kết quả. Nhận định “Trong phản ứng hạt nhân luôn bảo toàn khối lượng nghỉ.” sai vì trái với kiến thức trên.
}
\end{bt}

% ===== Câu 39 =====
% ID: L12C4B22-10-TLN
% Mức: VD
\begin{ex}
% ID: L12C4B22-10-TLN
% Mức: VD
Trong bốn phát biểu sau về nhận biết phản ứng hạt nhân và các định luật bảo toàn, có bao nhiêu phát biểu đúng? Chỉ ghi một số nguyên. (1) Trong phản ứng hạt nhân bảo toàn điện tích, số nucleon, năng lượng toàn phần và động lượng. (2) Cần đối chiếu kết quả với điều kiện bài toán. (3) Trong phản ứng hạt nhân luôn bảo toàn khối lượng nghỉ. (4) Mọi trường hợp đều cho cùng một kết quả.
\shortans{2}
\loigiai{
Đối chiếu từng phát biểu với kiến thức cốt lõi: Trong phản ứng hạt nhân bảo toàn điện tích, số nucleon, năng lượng toàn phần và động lượng. Có 2 phát biểu đúng.
}
\end{ex}

% ===== Câu 40 =====
% ID: L12C4B22-10-TN
% Mức: VD
\dangbt{Bài toán tổng hợp về phản ứng hạt nhân, phân hạch và nhiệt hạch}
\begin{ex}
% ID: L12C4B22-10-TN
% Mức: VD
Kết luận đúng khi vận dụng là về bài toán tổng hợp về phản ứng hạt nhân, phân hạch và nhiệt hạch?
\choice
{Mọi phản ứng hạt nhân đều tự phát và tỏa cùng một năng lượng.}
{\True Phân hạch hạt nhân nặng và nhiệt hạch hạt nhân nhẹ đều có thể tỏa năng lượng do sản phẩm bền hơn về năng lượng liên kết riêng.}
{Có thể bỏ qua điều kiện áp dụng và đơn vị trong mọi trường hợp.}
{Kết quả của bài toán không phụ thuộc dữ kiện đã cho.}
\loigiai{
Phân hạch hạt nhân nặng và nhiệt hạch hạt nhân nhẹ đều có thể tỏa năng lượng do sản phẩm bền hơn về năng lượng liên kết riêng. Vì vậy phương án $B$ đúng; các phương án còn lại trái với kiến thức hoặc điều kiện áp dụng.
}
\end{ex}

% ===== Câu 41 =====
% ID: L12C4B22-11-DS
% Mức: TH
\dangbt{Nhận biết phản ứng hạt nhân và các định luật bảo toàn}
\begin{ex}
% ID: L12C4B22-11-DS
% Mức: TH
Xét nhận biết phản ứng hạt nhân và các định luật bảo toàn (tình huống 2). Hãy xác định tính đúng, sai của bốn phát biểu sau.
\choiceTF
{Trong phản ứng hạt nhân luôn bảo toàn khối lượng nghỉ.}
{\True Trong phản ứng hạt nhân bảo toàn điện tích, số nucleon, năng lượng toàn phần và động lượng.}
{Có thể bỏ qua dữ kiện, đơn vị và ý nghĩa vật lí của kết quả.}
{\True Khi giải cần đọc đúng dữ kiện, dùng hệ đơn vị thống nhất và kiểm tra điều kiện áp dụng.}
\loigiai{
\begin{itemchoice}
\item (Sai) Trong phản ứng hạt nhân luôn bảo toàn khối lượng nghỉ. (Vì): Trong phản ứng hạt nhân luôn bảo toàn khối lượng nghỉ. b) Đúng: Trong phản ứng hạt nhân bảo toàn điện tích, số nucleon, năng lượng toàn phần và động lượng. c) Sai: Có thể bỏ qua dữ kiện, đơn vị và ý nghĩa vật lí của kết quả. d) Đúng: Khi giải cần đọc đúng dữ kiện, dùng hệ đơn vị thống nhất và kiểm tra điều kiện áp dụng. Kiến thức cốt lõi: Trong phản ứng hạt nhân bảo toàn điện tích, số nucleon, năng lượng toàn phần và động lượng
\item (Đúng) Trong phản ứng hạt nhân bảo toàn điện tích, số nucleon, năng lượng toàn phần và động lượng.
\item (Sai) Có thể bỏ qua dữ kiện, đơn vị và ý nghĩa vật lí của kết quả.
\item (Đúng) Khi giải cần đọc đúng dữ kiện, dùng hệ đơn vị thống nhất và kiểm tra điều kiện áp dụng.
\end{itemchoice}
}
\end{ex}

% ===== Câu 42 =====
% ID: L12C4B22-11-TL
% Mức: VD
\begin{bt}
% ID: L12C4B22-11-TL
% Mức: VD
Nêu một tình huống thực tiễn liên quan đến nhận biết phản ứng hạt nhân và các định luật bảo toàn và giải thích bằng kiến thức Vật lí.
\loigiai{
Cần nêu được: Trong phản ứng hạt nhân bảo toàn điện tích, số nucleon, năng lượng toàn phần và động lượng. Khi vận dụng phải xác định đúng dữ kiện, điều kiện áp dụng, hệ đơn vị và kiểm tra tính hợp lí của kết quả. Nhận định “Trong phản ứng hạt nhân luôn bảo toàn khối lượng nghỉ.” sai vì trái với kiến thức trên.
}
\end{bt}

% ===== Câu 43 =====
% ID: L12C4B22-11-TLN
% Mức: VD
\begin{ex}
% ID: L12C4B22-11-TLN
% Mức: VD
Trong bốn phát biểu sau về nhận biết phản ứng hạt nhân và các định luật bảo toàn, có bao nhiêu phát biểu đúng? Chỉ ghi một số nguyên. (1) Trong phản ứng hạt nhân luôn bảo toàn khối lượng nghỉ. (2) Cần đọc đúng dữ kiện và giả thiết. (3) Trong phản ứng hạt nhân bảo toàn điện tích, số nucleon, năng lượng toàn phần và động lượng. (4) Có thể thay số trước khi chọn công thức.
\shortans{2}
\loigiai{
Đối chiếu từng phát biểu với kiến thức cốt lõi: Trong phản ứng hạt nhân bảo toàn điện tích, số nucleon, năng lượng toàn phần và động lượng. Có 2 phát biểu đúng.
}
\end{ex}

% ===== Câu 44 =====
% ID: L12C4B22-11-TN
% Mức: NB
\dangbt{Hoàn thành phương trình phản ứng hạt nhân}
\begin{ex}
% ID: L12C4B22-11-TN
% Mức: NB
Phát biểu nào sau đây đúng về hoàn thành phương trình phản ứng hạt nhân?
\choice
{\True Muốn tìm hạt còn thiếu phải bảo toàn tổng số khối $A$ và tổng điện tích $Z$ ở hai vế.}
{Chỉ cần bảo toàn số khối, không cần bảo toàn điện tích.}
{Có thể bỏ qua điều kiện áp dụng và đơn vị trong mọi trường hợp.}
{Kết quả của bài toán không phụ thuộc dữ kiện đã cho.}
\loigiai{
Muốn tìm hạt còn thiếu phải bảo toàn tổng số khối $A$ và tổng điện tích $Z$ ở hai vế. Vì vậy phương án $A$ đúng; các phương án còn lại trái với kiến thức hoặc điều kiện áp dụng.
}
\end{ex}

% ===== Câu 45 =====
% ID: L12C4B22-12-DS
% Mức: VD
\dangbt{Nhận biết phản ứng hạt nhân và các định luật bảo toàn}
\begin{ex}
% ID: L12C4B22-12-DS
% Mức: VD
Xét nhận biết phản ứng hạt nhân và các định luật bảo toàn (tình huống 3). Hãy xác định tính đúng, sai của bốn phát biểu sau.
\choiceTF
{\True Trong phản ứng hạt nhân bảo toàn điện tích, số nucleon, năng lượng toàn phần và động lượng.}
{Trong phản ứng hạt nhân luôn bảo toàn khối lượng nghỉ.}
{\True Khi giải cần đọc đúng dữ kiện, dùng hệ đơn vị thống nhất và kiểm tra điều kiện áp dụng.}
{Có thể bỏ qua dữ kiện, đơn vị và ý nghĩa vật lí của kết quả.}
\loigiai{
\begin{itemchoice}
\item (Đúng) Trong phản ứng hạt nhân bảo toàn điện tích, số nucleon, năng lượng toàn phần và động lượng. (Vì): Trong phản ứng hạt nhân bảo toàn điện tích, số nucleon, năng lượng toàn phần và động lượng. b) Sai: Trong phản ứng hạt nhân luôn bảo toàn khối lượng nghỉ. c) Đúng: Khi giải cần đọc đúng dữ kiện, dùng hệ đơn vị thống nhất và kiểm tra điều kiện áp dụng. d) Sai: Có thể bỏ qua dữ kiện, đơn vị và ý nghĩa vật lí của kết quả. Kiến thức cốt lõi: Trong phản ứng hạt nhân bảo toàn điện tích, số nucleon, năng lượng toàn phần và động lượng
\item (Sai) Trong phản ứng hạt nhân luôn bảo toàn khối lượng nghỉ.
\item (Đúng) Khi giải cần đọc đúng dữ kiện, dùng hệ đơn vị thống nhất và kiểm tra điều kiện áp dụng.
\item (Sai) Có thể bỏ qua dữ kiện, đơn vị và ý nghĩa vật lí của kết quả.
\end{itemchoice}
}
\end{ex}

% ===== Câu 46 =====
% ID: L12C4B22-12-TL
% Mức: VDC
\begin{bt}
% ID: L12C4B22-12-TL
% Mức: VDC
So sánh nhận định đúng “Trong phản ứng hạt nhân bảo toàn điện tích, số nucleon, năng lượng toàn phần và động lượng.” với nhận định sai “Trong phản ứng hạt nhân luôn bảo toàn khối lượng nghỉ.”; chỉ rõ căn cứ Vật lí.
\loigiai{
Cần nêu được: Trong phản ứng hạt nhân bảo toàn điện tích, số nucleon, năng lượng toàn phần và động lượng. Khi vận dụng phải xác định đúng dữ kiện, điều kiện áp dụng, hệ đơn vị và kiểm tra tính hợp lí của kết quả. Nhận định “Trong phản ứng hạt nhân luôn bảo toàn khối lượng nghỉ.” sai vì trái với kiến thức trên.
}
\end{bt}

% ===== Câu 47 =====
% ID: L12C4B22-12-TLN
% Mức: VDC
\begin{ex}
% ID: L12C4B22-12-TLN
% Mức: VDC
Trong bốn phát biểu sau về nhận biết phản ứng hạt nhân và các định luật bảo toàn, có bao nhiêu phát biểu đúng? Chỉ ghi một số nguyên. (1) Kết quả cần có ý nghĩa vật lí. (2) Trong phản ứng hạt nhân luôn bảo toàn khối lượng nghỉ. (3) Phải dùng đúng định luật cho tình huống. (4) Trong phản ứng hạt nhân bảo toàn điện tích, số nucleon, năng lượng toàn phần và động lượng.
\shortans{3}
\loigiai{
Đối chiếu từng phát biểu với kiến thức cốt lõi: Trong phản ứng hạt nhân bảo toàn điện tích, số nucleon, năng lượng toàn phần và động lượng. Có 3 phát biểu đúng.
}
\end{ex}

% ===== Câu 48 =====
% ID: L12C4B22-12-TN
% Mức: NB
\dangbt{Hoàn thành phương trình phản ứng hạt nhân}
\begin{ex}
% ID: L12C4B22-12-TN
% Mức: NB
Nhận định nào phù hợp với kiến thức Vật lí về hoàn thành phương trình phản ứng hạt nhân?
\choice
{Chỉ cần bảo toàn số khối, không cần bảo toàn điện tích.}
{\True Muốn tìm hạt còn thiếu phải bảo toàn tổng số khối $A$ và tổng điện tích $Z$ ở hai vế.}
{Có thể bỏ qua điều kiện áp dụng và đơn vị trong mọi trường hợp.}
{Kết quả của bài toán không phụ thuộc dữ kiện đã cho.}
\loigiai{
Muốn tìm hạt còn thiếu phải bảo toàn tổng số khối $A$ và tổng điện tích $Z$ ở hai vế. Vì vậy phương án $B$ đúng; các phương án còn lại trái với kiến thức hoặc điều kiện áp dụng.
}
\end{ex}

% ===== Câu 49 =====
% ID: L12C4B22-13-DS
% Mức: VD
\dangbt{Nhận biết phản ứng hạt nhân và các định luật bảo toàn}
\begin{ex}
% ID: L12C4B22-13-DS
% Mức: VD
Xét nhận biết phản ứng hạt nhân và các định luật bảo toàn (tình huống 4). Hãy xác định tính đúng, sai của bốn phát biểu sau.
\choiceTF
{Trong phản ứng hạt nhân luôn bảo toàn khối lượng nghỉ.}
{\True Trong phản ứng hạt nhân bảo toàn điện tích, số nucleon, năng lượng toàn phần và động lượng.}
{Có thể bỏ qua dữ kiện, đơn vị và ý nghĩa vật lí của kết quả.}
{\True Khi giải cần đọc đúng dữ kiện, dùng hệ đơn vị thống nhất và kiểm tra điều kiện áp dụng.}
\loigiai{
\begin{itemchoice}
\item (Sai) Trong phản ứng hạt nhân luôn bảo toàn khối lượng nghỉ. (Vì): Trong phản ứng hạt nhân luôn bảo toàn khối lượng nghỉ. b) Đúng: Trong phản ứng hạt nhân bảo toàn điện tích, số nucleon, năng lượng toàn phần và động lượng. c) Sai: Có thể bỏ qua dữ kiện, đơn vị và ý nghĩa vật lí của kết quả. d) Đúng: Khi giải cần đọc đúng dữ kiện, dùng hệ đơn vị thống nhất và kiểm tra điều kiện áp dụng. Kiến thức cốt lõi: Trong phản ứng hạt nhân bảo toàn điện tích, số nucleon, năng lượng toàn phần và động lượng
\item (Đúng) Trong phản ứng hạt nhân bảo toàn điện tích, số nucleon, năng lượng toàn phần và động lượng.
\item (Sai) Có thể bỏ qua dữ kiện, đơn vị và ý nghĩa vật lí của kết quả.
\item (Đúng) Khi giải cần đọc đúng dữ kiện, dùng hệ đơn vị thống nhất và kiểm tra điều kiện áp dụng.
\end{itemchoice}
}
\end{ex}

% ===== Câu 50 =====
% ID: L12C4B22-13-TL
% Mức: TH
\dangbt{Hoàn thành phương trình phản ứng hạt nhân}
\begin{bt}
% ID: L12C4B22-13-TL
% Mức: TH
Trình bày kiến thức cốt lõi của hoàn thành phương trình phản ứng hạt nhân và nêu điều kiện áp dụng.
\loigiai{
Cần nêu được: Muốn tìm hạt còn thiếu phải bảo toàn tổng số khối $A$ và tổng điện tích $Z$ ở hai vế. Khi vận dụng phải xác định đúng dữ kiện, điều kiện áp dụng, hệ đơn vị và kiểm tra tính hợp lí của kết quả. Nhận định “Chỉ cần bảo toàn số khối, không cần bảo toàn điện tích.” sai vì trái với kiến thức trên.
}
\end{bt}

% ===== Câu 51 =====
% ID: L12C4B22-13-TLN
% Mức: TH
\begin{ex}
% ID: L12C4B22-13-TLN
% Mức: TH
Trong bốn phát biểu sau về hoàn thành phương trình phản ứng hạt nhân, có bao nhiêu phát biểu đúng? Chỉ ghi một số nguyên. (1) Muốn tìm hạt còn thiếu phải bảo toàn tổng số khối $A$ và tổng điện tích $Z$ ở hai vế. (2) Chỉ cần bảo toàn số khối, không cần bảo toàn điện tích. (3) Cần kiểm tra điều kiện áp dụng trước khi tính toán. (4) Có thể bỏ qua đơn vị của kết quả.
\shortans{2}
\loigiai{
Đối chiếu từng phát biểu với kiến thức cốt lõi: Muốn tìm hạt còn thiếu phải bảo toàn tổng số khối $A$ và tổng điện tích $Z$ ở hai vế. Có 2 phát biểu đúng.
}
\end{ex}

% ===== Câu 52 =====
% ID: L12C4B22-13-TN
% Mức: NB
\begin{ex}
% ID: L12C4B22-13-TN
% Mức: NB
Chọn kết luận chính xác về hoàn thành phương trình phản ứng hạt nhân?
\choice
{Chỉ cần bảo toàn số khối, không cần bảo toàn điện tích.}
{Có thể bỏ qua điều kiện áp dụng và đơn vị trong mọi trường hợp.}
{\True Muốn tìm hạt còn thiếu phải bảo toàn tổng số khối $A$ và tổng điện tích $Z$ ở hai vế.}
{Kết quả của bài toán không phụ thuộc dữ kiện đã cho.}
\loigiai{
Muốn tìm hạt còn thiếu phải bảo toàn tổng số khối $A$ và tổng điện tích $Z$ ở hai vế. Vì vậy phương án $C$ đúng; các phương án còn lại trái với kiến thức hoặc điều kiện áp dụng.
}
\end{ex}

% ===== Câu 53 =====
% ID: L12C4B22-14-DS
% Mức: VD
\dangbt{Nhận biết phản ứng hạt nhân và các định luật bảo toàn}
\begin{ex}
% ID: L12C4B22-14-DS
% Mức: VD
Xét nhận biết phản ứng hạt nhân và các định luật bảo toàn (tình huống 5). Hãy xác định tính đúng, sai của bốn phát biểu sau.
\choiceTF
{\True Trong phản ứng hạt nhân bảo toàn điện tích, số nucleon, năng lượng toàn phần và động lượng.}
{Có thể bỏ qua dữ kiện, đơn vị và ý nghĩa vật lí của kết quả.}
{Trong phản ứng hạt nhân luôn bảo toàn khối lượng nghỉ.}
{Có thể bỏ qua dữ kiện, đơn vị và ý nghĩa vật lí của kết quả.}
\loigiai{
\begin{itemchoice}
\item (Đúng) Trong phản ứng hạt nhân bảo toàn điện tích, số nucleon, năng lượng toàn phần và động lượng. (Vì): Trong phản ứng hạt nhân bảo toàn điện tích, số nucleon, năng lượng toàn phần và động lượng. b) Sai: Có thể bỏ qua dữ kiện, đơn vị và ý nghĩa vật lí của kết quả. c) Sai: Trong phản ứng hạt nhân luôn bảo toàn khối lượng nghỉ. d) Sai: Có thể bỏ qua dữ kiện, đơn vị và ý nghĩa vật lí của kết quả. Kiến thức cốt lõi: Trong phản ứng hạt nhân bảo toàn điện tích, số nucleon, năng lượng toàn phần và động lượng
\item (Sai) Có thể bỏ qua dữ kiện, đơn vị và ý nghĩa vật lí của kết quả.
\item (Sai) Trong phản ứng hạt nhân luôn bảo toàn khối lượng nghỉ.
\item (Sai) Có thể bỏ qua dữ kiện, đơn vị và ý nghĩa vật lí của kết quả.
\end{itemchoice}
}
\end{ex}

% ===== Câu 54 =====
% ID: L12C4B22-14-TL
% Mức: TH
\dangbt{Hoàn thành phương trình phản ứng hạt nhân}
\begin{bt}
% ID: L12C4B22-14-TL
% Mức: TH
Giải thích vì sao nhận định sau đúng: “Muốn tìm hạt còn thiếu phải bảo toàn tổng số khối $A$ và tổng điện tích $Z$ ở hai vế.”
\loigiai{
Cần nêu được: Muốn tìm hạt còn thiếu phải bảo toàn tổng số khối $A$ và tổng điện tích $Z$ ở hai vế. Khi vận dụng phải xác định đúng dữ kiện, điều kiện áp dụng, hệ đơn vị và kiểm tra tính hợp lí của kết quả. Nhận định “Chỉ cần bảo toàn số khối, không cần bảo toàn điện tích.” sai vì trái với kiến thức trên.
}
\end{bt}

% ===== Câu 55 =====
% ID: L12C4B22-14-TLN
% Mức: TH
\begin{ex}
% ID: L12C4B22-14-TLN
% Mức: TH
Trong bốn phát biểu sau về hoàn thành phương trình phản ứng hạt nhân, có bao nhiêu phát biểu đúng? Chỉ ghi một số nguyên. (1) Chỉ cần bảo toàn số khối, không cần bảo toàn điện tích. (2) Muốn tìm hạt còn thiếu phải bảo toàn tổng số khối $A$ và tổng điện tích $Z$ ở hai vế. (3) Kết quả phải phù hợp ý nghĩa vật lí. (4) Mọi đại lượng đều có cùng bản chất.
\shortans{2}
\loigiai{
Đối chiếu từng phát biểu với kiến thức cốt lõi: Muốn tìm hạt còn thiếu phải bảo toàn tổng số khối $A$ và tổng điện tích $Z$ ở hai vế. Có 2 phát biểu đúng.
}
\end{ex}

% ===== Câu 56 =====
% ID: L12C4B22-14-TN
% Mức: TH
\begin{ex}
% ID: L12C4B22-14-TN
% Mức: TH
Nội dung nào mô tả đúng nhất về hoàn thành phương trình phản ứng hạt nhân?
\choice
{Chỉ cần bảo toàn số khối, không cần bảo toàn điện tích.}
{Có thể bỏ qua điều kiện áp dụng và đơn vị trong mọi trường hợp.}
{Kết quả của bài toán không phụ thuộc dữ kiện đã cho.}
{\True Muốn tìm hạt còn thiếu phải bảo toàn tổng số khối $A$ và tổng điện tích $Z$ ở hai vế.}
\loigiai{
Muốn tìm hạt còn thiếu phải bảo toàn tổng số khối $A$ và tổng điện tích $Z$ ở hai vế. Vì vậy phương án $D$ đúng; các phương án còn lại trái với kiến thức hoặc điều kiện áp dụng.
}
\end{ex}

% ===== Câu 57 =====
% ID: L12C4B22-15-DS
% Mức: TH
\begin{ex}
% ID: L12C4B22-15-DS
% Mức: TH
Xét hoàn thành phương trình phản ứng hạt nhân (tình huống 1). Hãy xác định tính đúng, sai của bốn phát biểu sau.
\choiceTF
{\True Muốn tìm hạt còn thiếu phải bảo toàn tổng số khối $A$ và tổng điện tích $Z$ ở hai vế.}
{Chỉ cần bảo toàn số khối, không cần bảo toàn điện tích.}
{\True Khi giải cần đọc đúng dữ kiện, dùng hệ đơn vị thống nhất và kiểm tra điều kiện áp dụng.}
{Có thể bỏ qua dữ kiện, đơn vị và ý nghĩa vật lí của kết quả.}
\loigiai{
\begin{itemchoice}
\item (Đúng) Muốn tìm hạt còn thiếu phải bảo toàn tổng số khối $A$ và tổng điện tích $Z$ ở hai vế. (Vì): Muốn tìm hạt còn thiếu phải bảo toàn tổng số khối $A$ và tổng điện tích $Z$ ở hai vế. b) Sai: Chỉ cần bảo toàn số khối, không cần bảo toàn điện tích. c) Đúng: Khi giải cần đọc đúng dữ kiện, dùng hệ đơn vị thống nhất và kiểm tra điều kiện áp dụng. d) Sai: Có thể bỏ qua dữ kiện, đơn vị và ý nghĩa vật lí của kết quả. Kiến thức cốt lõi: Muốn tìm hạt còn thiếu phải bảo toàn tổng số khối $A$ và tổng điện tích $Z$ ở hai vế
\item (Sai) Chỉ cần bảo toàn số khối, không cần bảo toàn điện tích.
\item (Đúng) Khi giải cần đọc đúng dữ kiện, dùng hệ đơn vị thống nhất và kiểm tra điều kiện áp dụng.
\item (Sai) Có thể bỏ qua dữ kiện, đơn vị và ý nghĩa vật lí của kết quả.
\end{itemchoice}
}
\end{ex}

% ===== Câu 58 =====
% ID: L12C4B22-15-TL
% Mức: VD
\begin{bt}
% ID: L12C4B22-15-TL
% Mức: VD
Phân tích sai lầm trong nhận định: “Chỉ cần bảo toàn số khối, không cần bảo toàn điện tích.”
\loigiai{
Cần nêu được: Muốn tìm hạt còn thiếu phải bảo toàn tổng số khối $A$ và tổng điện tích $Z$ ở hai vế. Khi vận dụng phải xác định đúng dữ kiện, điều kiện áp dụng, hệ đơn vị và kiểm tra tính hợp lí của kết quả. Nhận định “Chỉ cần bảo toàn số khối, không cần bảo toàn điện tích.” sai vì trái với kiến thức trên.
}
\end{bt}

% ===== Câu 59 =====
% ID: L12C4B22-15-TLN
% Mức: TH
\begin{ex}
% ID: L12C4B22-15-TLN
% Mức: TH
Trong bốn phát biểu sau về hoàn thành phương trình phản ứng hạt nhân, có bao nhiêu phát biểu đúng? Chỉ ghi một số nguyên. (1) Cần đổi các đại lượng về hệ đơn vị thống nhất. (2) Muốn tìm hạt còn thiếu phải bảo toàn tổng số khối $A$ và tổng điện tích $Z$ ở hai vế. (3) Chỉ cần bảo toàn số khối, không cần bảo toàn điện tích. (4) Không cần kiểm tra dữ kiện.
\shortans{2}
\loigiai{
Đối chiếu từng phát biểu với kiến thức cốt lõi: Muốn tìm hạt còn thiếu phải bảo toàn tổng số khối $A$ và tổng điện tích $Z$ ở hai vế. Có 2 phát biểu đúng.
}
\end{ex}

% ===== Câu 60 =====
% ID: L12C4B22-15-TN
% Mức: TH
\begin{ex}
% ID: L12C4B22-15-TN
% Mức: TH
Khi phân tích, phát biểu nào đúng về hoàn thành phương trình phản ứng hạt nhân?
\choice
{\True Muốn tìm hạt còn thiếu phải bảo toàn tổng số khối $A$ và tổng điện tích $Z$ ở hai vế.}
{Chỉ cần bảo toàn số khối, không cần bảo toàn điện tích.}
{Có thể bỏ qua điều kiện áp dụng và đơn vị trong mọi trường hợp.}
{Kết quả của bài toán không phụ thuộc dữ kiện đã cho.}
\loigiai{
Muốn tìm hạt còn thiếu phải bảo toàn tổng số khối $A$ và tổng điện tích $Z$ ở hai vế. Vì vậy phương án $A$ đúng; các phương án còn lại trái với kiến thức hoặc điều kiện áp dụng.
}
\end{ex}

% ===== Câu 61 =====
% ID: L12C4B22-16-DS
% Mức: TH
\begin{ex}
% ID: L12C4B22-16-DS
% Mức: TH
Xét hoàn thành phương trình phản ứng hạt nhân (tình huống 2). Hãy xác định tính đúng, sai của bốn phát biểu sau.
\choiceTF
{Chỉ cần bảo toàn số khối, không cần bảo toàn điện tích.}
{\True Muốn tìm hạt còn thiếu phải bảo toàn tổng số khối $A$ và tổng điện tích $Z$ ở hai vế.}
{Có thể bỏ qua dữ kiện, đơn vị và ý nghĩa vật lí của kết quả.}
{\True Khi giải cần đọc đúng dữ kiện, dùng hệ đơn vị thống nhất và kiểm tra điều kiện áp dụng.}
\loigiai{
\begin{itemchoice}
\item (Sai) Chỉ cần bảo toàn số khối, không cần bảo toàn điện tích. (Vì): Chỉ cần bảo toàn số khối, không cần bảo toàn điện tích. b) Đúng: Muốn tìm hạt còn thiếu phải bảo toàn tổng số khối $A$ và tổng điện tích $Z$ ở hai vế. c) Sai: Có thể bỏ qua dữ kiện, đơn vị và ý nghĩa vật lí của kết quả. d) Đúng: Khi giải cần đọc đúng dữ kiện, dùng hệ đơn vị thống nhất và kiểm tra điều kiện áp dụng. Kiến thức cốt lõi: Muốn tìm hạt còn thiếu phải bảo toàn tổng số khối $A$ và tổng điện tích $Z$ ở hai vế
\item (Đúng) Muốn tìm hạt còn thiếu phải bảo toàn tổng số khối $A$ và tổng điện tích $Z$ ở hai vế.
\item (Sai) Có thể bỏ qua dữ kiện, đơn vị và ý nghĩa vật lí của kết quả.
\item (Đúng) Khi giải cần đọc đúng dữ kiện, dùng hệ đơn vị thống nhất và kiểm tra điều kiện áp dụng.
\end{itemchoice}
}
\end{ex}

% ===== Câu 62 =====
% ID: L12C4B22-16-TL
% Mức: VD
\begin{bt}
% ID: L12C4B22-16-TL
% Mức: VD
Đề xuất các bước giải một bài toán thuộc dạng hoàn thành phương trình phản ứng hạt nhân.
\loigiai{
Cần nêu được: Muốn tìm hạt còn thiếu phải bảo toàn tổng số khối $A$ và tổng điện tích $Z$ ở hai vế. Khi vận dụng phải xác định đúng dữ kiện, điều kiện áp dụng, hệ đơn vị và kiểm tra tính hợp lí của kết quả. Nhận định “Chỉ cần bảo toàn số khối, không cần bảo toàn điện tích.” sai vì trái với kiến thức trên.
}
\end{bt}

% ===== Câu 63 =====
% ID: L12C4B22-16-TLN
% Mức: VD
\begin{ex}
% ID: L12C4B22-16-TLN
% Mức: VD
Trong bốn phát biểu sau về hoàn thành phương trình phản ứng hạt nhân, có bao nhiêu phát biểu đúng? Chỉ ghi một số nguyên. (1) Muốn tìm hạt còn thiếu phải bảo toàn tổng số khối $A$ và tổng điện tích $Z$ ở hai vế. (2) Cần đối chiếu kết quả với điều kiện bài toán. (3) Chỉ cần bảo toàn số khối, không cần bảo toàn điện tích. (4) Mọi trường hợp đều cho cùng một kết quả.
\shortans{2}
\loigiai{
Đối chiếu từng phát biểu với kiến thức cốt lõi: Muốn tìm hạt còn thiếu phải bảo toàn tổng số khối $A$ và tổng điện tích $Z$ ở hai vế. Có 2 phát biểu đúng.
}
\end{ex}

% ===== Câu 64 =====
% ID: L12C4B22-16-TN
% Mức: TH
\begin{ex}
% ID: L12C4B22-16-TN
% Mức: TH
Kết luận nào của học sinh là đúng về hoàn thành phương trình phản ứng hạt nhân?
\choice
{Chỉ cần bảo toàn số khối, không cần bảo toàn điện tích.}
{\True Muốn tìm hạt còn thiếu phải bảo toàn tổng số khối $A$ và tổng điện tích $Z$ ở hai vế.}
{Có thể bỏ qua điều kiện áp dụng và đơn vị trong mọi trường hợp.}
{Kết quả của bài toán không phụ thuộc dữ kiện đã cho.}
\loigiai{
Muốn tìm hạt còn thiếu phải bảo toàn tổng số khối $A$ và tổng điện tích $Z$ ở hai vế. Vì vậy phương án $B$ đúng; các phương án còn lại trái với kiến thức hoặc điều kiện áp dụng.
}
\end{ex}

% ===== Câu 65 =====
% ID: L12C4B22-17-DS
% Mức: VD
\begin{ex}
% ID: L12C4B22-17-DS
% Mức: VD
Xét hoàn thành phương trình phản ứng hạt nhân (tình huống 3). Hãy xác định tính đúng, sai của bốn phát biểu sau.
\choiceTF
{\True Muốn tìm hạt còn thiếu phải bảo toàn tổng số khối $A$ và tổng điện tích $Z$ ở hai vế.}
{Chỉ cần bảo toàn số khối, không cần bảo toàn điện tích.}
{\True Khi giải cần đọc đúng dữ kiện, dùng hệ đơn vị thống nhất và kiểm tra điều kiện áp dụng.}
{Có thể bỏ qua dữ kiện, đơn vị và ý nghĩa vật lí của kết quả.}
\loigiai{
\begin{itemchoice}
\item (Đúng) Muốn tìm hạt còn thiếu phải bảo toàn tổng số khối $A$ và tổng điện tích $Z$ ở hai vế. (Vì): Muốn tìm hạt còn thiếu phải bảo toàn tổng số khối $A$ và tổng điện tích $Z$ ở hai vế. b) Sai: Chỉ cần bảo toàn số khối, không cần bảo toàn điện tích. c) Đúng: Khi giải cần đọc đúng dữ kiện, dùng hệ đơn vị thống nhất và kiểm tra điều kiện áp dụng. d) Sai: Có thể bỏ qua dữ kiện, đơn vị và ý nghĩa vật lí của kết quả. Kiến thức cốt lõi: Muốn tìm hạt còn thiếu phải bảo toàn tổng số khối $A$ và tổng điện tích $Z$ ở hai vế
\item (Sai) Chỉ cần bảo toàn số khối, không cần bảo toàn điện tích.
\item (Đúng) Khi giải cần đọc đúng dữ kiện, dùng hệ đơn vị thống nhất và kiểm tra điều kiện áp dụng.
\item (Sai) Có thể bỏ qua dữ kiện, đơn vị và ý nghĩa vật lí của kết quả.
\end{itemchoice}
}
\end{ex}

% ===== Câu 66 =====
% ID: L12C4B22-17-TL
% Mức: VD
\begin{bt}
% ID: L12C4B22-17-TL
% Mức: VD
Nêu một tình huống thực tiễn liên quan đến hoàn thành phương trình phản ứng hạt nhân và giải thích bằng kiến thức Vật lí.
\loigiai{
Cần nêu được: Muốn tìm hạt còn thiếu phải bảo toàn tổng số khối $A$ và tổng điện tích $Z$ ở hai vế. Khi vận dụng phải xác định đúng dữ kiện, điều kiện áp dụng, hệ đơn vị và kiểm tra tính hợp lí của kết quả. Nhận định “Chỉ cần bảo toàn số khối, không cần bảo toàn điện tích.” sai vì trái với kiến thức trên.
}
\end{bt}

% ===== Câu 67 =====
% ID: L12C4B22-17-TLN
% Mức: VD
\begin{ex}
% ID: L12C4B22-17-TLN
% Mức: VD
Trong bốn phát biểu sau về hoàn thành phương trình phản ứng hạt nhân, có bao nhiêu phát biểu đúng? Chỉ ghi một số nguyên. (1) Chỉ cần bảo toàn số khối, không cần bảo toàn điện tích. (2) Cần đọc đúng dữ kiện và giả thiết. (3) Muốn tìm hạt còn thiếu phải bảo toàn tổng số khối $A$ và tổng điện tích $Z$ ở hai vế. (4) Có thể thay số trước khi chọn công thức.
\shortans{2}
\loigiai{
Đối chiếu từng phát biểu với kiến thức cốt lõi: Muốn tìm hạt còn thiếu phải bảo toàn tổng số khối $A$ và tổng điện tích $Z$ ở hai vế. Có 2 phát biểu đúng.
}
\end{ex}

% ===== Câu 68 =====
% ID: L12C4B22-17-TN
% Mức: TH
\begin{ex}
% ID: L12C4B22-17-TN
% Mức: TH
Chọn phát biểu không trái với kiến thức về hoàn thành phương trình phản ứng hạt nhân?
\choice
{Chỉ cần bảo toàn số khối, không cần bảo toàn điện tích.}
{Có thể bỏ qua điều kiện áp dụng và đơn vị trong mọi trường hợp.}
{\True Muốn tìm hạt còn thiếu phải bảo toàn tổng số khối $A$ và tổng điện tích $Z$ ở hai vế.}
{Kết quả của bài toán không phụ thuộc dữ kiện đã cho.}
\loigiai{
Muốn tìm hạt còn thiếu phải bảo toàn tổng số khối $A$ và tổng điện tích $Z$ ở hai vế. Vì vậy phương án $C$ đúng; các phương án còn lại trái với kiến thức hoặc điều kiện áp dụng.
}
\end{ex}

% ===== Câu 69 =====
% ID: L12C4B22-18-DS
% Mức: VD
\begin{ex}
% ID: L12C4B22-18-DS
% Mức: VD
Xét hoàn thành phương trình phản ứng hạt nhân (tình huống 4). Hãy xác định tính đúng, sai của bốn phát biểu sau.
\choiceTF
{Chỉ cần bảo toàn số khối, không cần bảo toàn điện tích.}
{\True Muốn tìm hạt còn thiếu phải bảo toàn tổng số khối $A$ và tổng điện tích $Z$ ở hai vế.}
{Có thể bỏ qua dữ kiện, đơn vị và ý nghĩa vật lí của kết quả.}
{\True Khi giải cần đọc đúng dữ kiện, dùng hệ đơn vị thống nhất và kiểm tra điều kiện áp dụng.}
\loigiai{
\begin{itemchoice}
\item (Sai) Chỉ cần bảo toàn số khối, không cần bảo toàn điện tích. (Vì): Chỉ cần bảo toàn số khối, không cần bảo toàn điện tích. b) Đúng: Muốn tìm hạt còn thiếu phải bảo toàn tổng số khối $A$ và tổng điện tích $Z$ ở hai vế. c) Sai: Có thể bỏ qua dữ kiện, đơn vị và ý nghĩa vật lí của kết quả. d) Đúng: Khi giải cần đọc đúng dữ kiện, dùng hệ đơn vị thống nhất và kiểm tra điều kiện áp dụng. Kiến thức cốt lõi: Muốn tìm hạt còn thiếu phải bảo toàn tổng số khối $A$ và tổng điện tích $Z$ ở hai vế
\item (Đúng) Muốn tìm hạt còn thiếu phải bảo toàn tổng số khối $A$ và tổng điện tích $Z$ ở hai vế.
\item (Sai) Có thể bỏ qua dữ kiện, đơn vị và ý nghĩa vật lí của kết quả.
\item (Đúng) Khi giải cần đọc đúng dữ kiện, dùng hệ đơn vị thống nhất và kiểm tra điều kiện áp dụng.
\end{itemchoice}
}
\end{ex}

% ===== Câu 70 =====
% ID: L12C4B22-18-TL
% Mức: VDC
\begin{bt}
% ID: L12C4B22-18-TL
% Mức: VDC
So sánh nhận định đúng “Muốn tìm hạt còn thiếu phải bảo toàn tổng số khối $A$ và tổng điện tích $Z$ ở hai vế.” với nhận định sai “Chỉ cần bảo toàn số khối, không cần bảo toàn điện tích.”; chỉ rõ căn cứ Vật lí.
\loigiai{
Cần nêu được: Muốn tìm hạt còn thiếu phải bảo toàn tổng số khối $A$ và tổng điện tích $Z$ ở hai vế. Khi vận dụng phải xác định đúng dữ kiện, điều kiện áp dụng, hệ đơn vị và kiểm tra tính hợp lí của kết quả. Nhận định “Chỉ cần bảo toàn số khối, không cần bảo toàn điện tích.” sai vì trái với kiến thức trên.
}
\end{bt}

% ===== Câu 71 =====
% ID: L12C4B22-18-TLN
% Mức: VDC
\begin{ex}
% ID: L12C4B22-18-TLN
% Mức: VDC
Trong bốn phát biểu sau về hoàn thành phương trình phản ứng hạt nhân, có bao nhiêu phát biểu đúng? Chỉ ghi một số nguyên. (1) Kết quả cần có ý nghĩa vật lí. (2) Chỉ cần bảo toàn số khối, không cần bảo toàn điện tích. (3) Phải dùng đúng định luật cho tình huống. (4) Muốn tìm hạt còn thiếu phải bảo toàn tổng số khối $A$ và tổng điện tích $Z$ ở hai vế.
\shortans{3}
\loigiai{
Đối chiếu từng phát biểu với kiến thức cốt lõi: Muốn tìm hạt còn thiếu phải bảo toàn tổng số khối $A$ và tổng điện tích $Z$ ở hai vế. Có 3 phát biểu đúng.
}
\end{ex}

% ===== Câu 72 =====
% ID: L12C4B22-18-TN
% Mức: VD
\begin{ex}
% ID: L12C4B22-18-TN
% Mức: VD
Cơ sở Vật lí đúng để giải bài là gì về hoàn thành phương trình phản ứng hạt nhân?
\choice
{Chỉ cần bảo toàn số khối, không cần bảo toàn điện tích.}
{Có thể bỏ qua điều kiện áp dụng và đơn vị trong mọi trường hợp.}
{Kết quả của bài toán không phụ thuộc dữ kiện đã cho.}
{\True Muốn tìm hạt còn thiếu phải bảo toàn tổng số khối $A$ và tổng điện tích $Z$ ở hai vế.}
\loigiai{
Muốn tìm hạt còn thiếu phải bảo toàn tổng số khối $A$ và tổng điện tích $Z$ ở hai vế. Vì vậy phương án $D$ đúng; các phương án còn lại trái với kiến thức hoặc điều kiện áp dụng.
}
\end{ex}

% ===== Câu 73 =====
% ID: L12C4B22-19-DS
% Mức: VD
\begin{ex}
% ID: L12C4B22-19-DS
% Mức: VD
Xét hoàn thành phương trình phản ứng hạt nhân (tình huống 5). Hãy xác định tính đúng, sai của bốn phát biểu sau.
\choiceTF
{\True Muốn tìm hạt còn thiếu phải bảo toàn tổng số khối $A$ và tổng điện tích $Z$ ở hai vế.}
{Có thể bỏ qua dữ kiện, đơn vị và ý nghĩa vật lí của kết quả.}
{Chỉ cần bảo toàn số khối, không cần bảo toàn điện tích.}
{Có thể bỏ qua dữ kiện, đơn vị và ý nghĩa vật lí của kết quả.}
\loigiai{
\begin{itemchoice}
\item (Đúng) Muốn tìm hạt còn thiếu phải bảo toàn tổng số khối $A$ và tổng điện tích $Z$ ở hai vế. (Vì): Muốn tìm hạt còn thiếu phải bảo toàn tổng số khối $A$ và tổng điện tích $Z$ ở hai vế. b) Sai: Có thể bỏ qua dữ kiện, đơn vị và ý nghĩa vật lí của kết quả. c) Sai: Chỉ cần bảo toàn số khối, không cần bảo toàn điện tích. d) Sai: Có thể bỏ qua dữ kiện, đơn vị và ý nghĩa vật lí của kết quả. Kiến thức cốt lõi: Muốn tìm hạt còn thiếu phải bảo toàn tổng số khối $A$ và tổng điện tích $Z$ ở hai vế
\item (Sai) Có thể bỏ qua dữ kiện, đơn vị và ý nghĩa vật lí của kết quả.
\item (Sai) Chỉ cần bảo toàn số khối, không cần bảo toàn điện tích.
\item (Sai) Có thể bỏ qua dữ kiện, đơn vị và ý nghĩa vật lí của kết quả.
\end{itemchoice}
}
\end{ex}

% ===== Câu 74 =====
% ID: L12C4B22-19-TL
% Mức: TH
\dangbt{Tính độ hụt khối, năng lượng liên kết và năng lượng liên kết riêng}
\begin{bt}
% ID: L12C4B22-19-TL
% Mức: TH
Trình bày kiến thức cốt lõi của tính độ hụt khối, năng lượng liên kết và năng lượng liên kết riêng và nêu điều kiện áp dụng.
\loigiai{
Cần nêu được: Năng lượng liên kết $W_{lk}=\Delta m c^2$; năng lượng liên kết riêng bằng $W_{lk}/A$. Khi vận dụng phải xác định đúng dữ kiện, điều kiện áp dụng, hệ đơn vị và kiểm tra tính hợp lí của kết quả. Nhận định “Năng lượng liên kết riêng bằng năng lượng liên kết nhân với số khối.” sai vì trái với kiến thức trên.
}
\end{bt}

% ===== Câu 75 =====
% ID: L12C4B22-19-TLN
% Mức: TH
\begin{ex}
% ID: L12C4B22-19-TLN
% Mức: TH
Trong bốn phát biểu sau về tính độ hụt khối, năng lượng liên kết và năng lượng liên kết riêng, có bao nhiêu phát biểu đúng? Chỉ ghi một số nguyên. (1) Năng lượng liên kết $W_{lk}=\Delta m c^2$; năng lượng liên kết riêng bằng $W_{lk}/A$. (2) Năng lượng liên kết riêng bằng năng lượng liên kết nhân với số khối. (3) Cần kiểm tra điều kiện áp dụng trước khi tính toán. (4) Có thể bỏ qua đơn vị của kết quả.
\shortans{2}
\loigiai{
Đối chiếu từng phát biểu với kiến thức cốt lõi: Năng lượng liên kết $W_{lk}=\Delta m c^2$; năng lượng liên kết riêng bằng $W_{lk}/A$. Có 2 phát biểu đúng.
}
\end{ex}

% ===== Câu 76 =====
% ID: L12C4B22-19-TN
% Mức: VD
\dangbt{Hoàn thành phương trình phản ứng hạt nhân}
\begin{ex}
% ID: L12C4B22-19-TN
% Mức: VD
Trong một bài thực hành, nhận xét nào đúng về hoàn thành phương trình phản ứng hạt nhân?
\choice
{\True Muốn tìm hạt còn thiếu phải bảo toàn tổng số khối $A$ và tổng điện tích $Z$ ở hai vế.}
{Chỉ cần bảo toàn số khối, không cần bảo toàn điện tích.}
{Có thể bỏ qua điều kiện áp dụng và đơn vị trong mọi trường hợp.}
{Kết quả của bài toán không phụ thuộc dữ kiện đã cho.}
\loigiai{
Muốn tìm hạt còn thiếu phải bảo toàn tổng số khối $A$ và tổng điện tích $Z$ ở hai vế. Vì vậy phương án $A$ đúng; các phương án còn lại trái với kiến thức hoặc điều kiện áp dụng.
}
\end{ex}

% ===== Câu 77 =====
% ID: L12C4B22-20-DS
% Mức: TH
\dangbt{Tính độ hụt khối, năng lượng liên kết và năng lượng liên kết riêng}
\begin{ex}
% ID: L12C4B22-20-DS
% Mức: TH
Xét tính độ hụt khối, năng lượng liên kết và năng lượng liên kết riêng (tình huống 1). Hãy xác định tính đúng, sai của bốn phát biểu sau.
\choiceTF
{\True Năng lượng liên kết $W_{lk}=\Delta m c^2$; năng lượng liên kết riêng bằng $W_{lk}/A$.}
{Năng lượng liên kết riêng bằng năng lượng liên kết nhân với số khối.}
{\True Khi giải cần đọc đúng dữ kiện, dùng hệ đơn vị thống nhất và kiểm tra điều kiện áp dụng.}
{Có thể bỏ qua dữ kiện, đơn vị và ý nghĩa vật lí của kết quả.}
\loigiai{
\begin{itemchoice}
\item (Đúng) Năng lượng liên kết $W_{lk}=\Delta m c^2$; năng lượng liên kết riêng bằng $W_{lk}/A$. (Vì): Năng lượng liên kết $W_{lk}=\Delta m c^2$; năng lượng liên kết riêng bằng $W_{lk}/A$. b) Sai: Năng lượng liên kết riêng bằng năng lượng liên kết nhân với số khối. c) Đúng: Khi giải cần đọc đúng dữ kiện, dùng hệ đơn vị thống nhất và kiểm tra điều kiện áp dụng. d) Sai: Có thể bỏ qua dữ kiện, đơn vị và ý nghĩa vật lí của kết quả. Kiến thức cốt lõi: Năng lượng liên kết $W_{lk}=\Delta m c^2$; năng lượng liên kết riêng bằng $W_{lk}/A$
\item (Sai) Năng lượng liên kết riêng bằng năng lượng liên kết nhân với số khối.
\item (Đúng) Khi giải cần đọc đúng dữ kiện, dùng hệ đơn vị thống nhất và kiểm tra điều kiện áp dụng.
\item (Sai) Có thể bỏ qua dữ kiện, đơn vị và ý nghĩa vật lí của kết quả.
\end{itemchoice}
}
\end{ex}

% ===== Câu 78 =====
% ID: L12C4B22-20-TL
% Mức: TH
\begin{bt}
% ID: L12C4B22-20-TL
% Mức: TH
Giải thích vì sao nhận định sau đúng: “Năng lượng liên kết $W_{lk}=\Delta m c^2$; năng lượng liên kết riêng bằng $W_{lk}/A$.”
\loigiai{
Cần nêu được: Năng lượng liên kết $W_{lk}=\Delta m c^2$; năng lượng liên kết riêng bằng $W_{lk}/A$. Khi vận dụng phải xác định đúng dữ kiện, điều kiện áp dụng, hệ đơn vị và kiểm tra tính hợp lí của kết quả. Nhận định “Năng lượng liên kết riêng bằng năng lượng liên kết nhân với số khối.” sai vì trái với kiến thức trên.
}
\end{bt}

% ===== Câu 79 =====
% ID: L12C4B22-20-TLN
% Mức: TH
\begin{ex}
% ID: L12C4B22-20-TLN
% Mức: TH
Trong bốn phát biểu sau về tính độ hụt khối, năng lượng liên kết và năng lượng liên kết riêng, có bao nhiêu phát biểu đúng? Chỉ ghi một số nguyên. (1) Năng lượng liên kết riêng bằng năng lượng liên kết nhân với số khối. (2) Năng lượng liên kết $W_{lk}=\Delta m c^2$; năng lượng liên kết riêng bằng $W_{lk}/A$. (3) Kết quả phải phù hợp ý nghĩa vật lí. (4) Mọi đại lượng đều có cùng bản chất.
\shortans{2}
\loigiai{
Đối chiếu từng phát biểu với kiến thức cốt lõi: Năng lượng liên kết $W_{lk}=\Delta m c^2$; năng lượng liên kết riêng bằng $W_{lk}/A$. Có 2 phát biểu đúng.
}
\end{ex}

% ===== Câu 80 =====
% ID: L12C4B22-20-TN
% Mức: VD
\dangbt{Hoàn thành phương trình phản ứng hạt nhân}
\begin{ex}
% ID: L12C4B22-20-TN
% Mức: VD
Kết luận đúng khi vận dụng là về hoàn thành phương trình phản ứng hạt nhân?
\choice
{Chỉ cần bảo toàn số khối, không cần bảo toàn điện tích.}
{\True Muốn tìm hạt còn thiếu phải bảo toàn tổng số khối $A$ và tổng điện tích $Z$ ở hai vế.}
{Có thể bỏ qua điều kiện áp dụng và đơn vị trong mọi trường hợp.}
{Kết quả của bài toán không phụ thuộc dữ kiện đã cho.}
\loigiai{
Muốn tìm hạt còn thiếu phải bảo toàn tổng số khối $A$ và tổng điện tích $Z$ ở hai vế. Vì vậy phương án $B$ đúng; các phương án còn lại trái với kiến thức hoặc điều kiện áp dụng.
}
\end{ex}

% ===== Câu 81 =====
% ID: L12C4B22-21-DS
% Mức: TH
\dangbt{Tính độ hụt khối, năng lượng liên kết và năng lượng liên kết riêng}
\begin{ex}
% ID: L12C4B22-21-DS
% Mức: TH
Xét tính độ hụt khối, năng lượng liên kết và năng lượng liên kết riêng (tình huống 2). Hãy xác định tính đúng, sai của bốn phát biểu sau.
\choiceTF
{Năng lượng liên kết riêng bằng năng lượng liên kết nhân với số khối.}
{\True Năng lượng liên kết $W_{lk}=\Delta m c^2$; năng lượng liên kết riêng bằng $W_{lk}/A$.}
{Có thể bỏ qua dữ kiện, đơn vị và ý nghĩa vật lí của kết quả.}
{\True Khi giải cần đọc đúng dữ kiện, dùng hệ đơn vị thống nhất và kiểm tra điều kiện áp dụng.}
\loigiai{
\begin{itemchoice}
\item (Sai) Năng lượng liên kết riêng bằng năng lượng liên kết nhân với số khối. (Vì): Năng lượng liên kết riêng bằng năng lượng liên kết nhân với số khối. b) Đúng: Năng lượng liên kết $W_{lk}=\Delta m c^2$; năng lượng liên kết riêng bằng $W_{lk}/A$. c) Sai: Có thể bỏ qua dữ kiện, đơn vị và ý nghĩa vật lí của kết quả. d) Đúng: Khi giải cần đọc đúng dữ kiện, dùng hệ đơn vị thống nhất và kiểm tra điều kiện áp dụng. Kiến thức cốt lõi: Năng lượng liên kết $W_{lk}=\Delta m c^2$; năng lượng liên kết riêng bằng $W_{lk}/A$
\item (Đúng) Năng lượng liên kết $W_{lk}=\Delta m c^2$; năng lượng liên kết riêng bằng $W_{lk}/A$.
\item (Sai) Có thể bỏ qua dữ kiện, đơn vị và ý nghĩa vật lí của kết quả.
\item (Đúng) Khi giải cần đọc đúng dữ kiện, dùng hệ đơn vị thống nhất và kiểm tra điều kiện áp dụng.
\end{itemchoice}
}
\end{ex}

% ===== Câu 82 =====
% ID: L12C4B22-21-TL
% Mức: VD
\begin{bt}
% ID: L12C4B22-21-TL
% Mức: VD
Phân tích sai lầm trong nhận định: “Năng lượng liên kết riêng bằng năng lượng liên kết nhân với số khối.”
\loigiai{
Cần nêu được: Năng lượng liên kết $W_{lk}=\Delta m c^2$; năng lượng liên kết riêng bằng $W_{lk}/A$. Khi vận dụng phải xác định đúng dữ kiện, điều kiện áp dụng, hệ đơn vị và kiểm tra tính hợp lí của kết quả. Nhận định “Năng lượng liên kết riêng bằng năng lượng liên kết nhân với số khối.” sai vì trái với kiến thức trên.
}
\end{bt}

% ===== Câu 83 =====
% ID: L12C4B22-21-TLN
% Mức: TH
\begin{ex}
% ID: L12C4B22-21-TLN
% Mức: TH
Trong bốn phát biểu sau về tính độ hụt khối, năng lượng liên kết và năng lượng liên kết riêng, có bao nhiêu phát biểu đúng? Chỉ ghi một số nguyên. (1) Cần đổi các đại lượng về hệ đơn vị thống nhất. (2) Năng lượng liên kết $W_{lk}=\Delta m c^2$; năng lượng liên kết riêng bằng $W_{lk}/A$. (3) Năng lượng liên kết riêng bằng năng lượng liên kết nhân với số khối. (4) Không cần kiểm tra dữ kiện.
\shortans{2}
\loigiai{
Đối chiếu từng phát biểu với kiến thức cốt lõi: Năng lượng liên kết $W_{lk}=\Delta m c^2$; năng lượng liên kết riêng bằng $W_{lk}/A$. Có 2 phát biểu đúng.
}
\end{ex}

% ===== Câu 84 =====
% ID: L12C4B22-21-TN
% Mức: NB
\dangbt{Nhận biết từ trường, nguồn gây ra từ trường và tương tác từ}
\begin{ex}
% ID: L12C4B22-21-TN
% Mức: NB
Phát biểu nào sau đây đúng về nhận biết từ trường, nguồn gây ra từ trường và tương tác từ?
\choice
{\True Từ trường tồn tại xung quanh nam châm và dòng điện; nó gây lực từ lên nam châm hoặc dòng điện khác đặt trong đó.}
{Mọi điện tích đứng yên đều tạo ra từ trường giống dòng điện.}
{Đại lượng đang xét không phụ thuộc các điều kiện vật lí nêu trong bài.}
{Có thể bỏ qua phương, chiều và đơn vị của các đại lượng trong mọi trường hợp.}
\loigiai{
Từ trường tồn tại xung quanh nam châm và dòng điện; nó gây lực từ lên nam châm hoặc dòng điện khác đặt trong đó. Vì vậy phương án $A$ đúng; các phương án còn lại trái với định nghĩa, định luật hoặc điều kiện áp dụng.
}
\end{ex}

% ===== Câu 85 =====
% ID: L12C4B22-22-DS
% Mức: VD
\dangbt{Tính độ hụt khối, năng lượng liên kết và năng lượng liên kết riêng}
\begin{ex}
% ID: L12C4B22-22-DS
% Mức: VD
Xét tính độ hụt khối, năng lượng liên kết và năng lượng liên kết riêng (tình huống 3). Hãy xác định tính đúng, sai của bốn phát biểu sau.
\choiceTF
{\True Năng lượng liên kết $W_{lk}=\Delta m c^2$; năng lượng liên kết riêng bằng $W_{lk}/A$.}
{Năng lượng liên kết riêng bằng năng lượng liên kết nhân với số khối.}
{\True Khi giải cần đọc đúng dữ kiện, dùng hệ đơn vị thống nhất và kiểm tra điều kiện áp dụng.}
{Có thể bỏ qua dữ kiện, đơn vị và ý nghĩa vật lí của kết quả.}
\loigiai{
\begin{itemchoice}
\item (Đúng) Năng lượng liên kết $W_{lk}=\Delta m c^2$; năng lượng liên kết riêng bằng $W_{lk}/A$. (Vì): Năng lượng liên kết $W_{lk}=\Delta m c^2$; năng lượng liên kết riêng bằng $W_{lk}/A$. b) Sai: Năng lượng liên kết riêng bằng năng lượng liên kết nhân với số khối. c) Đúng: Khi giải cần đọc đúng dữ kiện, dùng hệ đơn vị thống nhất và kiểm tra điều kiện áp dụng. d) Sai: Có thể bỏ qua dữ kiện, đơn vị và ý nghĩa vật lí của kết quả. Kiến thức cốt lõi: Năng lượng liên kết $W_{lk}=\Delta m c^2$; năng lượng liên kết riêng bằng $W_{lk}/A$
\item (Sai) Năng lượng liên kết riêng bằng năng lượng liên kết nhân với số khối.
\item (Đúng) Khi giải cần đọc đúng dữ kiện, dùng hệ đơn vị thống nhất và kiểm tra điều kiện áp dụng.
\item (Sai) Có thể bỏ qua dữ kiện, đơn vị và ý nghĩa vật lí của kết quả.
\end{itemchoice}
}
\end{ex}

% ===== Câu 86 =====
% ID: L12C4B22-22-TL
% Mức: VD
\begin{bt}
% ID: L12C4B22-22-TL
% Mức: VD
Đề xuất các bước giải một bài toán thuộc dạng tính độ hụt khối, năng lượng liên kết và năng lượng liên kết riêng.
\loigiai{
Cần nêu được: Năng lượng liên kết $W_{lk}=\Delta m c^2$; năng lượng liên kết riêng bằng $W_{lk}/A$. Khi vận dụng phải xác định đúng dữ kiện, điều kiện áp dụng, hệ đơn vị và kiểm tra tính hợp lí của kết quả. Nhận định “Năng lượng liên kết riêng bằng năng lượng liên kết nhân với số khối.” sai vì trái với kiến thức trên.
}
\end{bt}

% ===== Câu 87 =====
% ID: L12C4B22-22-TLN
% Mức: VD
\begin{ex}
% ID: L12C4B22-22-TLN
% Mức: VD
Trong bốn phát biểu sau về tính độ hụt khối, năng lượng liên kết và năng lượng liên kết riêng, có bao nhiêu phát biểu đúng? Chỉ ghi một số nguyên. (1) Năng lượng liên kết $W_{lk}=\Delta m c^2$; năng lượng liên kết riêng bằng $W_{lk}/A$. (2) Cần đối chiếu kết quả với điều kiện bài toán. (3) Năng lượng liên kết riêng bằng năng lượng liên kết nhân với số khối. (4) Mọi trường hợp đều cho cùng một kết quả.
\shortans{2}
\loigiai{
Đối chiếu từng phát biểu với kiến thức cốt lõi: Năng lượng liên kết $W_{lk}=\Delta m c^2$; năng lượng liên kết riêng bằng $W_{lk}/A$. Có 2 phát biểu đúng.
}
\end{ex}

% ===== Câu 88 =====
% ID: L12C4B22-22-TN
% Mức: NB
\dangbt{Nhận biết phản ứng hạt nhân và các định luật bảo toàn}
\begin{ex}
% ID: L12C4B22-22-TN
% Mức: NB
Phát biểu nào sau đây đúng về nhận biết phản ứng hạt nhân và các định luật bảo toàn?
\choice
{\True Trong phản ứng hạt nhân bảo toàn điện tích, số nucleon, năng lượng toàn phần và động lượng.}
{Trong phản ứng hạt nhân luôn bảo toàn khối lượng nghỉ.}
{Có thể bỏ qua điều kiện áp dụng và đơn vị trong mọi trường hợp.}
{Kết quả của bài toán không phụ thuộc dữ kiện đã cho.}
\loigiai{
Trong phản ứng hạt nhân bảo toàn điện tích, số nucleon, năng lượng toàn phần và động lượng. Vì vậy phương án $A$ đúng; các phương án còn lại trái với kiến thức hoặc điều kiện áp dụng.
}
\end{ex}

% ===== Câu 89 =====
% ID: L12C4B22-23-DS
% Mức: VD
\dangbt{Tính độ hụt khối, năng lượng liên kết và năng lượng liên kết riêng}
\begin{ex}
% ID: L12C4B22-23-DS
% Mức: VD
Xét tính độ hụt khối, năng lượng liên kết và năng lượng liên kết riêng (tình huống 4). Hãy xác định tính đúng, sai của bốn phát biểu sau.
\choiceTF
{Năng lượng liên kết riêng bằng năng lượng liên kết nhân với số khối.}
{\True Năng lượng liên kết $W_{lk}=\Delta m c^2$; năng lượng liên kết riêng bằng $W_{lk}/A$.}
{Có thể bỏ qua dữ kiện, đơn vị và ý nghĩa vật lí của kết quả.}
{\True Khi giải cần đọc đúng dữ kiện, dùng hệ đơn vị thống nhất và kiểm tra điều kiện áp dụng.}
\loigiai{
\begin{itemchoice}
\item (Sai) Năng lượng liên kết riêng bằng năng lượng liên kết nhân với số khối. (Vì): Năng lượng liên kết riêng bằng năng lượng liên kết nhân với số khối. b) Đúng: Năng lượng liên kết $W_{lk}=\Delta m c^2$; năng lượng liên kết riêng bằng $W_{lk}/A$. c) Sai: Có thể bỏ qua dữ kiện, đơn vị và ý nghĩa vật lí của kết quả. d) Đúng: Khi giải cần đọc đúng dữ kiện, dùng hệ đơn vị thống nhất và kiểm tra điều kiện áp dụng. Kiến thức cốt lõi: Năng lượng liên kết $W_{lk}=\Delta m c^2$; năng lượng liên kết riêng bằng $W_{lk}/A$
\item (Đúng) Năng lượng liên kết $W_{lk}=\Delta m c^2$; năng lượng liên kết riêng bằng $W_{lk}/A$.
\item (Sai) Có thể bỏ qua dữ kiện, đơn vị và ý nghĩa vật lí của kết quả.
\item (Đúng) Khi giải cần đọc đúng dữ kiện, dùng hệ đơn vị thống nhất và kiểm tra điều kiện áp dụng.
\end{itemchoice}
}
\end{ex}

% ===== Câu 90 =====
% ID: L12C4B22-23-TL
% Mức: VD
\begin{bt}
% ID: L12C4B22-23-TL
% Mức: VD
Nêu một tình huống thực tiễn liên quan đến tính độ hụt khối, năng lượng liên kết và năng lượng liên kết riêng và giải thích bằng kiến thức Vật lí.
\loigiai{
Cần nêu được: Năng lượng liên kết $W_{lk}=\Delta m c^2$; năng lượng liên kết riêng bằng $W_{lk}/A$. Khi vận dụng phải xác định đúng dữ kiện, điều kiện áp dụng, hệ đơn vị và kiểm tra tính hợp lí của kết quả. Nhận định “Năng lượng liên kết riêng bằng năng lượng liên kết nhân với số khối.” sai vì trái với kiến thức trên.
}
\end{bt}

% ===== Câu 91 =====
% ID: L12C4B22-23-TLN
% Mức: VD
\begin{ex}
% ID: L12C4B22-23-TLN
% Mức: VD
Trong bốn phát biểu sau về tính độ hụt khối, năng lượng liên kết và năng lượng liên kết riêng, có bao nhiêu phát biểu đúng? Chỉ ghi một số nguyên. (1) Năng lượng liên kết riêng bằng năng lượng liên kết nhân với số khối. (2) Cần đọc đúng dữ kiện và giả thiết. (3) Năng lượng liên kết $W_{lk}=\Delta m c^2$; năng lượng liên kết riêng bằng $W_{lk}/A$. (4) Có thể thay số trước khi chọn công thức.
\shortans{2}
\loigiai{
Đối chiếu từng phát biểu với kiến thức cốt lõi: Năng lượng liên kết $W_{lk}=\Delta m c^2$; năng lượng liên kết riêng bằng $W_{lk}/A$. Có 2 phát biểu đúng.
}
\end{ex}

% ===== Câu 92 =====
% ID: L12C4B22-23-TN
% Mức: NB
\dangbt{Nhận biết phản ứng hạt nhân và các định luật bảo toàn}
\begin{ex}
% ID: L12C4B22-23-TN
% Mức: NB
Nhận định nào phù hợp với kiến thức Vật lí về nhận biết phản ứng hạt nhân và các định luật bảo toàn?
\choice
{Trong phản ứng hạt nhân luôn bảo toàn khối lượng nghỉ.}
{\True Trong phản ứng hạt nhân bảo toàn điện tích, số nucleon, năng lượng toàn phần và động lượng.}
{Có thể bỏ qua điều kiện áp dụng và đơn vị trong mọi trường hợp.}
{Kết quả của bài toán không phụ thuộc dữ kiện đã cho.}
\loigiai{
Trong phản ứng hạt nhân bảo toàn điện tích, số nucleon, năng lượng toàn phần và động lượng. Vì vậy phương án $B$ đúng; các phương án còn lại trái với kiến thức hoặc điều kiện áp dụng.
}
\end{ex}

% ===== Câu 93 =====
% ID: L12C4B22-24-DS
% Mức: VD
\dangbt{Tính độ hụt khối, năng lượng liên kết và năng lượng liên kết riêng}
\begin{ex}
% ID: L12C4B22-24-DS
% Mức: VD
Xét tính độ hụt khối, năng lượng liên kết và năng lượng liên kết riêng (tình huống 5). Hãy xác định tính đúng, sai của bốn phát biểu sau.
\choiceTF
{\True Năng lượng liên kết $W_{lk}=\Delta m c^2$; năng lượng liên kết riêng bằng $W_{lk}/A$.}
{Có thể bỏ qua dữ kiện, đơn vị và ý nghĩa vật lí của kết quả.}
{Năng lượng liên kết riêng bằng năng lượng liên kết nhân với số khối.}
{Có thể bỏ qua dữ kiện, đơn vị và ý nghĩa vật lí của kết quả.}
\loigiai{
\begin{itemchoice}
\item (Đúng) Năng lượng liên kết $W_{lk}=\Delta m c^2$; năng lượng liên kết riêng bằng $W_{lk}/A$. (Vì): Năng lượng liên kết $W_{lk}=\Delta m c^2$; năng lượng liên kết riêng bằng $W_{lk}/A$. b) Sai: Có thể bỏ qua dữ kiện, đơn vị và ý nghĩa vật lí của kết quả. c) Sai: Năng lượng liên kết riêng bằng năng lượng liên kết nhân với số khối. d) Sai: Có thể bỏ qua dữ kiện, đơn vị và ý nghĩa vật lí của kết quả. Kiến thức cốt lõi: Năng lượng liên kết $W_{lk}=\Delta m c^2$; năng lượng liên kết riêng bằng $W_{lk}/A$
\item (Sai) Có thể bỏ qua dữ kiện, đơn vị và ý nghĩa vật lí của kết quả.
\item (Sai) Năng lượng liên kết riêng bằng năng lượng liên kết nhân với số khối.
\item (Sai) Có thể bỏ qua dữ kiện, đơn vị và ý nghĩa vật lí của kết quả.
\end{itemchoice}
}
\end{ex}

% ===== Câu 94 =====
% ID: L12C4B22-24-TL
% Mức: VDC
\begin{bt}
% ID: L12C4B22-24-TL
% Mức: VDC
So sánh nhận định đúng “Năng lượng liên kết $W_{lk}=\Delta m c^2$; năng lượng liên kết riêng bằng $W_{lk}/A$.” với nhận định sai “Năng lượng liên kết riêng bằng năng lượng liên kết nhân với số khối.”; chỉ rõ căn cứ Vật lí.
\loigiai{
Cần nêu được: Năng lượng liên kết $W_{lk}=\Delta m c^2$; năng lượng liên kết riêng bằng $W_{lk}/A$. Khi vận dụng phải xác định đúng dữ kiện, điều kiện áp dụng, hệ đơn vị và kiểm tra tính hợp lí của kết quả. Nhận định “Năng lượng liên kết riêng bằng năng lượng liên kết nhân với số khối.” sai vì trái với kiến thức trên.
}
\end{bt}

% ===== Câu 95 =====
% ID: L12C4B22-24-TLN
% Mức: VDC
\begin{ex}
% ID: L12C4B22-24-TLN
% Mức: VDC
Trong bốn phát biểu sau về tính độ hụt khối, năng lượng liên kết và năng lượng liên kết riêng, có bao nhiêu phát biểu đúng? Chỉ ghi một số nguyên. (1) Kết quả cần có ý nghĩa vật lí. (2) Năng lượng liên kết riêng bằng năng lượng liên kết nhân với số khối. (3) Phải dùng đúng định luật cho tình huống. (4) Năng lượng liên kết $W_{lk}=\Delta m c^2$; năng lượng liên kết riêng bằng $W_{lk}/A$.
\shortans{3}
\loigiai{
Đối chiếu từng phát biểu với kiến thức cốt lõi: Năng lượng liên kết $W_{lk}=\Delta m c^2$; năng lượng liên kết riêng bằng $W_{lk}/A$. Có 3 phát biểu đúng.
}
\end{ex}

% ===== Câu 96 =====
% ID: L12C4B22-24-TN
% Mức: NB
\dangbt{Nhận biết phản ứng hạt nhân và các định luật bảo toàn}
\begin{ex}
% ID: L12C4B22-24-TN
% Mức: NB
Chọn kết luận chính xác về nhận biết phản ứng hạt nhân và các định luật bảo toàn?
\choice
{Trong phản ứng hạt nhân luôn bảo toàn khối lượng nghỉ.}
{Có thể bỏ qua điều kiện áp dụng và đơn vị trong mọi trường hợp.}
{\True Trong phản ứng hạt nhân bảo toàn điện tích, số nucleon, năng lượng toàn phần và động lượng.}
{Kết quả của bài toán không phụ thuộc dữ kiện đã cho.}
\loigiai{
Trong phản ứng hạt nhân bảo toàn điện tích, số nucleon, năng lượng toàn phần và động lượng. Vì vậy phương án $C$ đúng; các phương án còn lại trái với kiến thức hoặc điều kiện áp dụng.
}
\end{ex}

% ===== Câu 97 =====
% ID: L12C4B22-25-DS
% Mức: TH
\dangbt{So sánh độ bền vững của các hạt nhân}
\begin{ex}
% ID: L12C4B22-25-DS
% Mức: TH
Xét so sánh độ bền vững của các hạt nhân (tình huống 1). Hãy xác định tính đúng, sai của bốn phát biểu sau.
\choiceTF
{\True Trong phạm vi so sánh, hạt nhân có năng lượng liên kết riêng lớn hơn thường bền vững hơn.}
{Hạt nhân có số khối lớn hơn luôn bền vững hơn.}
{\True Khi giải cần đọc đúng dữ kiện, dùng hệ đơn vị thống nhất và kiểm tra điều kiện áp dụng.}
{Có thể bỏ qua dữ kiện, đơn vị và ý nghĩa vật lí của kết quả.}
\loigiai{
\begin{itemchoice}
\item (Đúng) Trong phạm vi so sánh, hạt nhân có năng lượng liên kết riêng lớn hơn thường bền vững hơn. (Vì): Trong phạm vi so sánh, hạt nhân có năng lượng liên kết riêng lớn hơn thường bền vững hơn. b) Sai: Hạt nhân có số khối lớn hơn luôn bền vững hơn. c) Đúng: Khi giải cần đọc đúng dữ kiện, dùng hệ đơn vị thống nhất và kiểm tra điều kiện áp dụng. d) Sai: Có thể bỏ qua dữ kiện, đơn vị và ý nghĩa vật lí của kết quả. Kiến thức cốt lõi: Trong phạm vi so sánh, hạt nhân có năng lượng liên kết riêng lớn hơn thường bền vững hơn
\item (Sai) Hạt nhân có số khối lớn hơn luôn bền vững hơn.
\item (Đúng) Khi giải cần đọc đúng dữ kiện, dùng hệ đơn vị thống nhất và kiểm tra điều kiện áp dụng.
\item (Sai) Có thể bỏ qua dữ kiện, đơn vị và ý nghĩa vật lí của kết quả.
\end{itemchoice}
}
\end{ex}

% ===== Câu 98 =====
% ID: L12C4B22-25-TL
% Mức: TH
\begin{bt}
% ID: L12C4B22-25-TL
% Mức: TH
Trình bày kiến thức cốt lõi của so sánh độ bền vững của các hạt nhân và nêu điều kiện áp dụng.
\loigiai{
Cần nêu được: Trong phạm vi so sánh, hạt nhân có năng lượng liên kết riêng lớn hơn thường bền vững hơn. Khi vận dụng phải xác định đúng dữ kiện, điều kiện áp dụng, hệ đơn vị và kiểm tra tính hợp lí của kết quả. Nhận định “Hạt nhân có số khối lớn hơn luôn bền vững hơn.” sai vì trái với kiến thức trên.
}
\end{bt}

% ===== Câu 99 =====
% ID: L12C4B22-25-TLN
% Mức: TH
\begin{ex}
% ID: L12C4B22-25-TLN
% Mức: TH
Trong bốn phát biểu sau về so sánh độ bền vững của các hạt nhân, có bao nhiêu phát biểu đúng? Chỉ ghi một số nguyên. (1) Trong phạm vi so sánh, hạt nhân có năng lượng liên kết riêng lớn hơn thường bền vững hơn. (2) Hạt nhân có số khối lớn hơn luôn bền vững hơn. (3) Cần kiểm tra điều kiện áp dụng trước khi tính toán. (4) Có thể bỏ qua đơn vị của kết quả.
\shortans{2}
\loigiai{
Đối chiếu từng phát biểu với kiến thức cốt lõi: Trong phạm vi so sánh, hạt nhân có năng lượng liên kết riêng lớn hơn thường bền vững hơn. Có 2 phát biểu đúng.
}
\end{ex}

% ===== Câu 100 =====
% ID: L12C4B22-25-TN
% Mức: TH
\dangbt{Nhận biết phản ứng hạt nhân và các định luật bảo toàn}
\begin{ex}
% ID: L12C4B22-25-TN
% Mức: TH
Nội dung nào mô tả đúng nhất về nhận biết phản ứng hạt nhân và các định luật bảo toàn?
\choice
{Trong phản ứng hạt nhân luôn bảo toàn khối lượng nghỉ.}
{Có thể bỏ qua điều kiện áp dụng và đơn vị trong mọi trường hợp.}
{Kết quả của bài toán không phụ thuộc dữ kiện đã cho.}
{\True Trong phản ứng hạt nhân bảo toàn điện tích, số nucleon, năng lượng toàn phần và động lượng.}
\loigiai{
Trong phản ứng hạt nhân bảo toàn điện tích, số nucleon, năng lượng toàn phần và động lượng. Vì vậy phương án $D$ đúng; các phương án còn lại trái với kiến thức hoặc điều kiện áp dụng.
}
\end{ex}

% ===== Câu 101 =====
% ID: L12C4B22-26-DS
% Mức: TH
\dangbt{So sánh độ bền vững của các hạt nhân}
\begin{ex}
% ID: L12C4B22-26-DS
% Mức: TH
Xét so sánh độ bền vững của các hạt nhân (tình huống 2). Hãy xác định tính đúng, sai của bốn phát biểu sau.
\choiceTF
{Hạt nhân có số khối lớn hơn luôn bền vững hơn.}
{\True Trong phạm vi so sánh, hạt nhân có năng lượng liên kết riêng lớn hơn thường bền vững hơn.}
{Có thể bỏ qua dữ kiện, đơn vị và ý nghĩa vật lí của kết quả.}
{\True Khi giải cần đọc đúng dữ kiện, dùng hệ đơn vị thống nhất và kiểm tra điều kiện áp dụng.}
\loigiai{
\begin{itemchoice}
\item (Sai) Hạt nhân có số khối lớn hơn luôn bền vững hơn. (Vì): Hạt nhân có số khối lớn hơn luôn bền vững hơn. b) Đúng: Trong phạm vi so sánh, hạt nhân có năng lượng liên kết riêng lớn hơn thường bền vững hơn. c) Sai: Có thể bỏ qua dữ kiện, đơn vị và ý nghĩa vật lí của kết quả. d) Đúng: Khi giải cần đọc đúng dữ kiện, dùng hệ đơn vị thống nhất và kiểm tra điều kiện áp dụng. Kiến thức cốt lõi: Trong phạm vi so sánh, hạt nhân có năng lượng liên kết riêng lớn hơn thường bền vững hơn
\item (Đúng) Trong phạm vi so sánh, hạt nhân có năng lượng liên kết riêng lớn hơn thường bền vững hơn.
\item (Sai) Có thể bỏ qua dữ kiện, đơn vị và ý nghĩa vật lí của kết quả.
\item (Đúng) Khi giải cần đọc đúng dữ kiện, dùng hệ đơn vị thống nhất và kiểm tra điều kiện áp dụng.
\end{itemchoice}
}
\end{ex}

% ===== Câu 102 =====
% ID: L12C4B22-26-TL
% Mức: TH
\begin{bt}
% ID: L12C4B22-26-TL
% Mức: TH
Giải thích vì sao nhận định sau đúng: “Trong phạm vi so sánh, hạt nhân có năng lượng liên kết riêng lớn hơn thường bền vững hơn.”
\loigiai{
Cần nêu được: Trong phạm vi so sánh, hạt nhân có năng lượng liên kết riêng lớn hơn thường bền vững hơn. Khi vận dụng phải xác định đúng dữ kiện, điều kiện áp dụng, hệ đơn vị và kiểm tra tính hợp lí của kết quả. Nhận định “Hạt nhân có số khối lớn hơn luôn bền vững hơn.” sai vì trái với kiến thức trên.
}
\end{bt}

% ===== Câu 103 =====
% ID: L12C4B22-26-TLN
% Mức: TH
\begin{ex}
% ID: L12C4B22-26-TLN
% Mức: TH
Trong bốn phát biểu sau về so sánh độ bền vững của các hạt nhân, có bao nhiêu phát biểu đúng? Chỉ ghi một số nguyên. (1) Hạt nhân có số khối lớn hơn luôn bền vững hơn. (2) Trong phạm vi so sánh, hạt nhân có năng lượng liên kết riêng lớn hơn thường bền vững hơn. (3) Kết quả phải phù hợp ý nghĩa vật lí. (4) Mọi đại lượng đều có cùng bản chất.
\shortans{2}
\loigiai{
Đối chiếu từng phát biểu với kiến thức cốt lõi: Trong phạm vi so sánh, hạt nhân có năng lượng liên kết riêng lớn hơn thường bền vững hơn. Có 2 phát biểu đúng.
}
\end{ex}

% ===== Câu 104 =====
% ID: L12C4B22-26-TN
% Mức: TH
\dangbt{Nhận biết phản ứng hạt nhân và các định luật bảo toàn}
\begin{ex}
% ID: L12C4B22-26-TN
% Mức: TH
Khi phân tích, phát biểu nào đúng về nhận biết phản ứng hạt nhân và các định luật bảo toàn?
\choice
{\True Trong phản ứng hạt nhân bảo toàn điện tích, số nucleon, năng lượng toàn phần và động lượng.}
{Trong phản ứng hạt nhân luôn bảo toàn khối lượng nghỉ.}
{Có thể bỏ qua điều kiện áp dụng và đơn vị trong mọi trường hợp.}
{Kết quả của bài toán không phụ thuộc dữ kiện đã cho.}
\loigiai{
Trong phản ứng hạt nhân bảo toàn điện tích, số nucleon, năng lượng toàn phần và động lượng. Vì vậy phương án $A$ đúng; các phương án còn lại trái với kiến thức hoặc điều kiện áp dụng.
}
\end{ex}

% ===== Câu 105 =====
% ID: L12C4B22-27-DS
% Mức: VD
\dangbt{So sánh độ bền vững của các hạt nhân}
\begin{ex}
% ID: L12C4B22-27-DS
% Mức: VD
Xét so sánh độ bền vững của các hạt nhân (tình huống 3). Hãy xác định tính đúng, sai của bốn phát biểu sau.
\choiceTF
{\True Trong phạm vi so sánh, hạt nhân có năng lượng liên kết riêng lớn hơn thường bền vững hơn.}
{Hạt nhân có số khối lớn hơn luôn bền vững hơn.}
{\True Khi giải cần đọc đúng dữ kiện, dùng hệ đơn vị thống nhất và kiểm tra điều kiện áp dụng.}
{Có thể bỏ qua dữ kiện, đơn vị và ý nghĩa vật lí của kết quả.}
\loigiai{
\begin{itemchoice}
\item (Đúng) Trong phạm vi so sánh, hạt nhân có năng lượng liên kết riêng lớn hơn thường bền vững hơn. (Vì): Trong phạm vi so sánh, hạt nhân có năng lượng liên kết riêng lớn hơn thường bền vững hơn. b) Sai: Hạt nhân có số khối lớn hơn luôn bền vững hơn. c) Đúng: Khi giải cần đọc đúng dữ kiện, dùng hệ đơn vị thống nhất và kiểm tra điều kiện áp dụng. d) Sai: Có thể bỏ qua dữ kiện, đơn vị và ý nghĩa vật lí của kết quả. Kiến thức cốt lõi: Trong phạm vi so sánh, hạt nhân có năng lượng liên kết riêng lớn hơn thường bền vững hơn
\item (Sai) Hạt nhân có số khối lớn hơn luôn bền vững hơn.
\item (Đúng) Khi giải cần đọc đúng dữ kiện, dùng hệ đơn vị thống nhất và kiểm tra điều kiện áp dụng.
\item (Sai) Có thể bỏ qua dữ kiện, đơn vị và ý nghĩa vật lí của kết quả.
\end{itemchoice}
}
\end{ex}

% ===== Câu 106 =====
% ID: L12C4B22-27-TL
% Mức: VD
\begin{bt}
% ID: L12C4B22-27-TL
% Mức: VD
Phân tích sai lầm trong nhận định: “Hạt nhân có số khối lớn hơn luôn bền vững hơn.”
\loigiai{
Cần nêu được: Trong phạm vi so sánh, hạt nhân có năng lượng liên kết riêng lớn hơn thường bền vững hơn. Khi vận dụng phải xác định đúng dữ kiện, điều kiện áp dụng, hệ đơn vị và kiểm tra tính hợp lí của kết quả. Nhận định “Hạt nhân có số khối lớn hơn luôn bền vững hơn.” sai vì trái với kiến thức trên.
}
\end{bt}

% ===== Câu 107 =====
% ID: L12C4B22-27-TLN
% Mức: TH
\begin{ex}
% ID: L12C4B22-27-TLN
% Mức: TH
Trong bốn phát biểu sau về so sánh độ bền vững của các hạt nhân, có bao nhiêu phát biểu đúng? Chỉ ghi một số nguyên. (1) Cần đổi các đại lượng về hệ đơn vị thống nhất. (2) Trong phạm vi so sánh, hạt nhân có năng lượng liên kết riêng lớn hơn thường bền vững hơn. (3) Hạt nhân có số khối lớn hơn luôn bền vững hơn. (4) Không cần kiểm tra dữ kiện.
\shortans{2}
\loigiai{
Đối chiếu từng phát biểu với kiến thức cốt lõi: Trong phạm vi so sánh, hạt nhân có năng lượng liên kết riêng lớn hơn thường bền vững hơn. Có 2 phát biểu đúng.
}
\end{ex}

% ===== Câu 108 =====
% ID: L12C4B22-27-TN
% Mức: TH
\dangbt{Nhận biết phản ứng hạt nhân và các định luật bảo toàn}
\begin{ex}
% ID: L12C4B22-27-TN
% Mức: TH
Kết luận nào của học sinh là đúng về nhận biết phản ứng hạt nhân và các định luật bảo toàn?
\choice
{Trong phản ứng hạt nhân luôn bảo toàn khối lượng nghỉ.}
{\True Trong phản ứng hạt nhân bảo toàn điện tích, số nucleon, năng lượng toàn phần và động lượng.}
{Có thể bỏ qua điều kiện áp dụng và đơn vị trong mọi trường hợp.}
{Kết quả của bài toán không phụ thuộc dữ kiện đã cho.}
\loigiai{
Trong phản ứng hạt nhân bảo toàn điện tích, số nucleon, năng lượng toàn phần và động lượng. Vì vậy phương án $B$ đúng; các phương án còn lại trái với kiến thức hoặc điều kiện áp dụng.
}
\end{ex}

% ===== Câu 109 =====
% ID: L12C4B22-28-DS
% Mức: VD
\dangbt{So sánh độ bền vững của các hạt nhân}
\begin{ex}
% ID: L12C4B22-28-DS
% Mức: VD
Xét so sánh độ bền vững của các hạt nhân (tình huống 4). Hãy xác định tính đúng, sai của bốn phát biểu sau.
\choiceTF
{Hạt nhân có số khối lớn hơn luôn bền vững hơn.}
{\True Trong phạm vi so sánh, hạt nhân có năng lượng liên kết riêng lớn hơn thường bền vững hơn.}
{Có thể bỏ qua dữ kiện, đơn vị và ý nghĩa vật lí của kết quả.}
{\True Khi giải cần đọc đúng dữ kiện, dùng hệ đơn vị thống nhất và kiểm tra điều kiện áp dụng.}
\loigiai{
\begin{itemchoice}
\item (Sai) Hạt nhân có số khối lớn hơn luôn bền vững hơn. (Vì): Hạt nhân có số khối lớn hơn luôn bền vững hơn. b) Đúng: Trong phạm vi so sánh, hạt nhân có năng lượng liên kết riêng lớn hơn thường bền vững hơn. c) Sai: Có thể bỏ qua dữ kiện, đơn vị và ý nghĩa vật lí của kết quả. d) Đúng: Khi giải cần đọc đúng dữ kiện, dùng hệ đơn vị thống nhất và kiểm tra điều kiện áp dụng. Kiến thức cốt lõi: Trong phạm vi so sánh, hạt nhân có năng lượng liên kết riêng lớn hơn thường bền vững hơn
\item (Đúng) Trong phạm vi so sánh, hạt nhân có năng lượng liên kết riêng lớn hơn thường bền vững hơn.
\item (Sai) Có thể bỏ qua dữ kiện, đơn vị và ý nghĩa vật lí của kết quả.
\item (Đúng) Khi giải cần đọc đúng dữ kiện, dùng hệ đơn vị thống nhất và kiểm tra điều kiện áp dụng.
\end{itemchoice}
}
\end{ex}

% ===== Câu 110 =====
% ID: L12C4B22-28-TL
% Mức: VD
\begin{bt}
% ID: L12C4B22-28-TL
% Mức: VD
Đề xuất các bước giải một bài toán thuộc dạng so sánh độ bền vững của các hạt nhân.
\loigiai{
Cần nêu được: Trong phạm vi so sánh, hạt nhân có năng lượng liên kết riêng lớn hơn thường bền vững hơn. Khi vận dụng phải xác định đúng dữ kiện, điều kiện áp dụng, hệ đơn vị và kiểm tra tính hợp lí của kết quả. Nhận định “Hạt nhân có số khối lớn hơn luôn bền vững hơn.” sai vì trái với kiến thức trên.
}
\end{bt}

% ===== Câu 111 =====
% ID: L12C4B22-28-TLN
% Mức: VD
\begin{ex}
% ID: L12C4B22-28-TLN
% Mức: VD
Trong bốn phát biểu sau về so sánh độ bền vững của các hạt nhân, có bao nhiêu phát biểu đúng? Chỉ ghi một số nguyên. (1) Trong phạm vi so sánh, hạt nhân có năng lượng liên kết riêng lớn hơn thường bền vững hơn. (2) Cần đối chiếu kết quả với điều kiện bài toán. (3) Hạt nhân có số khối lớn hơn luôn bền vững hơn. (4) Mọi trường hợp đều cho cùng một kết quả.
\shortans{2}
\loigiai{
Đối chiếu từng phát biểu với kiến thức cốt lõi: Trong phạm vi so sánh, hạt nhân có năng lượng liên kết riêng lớn hơn thường bền vững hơn. Có 2 phát biểu đúng.
}
\end{ex}

% ===== Câu 112 =====
% ID: L12C4B22-28-TN
% Mức: TH
\dangbt{Nhận biết phản ứng hạt nhân và các định luật bảo toàn}
\begin{ex}
% ID: L12C4B22-28-TN
% Mức: TH
Chọn phát biểu không trái với kiến thức về nhận biết phản ứng hạt nhân và các định luật bảo toàn?
\choice
{Trong phản ứng hạt nhân luôn bảo toàn khối lượng nghỉ.}
{Có thể bỏ qua điều kiện áp dụng và đơn vị trong mọi trường hợp.}
{\True Trong phản ứng hạt nhân bảo toàn điện tích, số nucleon, năng lượng toàn phần và động lượng.}
{Kết quả của bài toán không phụ thuộc dữ kiện đã cho.}
\loigiai{
Trong phản ứng hạt nhân bảo toàn điện tích, số nucleon, năng lượng toàn phần và động lượng. Vì vậy phương án $C$ đúng; các phương án còn lại trái với kiến thức hoặc điều kiện áp dụng.
}
\end{ex}

% ===== Câu 113 =====
% ID: L12C4B22-29-DS
% Mức: VD
\dangbt{So sánh độ bền vững của các hạt nhân}
\begin{ex}
% ID: L12C4B22-29-DS
% Mức: VD
Xét so sánh độ bền vững của các hạt nhân (tình huống 5). Hãy xác định tính đúng, sai của bốn phát biểu sau.
\choiceTF
{\True Trong phạm vi so sánh, hạt nhân có năng lượng liên kết riêng lớn hơn thường bền vững hơn.}
{Có thể bỏ qua dữ kiện, đơn vị và ý nghĩa vật lí của kết quả.}
{Hạt nhân có số khối lớn hơn luôn bền vững hơn.}
{Có thể bỏ qua dữ kiện, đơn vị và ý nghĩa vật lí của kết quả.}
\loigiai{
\begin{itemchoice}
\item (Đúng) Trong phạm vi so sánh, hạt nhân có năng lượng liên kết riêng lớn hơn thường bền vững hơn. (Vì): Trong phạm vi so sánh, hạt nhân có năng lượng liên kết riêng lớn hơn thường bền vững hơn. b) Sai: Có thể bỏ qua dữ kiện, đơn vị và ý nghĩa vật lí của kết quả. c) Sai: Hạt nhân có số khối lớn hơn luôn bền vững hơn. d) Sai: Có thể bỏ qua dữ kiện, đơn vị và ý nghĩa vật lí của kết quả. Kiến thức cốt lõi: Trong phạm vi so sánh, hạt nhân có năng lượng liên kết riêng lớn hơn thường bền vững hơn
\item (Sai) Có thể bỏ qua dữ kiện, đơn vị và ý nghĩa vật lí của kết quả.
\item (Sai) Hạt nhân có số khối lớn hơn luôn bền vững hơn.
\item (Sai) Có thể bỏ qua dữ kiện, đơn vị và ý nghĩa vật lí của kết quả.
\end{itemchoice}
}
\end{ex}

% ===== Câu 114 =====
% ID: L12C4B22-29-TL
% Mức: VD
\begin{bt}
% ID: L12C4B22-29-TL
% Mức: VD
Nêu một tình huống thực tiễn liên quan đến so sánh độ bền vững của các hạt nhân và giải thích bằng kiến thức Vật lí.
\loigiai{
Cần nêu được: Trong phạm vi so sánh, hạt nhân có năng lượng liên kết riêng lớn hơn thường bền vững hơn. Khi vận dụng phải xác định đúng dữ kiện, điều kiện áp dụng, hệ đơn vị và kiểm tra tính hợp lí của kết quả. Nhận định “Hạt nhân có số khối lớn hơn luôn bền vững hơn.” sai vì trái với kiến thức trên.
}
\end{bt}

% ===== Câu 115 =====
% ID: L12C4B22-29-TLN
% Mức: VD
\begin{ex}
% ID: L12C4B22-29-TLN
% Mức: VD
Trong bốn phát biểu sau về so sánh độ bền vững của các hạt nhân, có bao nhiêu phát biểu đúng? Chỉ ghi một số nguyên. (1) Hạt nhân có số khối lớn hơn luôn bền vững hơn. (2) Cần đọc đúng dữ kiện và giả thiết. (3) Trong phạm vi so sánh, hạt nhân có năng lượng liên kết riêng lớn hơn thường bền vững hơn. (4) Có thể thay số trước khi chọn công thức.
\shortans{2}
\loigiai{
Đối chiếu từng phát biểu với kiến thức cốt lõi: Trong phạm vi so sánh, hạt nhân có năng lượng liên kết riêng lớn hơn thường bền vững hơn. Có 2 phát biểu đúng.
}
\end{ex}

% ===== Câu 116 =====
% ID: L12C4B22-29-TN
% Mức: VD
\dangbt{Nhận biết phản ứng hạt nhân và các định luật bảo toàn}
\begin{ex}
% ID: L12C4B22-29-TN
% Mức: VD
Cơ sở Vật lí đúng để giải bài là gì về nhận biết phản ứng hạt nhân và các định luật bảo toàn?
\choice
{Trong phản ứng hạt nhân luôn bảo toàn khối lượng nghỉ.}
{Có thể bỏ qua điều kiện áp dụng và đơn vị trong mọi trường hợp.}
{Kết quả của bài toán không phụ thuộc dữ kiện đã cho.}
{\True Trong phản ứng hạt nhân bảo toàn điện tích, số nucleon, năng lượng toàn phần và động lượng.}
\loigiai{
Trong phản ứng hạt nhân bảo toàn điện tích, số nucleon, năng lượng toàn phần và động lượng. Vì vậy phương án $D$ đúng; các phương án còn lại trái với kiến thức hoặc điều kiện áp dụng.
}
\end{ex}

% ===== Câu 117 =====
% ID: L12C4B22-30-DS
% Mức: TH
\dangbt{Tính năng lượng tỏa ra hoặc thu vào trong phản ứng hạt nhân}
\begin{ex}
% ID: L12C4B22-30-DS
% Mức: TH
Xét tính năng lượng tỏa ra hoặc thu vào trong phản ứng hạt nhân (tình huống 1). Hãy xác định tính đúng, sai của bốn phát biểu sau.
\choiceTF
{\True Năng lượng phản ứng $Q=(m_{trước}-m_{sau})c^2$; $Q>0$ là tỏa năng lượng, $Q<0$ là thu năng lượng.}
{Phản ứng có $m_{sau}\gt m_{trước}$ luôn tỏa năng lượng.}
{\True Khi giải cần đọc đúng dữ kiện, dùng hệ đơn vị thống nhất và kiểm tra điều kiện áp dụng.}
{Có thể bỏ qua dữ kiện, đơn vị và ý nghĩa vật lí của kết quả.}
\loigiai{
\begin{itemchoice}
\item (Đúng) Năng lượng phản ứng $Q=(m_{trước}-m_{sau})c^2$; $Q>0$ là tỏa năng lượng, $Q<0$ là thu năng lượng. (Vì): Năng lượng phản ứng $Q=(m_{trước}-m_{sau})c^2$; $Q>0$ là tỏa năng lượng, $Q<0$ là thu năng lượng. b) Sai: Phản ứng có $m_{sau}\gt m_{trước}$ luôn tỏa năng lượng. c) Đúng: Khi giải cần đọc đúng dữ kiện, dùng hệ đơn vị thống nhất và kiểm tra điều kiện áp dụng. d) Sai: Có thể bỏ qua dữ kiện, đơn vị và ý nghĩa vật lí của kết quả. Kiến thức cốt lõi: Năng lượng phản ứng $Q=(m_{trước}-m_{sau})c^2$; $Q>0$ là tỏa năng lượng, $Q<0$ là thu năng lượng
\item (Sai) Phản ứng có $m_{sau}\gt m_{trước}$ luôn tỏa năng lượng.
\item (Đúng) Khi giải cần đọc đúng dữ kiện, dùng hệ đơn vị thống nhất và kiểm tra điều kiện áp dụng.
\item (Sai) Có thể bỏ qua dữ kiện, đơn vị và ý nghĩa vật lí của kết quả.
\end{itemchoice}
}
\end{ex}

% ===== Câu 118 =====
% ID: L12C4B22-30-TL
% Mức: VDC
\dangbt{So sánh độ bền vững của các hạt nhân}
\begin{bt}
% ID: L12C4B22-30-TL
% Mức: VDC
So sánh nhận định đúng “Trong phạm vi so sánh, hạt nhân có năng lượng liên kết riêng lớn hơn thường bền vững hơn.” với nhận định sai “Hạt nhân có số khối lớn hơn luôn bền vững hơn.”; chỉ rõ căn cứ Vật lí.
\loigiai{
Cần nêu được: Trong phạm vi so sánh, hạt nhân có năng lượng liên kết riêng lớn hơn thường bền vững hơn. Khi vận dụng phải xác định đúng dữ kiện, điều kiện áp dụng, hệ đơn vị và kiểm tra tính hợp lí của kết quả. Nhận định “Hạt nhân có số khối lớn hơn luôn bền vững hơn.” sai vì trái với kiến thức trên.
}
\end{bt}

% ===== Câu 119 =====
% ID: L12C4B22-30-TLN
% Mức: VDC
\begin{ex}
% ID: L12C4B22-30-TLN
% Mức: VDC
Trong bốn phát biểu sau về so sánh độ bền vững của các hạt nhân, có bao nhiêu phát biểu đúng? Chỉ ghi một số nguyên. (1) Kết quả cần có ý nghĩa vật lí. (2) Hạt nhân có số khối lớn hơn luôn bền vững hơn. (3) Phải dùng đúng định luật cho tình huống. (4) Trong phạm vi so sánh, hạt nhân có năng lượng liên kết riêng lớn hơn thường bền vững hơn.
\shortans{3}
\loigiai{
Đối chiếu từng phát biểu với kiến thức cốt lõi: Trong phạm vi so sánh, hạt nhân có năng lượng liên kết riêng lớn hơn thường bền vững hơn. Có 3 phát biểu đúng.
}
\end{ex}

% ===== Câu 120 =====
% ID: L12C4B22-30-TN
% Mức: VD
\dangbt{Nhận biết phản ứng hạt nhân và các định luật bảo toàn}
\begin{ex}
% ID: L12C4B22-30-TN
% Mức: VD
Trong một bài thực hành, nhận xét nào đúng về nhận biết phản ứng hạt nhân và các định luật bảo toàn?
\choice
{\True Trong phản ứng hạt nhân bảo toàn điện tích, số nucleon, năng lượng toàn phần và động lượng.}
{Trong phản ứng hạt nhân luôn bảo toàn khối lượng nghỉ.}
{Có thể bỏ qua điều kiện áp dụng và đơn vị trong mọi trường hợp.}
{Kết quả của bài toán không phụ thuộc dữ kiện đã cho.}
\loigiai{
Trong phản ứng hạt nhân bảo toàn điện tích, số nucleon, năng lượng toàn phần và động lượng. Vì vậy phương án $A$ đúng; các phương án còn lại trái với kiến thức hoặc điều kiện áp dụng.
}
\end{ex}

% ===== Câu 121 =====
% ID: L12C4B22-31-DS
% Mức: TH
\dangbt{Tính năng lượng tỏa ra hoặc thu vào trong phản ứng hạt nhân}
\begin{ex}
% ID: L12C4B22-31-DS
% Mức: TH
Xét tính năng lượng tỏa ra hoặc thu vào trong phản ứng hạt nhân (tình huống 2). Hãy xác định tính đúng, sai của bốn phát biểu sau.
\choiceTF
{Phản ứng có $m_{sau}\gt m_{trước}$ luôn tỏa năng lượng.}
{\True Năng lượng phản ứng $Q=(m_{trước}-m_{sau})c^2$; $Q>0$ là tỏa năng lượng, $Q<0$ là thu năng lượng.}
{Có thể bỏ qua dữ kiện, đơn vị và ý nghĩa vật lí của kết quả.}
{\True Khi giải cần đọc đúng dữ kiện, dùng hệ đơn vị thống nhất và kiểm tra điều kiện áp dụng.}
\loigiai{
\begin{itemchoice}
\item (Sai) Phản ứng có $m_{sau}\gt m_{trước}$ luôn tỏa năng lượng. (Vì): Phản ứng có $m_{sau}\gt m_{trước}$ luôn tỏa năng lượng. b) Đúng: Năng lượng phản ứng $Q=(m_{trước}-m_{sau})c^2$; $Q>0$ là tỏa năng lượng, $Q<0$ là thu năng lượng. c) Sai: Có thể bỏ qua dữ kiện, đơn vị và ý nghĩa vật lí của kết quả. d) Đúng: Khi giải cần đọc đúng dữ kiện, dùng hệ đơn vị thống nhất và kiểm tra điều kiện áp dụng. Kiến thức cốt lõi: Năng lượng phản ứng $Q=(m_{trước}-m_{sau})c^2$; $Q>0$ là tỏa năng lượng, $Q<0$ là thu năng lượng
\item (Đúng) Năng lượng phản ứng $Q=(m_{trước}-m_{sau})c^2$; $Q>0$ là tỏa năng lượng, $Q<0$ là thu năng lượng.
\item (Sai) Có thể bỏ qua dữ kiện, đơn vị và ý nghĩa vật lí của kết quả.
\item (Đúng) Khi giải cần đọc đúng dữ kiện, dùng hệ đơn vị thống nhất và kiểm tra điều kiện áp dụng.
\end{itemchoice}
}
\end{ex}

% ===== Câu 122 =====
% ID: L12C4B22-31-TL
% Mức: TH
\begin{bt}
% ID: L12C4B22-31-TL
% Mức: TH
Trình bày kiến thức cốt lõi của tính năng lượng tỏa ra hoặc thu vào trong phản ứng hạt nhân và nêu điều kiện áp dụng.
\loigiai{
Cần nêu được: Năng lượng phản ứng $Q=(m_{trước}-m_{sau})c^2$; $Q>0$ là tỏa năng lượng, $Q<0$ là thu năng lượng. Khi vận dụng phải xác định đúng dữ kiện, điều kiện áp dụng, hệ đơn vị và kiểm tra tính hợp lí của kết quả. Nhận định “Phản ứng có $m_{sau}\gt m_{trước}$ luôn tỏa năng lượng.” sai vì trái với kiến thức trên.
}
\end{bt}

% ===== Câu 123 =====
% ID: L12C4B22-31-TLN
% Mức: TH
\begin{ex}
% ID: L12C4B22-31-TLN
% Mức: TH
Trong bốn phát biểu sau về tính năng lượng tỏa ra hoặc thu vào trong phản ứng hạt nhân, có bao nhiêu phát biểu đúng? Chỉ ghi một số nguyên. (1) Năng lượng phản ứng $Q=(m_{trước}-m_{sau})c^2$; $Q>0$ là tỏa năng lượng, $Q<0$ là thu năng lượng. (2) Phản ứng có $m_{sau}\gt m_{trước}$ luôn tỏa năng lượng. (3) Cần kiểm tra điều kiện áp dụng trước khi tính toán. (4) Có thể bỏ qua đơn vị của kết quả.
\shortans{2}
\loigiai{
Đối chiếu từng phát biểu với kiến thức cốt lõi: Năng lượng phản ứng $Q=(m_{trước}-m_{sau})c^2$; $Q>0$ là tỏa năng lượng, $Q<0$ là thu năng lượng. Có 2 phát biểu đúng.
}
\end{ex}

% ===== Câu 124 =====
% ID: L12C4B22-31-TN
% Mức: VD
\dangbt{Nhận biết phản ứng hạt nhân và các định luật bảo toàn}
\begin{ex}
% ID: L12C4B22-31-TN
% Mức: VD
Kết luận đúng khi vận dụng là về nhận biết phản ứng hạt nhân và các định luật bảo toàn?
\choice
{Trong phản ứng hạt nhân luôn bảo toàn khối lượng nghỉ.}
{\True Trong phản ứng hạt nhân bảo toàn điện tích, số nucleon, năng lượng toàn phần và động lượng.}
{Có thể bỏ qua điều kiện áp dụng và đơn vị trong mọi trường hợp.}
{Kết quả của bài toán không phụ thuộc dữ kiện đã cho.}
\loigiai{
Trong phản ứng hạt nhân bảo toàn điện tích, số nucleon, năng lượng toàn phần và động lượng. Vì vậy phương án $B$ đúng; các phương án còn lại trái với kiến thức hoặc điều kiện áp dụng.
}
\end{ex}

% ===== Câu 125 =====
% ID: L12C4B22-32-DS
% Mức: VD
\dangbt{Tính năng lượng tỏa ra hoặc thu vào trong phản ứng hạt nhân}
\begin{ex}
% ID: L12C4B22-32-DS
% Mức: VD
Xét tính năng lượng tỏa ra hoặc thu vào trong phản ứng hạt nhân (tình huống 3). Hãy xác định tính đúng, sai của bốn phát biểu sau.
\choiceTF
{\True Năng lượng phản ứng $Q=(m_{trước}-m_{sau})c^2$; $Q>0$ là tỏa năng lượng, $Q<0$ là thu năng lượng.}
{Phản ứng có $m_{sau}\gt m_{trước}$ luôn tỏa năng lượng.}
{\True Khi giải cần đọc đúng dữ kiện, dùng hệ đơn vị thống nhất và kiểm tra điều kiện áp dụng.}
{Có thể bỏ qua dữ kiện, đơn vị và ý nghĩa vật lí của kết quả.}
\loigiai{
\begin{itemchoice}
\item (Đúng) Năng lượng phản ứng $Q=(m_{trước}-m_{sau})c^2$; $Q>0$ là tỏa năng lượng, $Q<0$ là thu năng lượng. (Vì): Năng lượng phản ứng $Q=(m_{trước}-m_{sau})c^2$; $Q>0$ là tỏa năng lượng, $Q<0$ là thu năng lượng. b) Sai: Phản ứng có $m_{sau}\gt m_{trước}$ luôn tỏa năng lượng. c) Đúng: Khi giải cần đọc đúng dữ kiện, dùng hệ đơn vị thống nhất và kiểm tra điều kiện áp dụng. d) Sai: Có thể bỏ qua dữ kiện, đơn vị và ý nghĩa vật lí của kết quả. Kiến thức cốt lõi: Năng lượng phản ứng $Q=(m_{trước}-m_{sau})c^2$; $Q>0$ là tỏa năng lượng, $Q<0$ là thu năng lượng
\item (Sai) Phản ứng có $m_{sau}\gt m_{trước}$ luôn tỏa năng lượng.
\item (Đúng) Khi giải cần đọc đúng dữ kiện, dùng hệ đơn vị thống nhất và kiểm tra điều kiện áp dụng.
\item (Sai) Có thể bỏ qua dữ kiện, đơn vị và ý nghĩa vật lí của kết quả.
\end{itemchoice}
}
\end{ex}

% ===== Câu 126 =====
% ID: L12C4B22-32-TL
% Mức: TH
\begin{bt}
% ID: L12C4B22-32-TL
% Mức: TH
Giải thích vì sao nhận định sau đúng: “Năng lượng phản ứng $Q=(m_{trước}-m_{sau})c^2$; $Q>0$ là tỏa năng lượng, $Q<0$ là thu năng lượng.”
\loigiai{
Cần nêu được: Năng lượng phản ứng $Q=(m_{trước}-m_{sau})c^2$; $Q>0$ là tỏa năng lượng, $Q<0$ là thu năng lượng. Khi vận dụng phải xác định đúng dữ kiện, điều kiện áp dụng, hệ đơn vị và kiểm tra tính hợp lí của kết quả. Nhận định “Phản ứng có $m_{sau}\gt m_{trước}$ luôn tỏa năng lượng.” sai vì trái với kiến thức trên.
}
\end{bt}

% ===== Câu 127 =====
% ID: L12C4B22-32-TLN
% Mức: TH
\begin{ex}
% ID: L12C4B22-32-TLN
% Mức: TH
Trong bốn phát biểu sau về tính năng lượng tỏa ra hoặc thu vào trong phản ứng hạt nhân, có bao nhiêu phát biểu đúng? Chỉ ghi một số nguyên. (1) Phản ứng có $m_{sau}\gt m_{trước}$ luôn tỏa năng lượng. (2) Năng lượng phản ứng $Q=(m_{trước}-m_{sau})c^2$; $Q>0$ là tỏa năng lượng, $Q<0$ là thu năng lượng. (3) Kết quả phải phù hợp ý nghĩa vật lí. (4) Mọi đại lượng đều có cùng bản chất.
\shortans{2}
\loigiai{
Đối chiếu từng phát biểu với kiến thức cốt lõi: Năng lượng phản ứng $Q=(m_{trước}-m_{sau})c^2$; $Q>0$ là tỏa năng lượng, $Q<0$ là thu năng lượng. Có 2 phát biểu đúng.
}
\end{ex}

% ===== Câu 128 =====
% ID: L12C4B22-32-TN
% Mức: NB
\dangbt{Hoàn thành phương trình phản ứng hạt nhân}
\begin{ex}
% ID: L12C4B22-32-TN
% Mức: NB
Phát biểu nào sau đây đúng về hoàn thành phương trình phản ứng hạt nhân?
\choice
{\True Muốn tìm hạt còn thiếu phải bảo toàn tổng số khối $A$ và tổng điện tích $Z$ ở hai vế.}
{Chỉ cần bảo toàn số khối, không cần bảo toàn điện tích.}
{Có thể bỏ qua điều kiện áp dụng và đơn vị trong mọi trường hợp.}
{Kết quả của bài toán không phụ thuộc dữ kiện đã cho.}
\loigiai{
Muốn tìm hạt còn thiếu phải bảo toàn tổng số khối $A$ và tổng điện tích $Z$ ở hai vế. Vì vậy phương án $A$ đúng; các phương án còn lại trái với kiến thức hoặc điều kiện áp dụng.
}
\end{ex}

% ===== Câu 129 =====
% ID: L12C4B22-33-DS
% Mức: VD
\dangbt{Tính năng lượng tỏa ra hoặc thu vào trong phản ứng hạt nhân}
\begin{ex}
% ID: L12C4B22-33-DS
% Mức: VD
Xét tính năng lượng tỏa ra hoặc thu vào trong phản ứng hạt nhân (tình huống 4). Hãy xác định tính đúng, sai của bốn phát biểu sau.
\choiceTF
{Phản ứng có $m_{sau}\gt m_{trước}$ luôn tỏa năng lượng.}
{\True Năng lượng phản ứng $Q=(m_{trước}-m_{sau})c^2$; $Q>0$ là tỏa năng lượng, $Q<0$ là thu năng lượng.}
{Có thể bỏ qua dữ kiện, đơn vị và ý nghĩa vật lí của kết quả.}
{\True Khi giải cần đọc đúng dữ kiện, dùng hệ đơn vị thống nhất và kiểm tra điều kiện áp dụng.}
\loigiai{
\begin{itemchoice}
\item (Sai) Phản ứng có $m_{sau}\gt m_{trước}$ luôn tỏa năng lượng. (Vì): Phản ứng có $m_{sau}\gt m_{trước}$ luôn tỏa năng lượng. b) Đúng: Năng lượng phản ứng $Q=(m_{trước}-m_{sau})c^2$; $Q>0$ là tỏa năng lượng, $Q<0$ là thu năng lượng. c) Sai: Có thể bỏ qua dữ kiện, đơn vị và ý nghĩa vật lí của kết quả. d) Đúng: Khi giải cần đọc đúng dữ kiện, dùng hệ đơn vị thống nhất và kiểm tra điều kiện áp dụng. Kiến thức cốt lõi: Năng lượng phản ứng $Q=(m_{trước}-m_{sau})c^2$; $Q>0$ là tỏa năng lượng, $Q<0$ là thu năng lượng
\item (Đúng) Năng lượng phản ứng $Q=(m_{trước}-m_{sau})c^2$; $Q>0$ là tỏa năng lượng, $Q<0$ là thu năng lượng.
\item (Sai) Có thể bỏ qua dữ kiện, đơn vị và ý nghĩa vật lí của kết quả.
\item (Đúng) Khi giải cần đọc đúng dữ kiện, dùng hệ đơn vị thống nhất và kiểm tra điều kiện áp dụng.
\end{itemchoice}
}
\end{ex}

% ===== Câu 130 =====
% ID: L12C4B22-33-TL
% Mức: VD
\begin{bt}
% ID: L12C4B22-33-TL
% Mức: VD
Phân tích sai lầm trong nhận định: “Phản ứng có $m_{sau}\gt m_{trước}$ luôn tỏa năng lượng.”
\loigiai{
Cần nêu được: Năng lượng phản ứng $Q=(m_{trước}-m_{sau})c^2$; $Q>0$ là tỏa năng lượng, $Q<0$ là thu năng lượng. Khi vận dụng phải xác định đúng dữ kiện, điều kiện áp dụng, hệ đơn vị và kiểm tra tính hợp lí của kết quả. Nhận định “Phản ứng có $m_{sau}\gt m_{trước}$ luôn tỏa năng lượng.” sai vì trái với kiến thức trên.
}
\end{bt}

% ===== Câu 131 =====
% ID: L12C4B22-33-TLN
% Mức: TH
\begin{ex}
% ID: L12C4B22-33-TLN
% Mức: TH
Trong bốn phát biểu sau về tính năng lượng tỏa ra hoặc thu vào trong phản ứng hạt nhân, có bao nhiêu phát biểu đúng? Chỉ ghi một số nguyên. (1) Cần đổi các đại lượng về hệ đơn vị thống nhất. (2) Năng lượng phản ứng $Q=(m_{trước}-m_{sau})c^2$; $Q>0$ là tỏa năng lượng, $Q<0$ là thu năng lượng. (3) Phản ứng có $m_{sau}\gt m_{trước}$ luôn tỏa năng lượng. (4) Không cần kiểm tra dữ kiện.
\shortans{2}
\loigiai{
Đối chiếu từng phát biểu với kiến thức cốt lõi: Năng lượng phản ứng $Q=(m_{trước}-m_{sau})c^2$; $Q>0$ là tỏa năng lượng, $Q<0$ là thu năng lượng. Có 2 phát biểu đúng.
}
\end{ex}

% ===== Câu 132 =====
% ID: L12C4B22-33-TN
% Mức: NB
\dangbt{Hoàn thành phương trình phản ứng hạt nhân}
\begin{ex}
% ID: L12C4B22-33-TN
% Mức: NB
Nhận định nào phù hợp với kiến thức Vật lí về hoàn thành phương trình phản ứng hạt nhân?
\choice
{Chỉ cần bảo toàn số khối, không cần bảo toàn điện tích.}
{\True Muốn tìm hạt còn thiếu phải bảo toàn tổng số khối $A$ và tổng điện tích $Z$ ở hai vế.}
{Có thể bỏ qua điều kiện áp dụng và đơn vị trong mọi trường hợp.}
{Kết quả của bài toán không phụ thuộc dữ kiện đã cho.}
\loigiai{
Muốn tìm hạt còn thiếu phải bảo toàn tổng số khối $A$ và tổng điện tích $Z$ ở hai vế. Vì vậy phương án $B$ đúng; các phương án còn lại trái với kiến thức hoặc điều kiện áp dụng.
}
\end{ex}

% ===== Câu 133 =====
% ID: L12C4B22-34-DS
% Mức: VD
\dangbt{Tính năng lượng tỏa ra hoặc thu vào trong phản ứng hạt nhân}
\begin{ex}
% ID: L12C4B22-34-DS
% Mức: VD
Xét tính năng lượng tỏa ra hoặc thu vào trong phản ứng hạt nhân (tình huống 5). Hãy xác định tính đúng, sai của bốn phát biểu sau.
\choiceTF
{\True Năng lượng phản ứng $Q=(m_{trước}-m_{sau})c^2$; $Q>0$ là tỏa năng lượng, $Q<0$ là thu năng lượng.}
{Có thể bỏ qua dữ kiện, đơn vị và ý nghĩa vật lí của kết quả.}
{Phản ứng có $m_{sau}\gt m_{trước}$ luôn tỏa năng lượng.}
{Có thể bỏ qua dữ kiện, đơn vị và ý nghĩa vật lí của kết quả.}
\loigiai{
\begin{itemchoice}
\item (Đúng) Năng lượng phản ứng $Q=(m_{trước}-m_{sau})c^2$; $Q>0$ là tỏa năng lượng, $Q<0$ là thu năng lượng. (Vì): Năng lượng phản ứng $Q=(m_{trước}-m_{sau})c^2$; $Q>0$ là tỏa năng lượng, $Q<0$ là thu năng lượng. b) Sai: Có thể bỏ qua dữ kiện, đơn vị và ý nghĩa vật lí của kết quả. c) Sai: Phản ứng có $m_{sau}\gt m_{trước}$ luôn tỏa năng lượng. d) Sai: Có thể bỏ qua dữ kiện, đơn vị và ý nghĩa vật lí của kết quả. Kiến thức cốt lõi: Năng lượng phản ứng $Q=(m_{trước}-m_{sau})c^2$; $Q>0$ là tỏa năng lượng, $Q<0$ là thu năng lượng
\item (Sai) Có thể bỏ qua dữ kiện, đơn vị và ý nghĩa vật lí của kết quả.
\item (Sai) Phản ứng có $m_{sau}\gt m_{trước}$ luôn tỏa năng lượng.
\item (Sai) Có thể bỏ qua dữ kiện, đơn vị và ý nghĩa vật lí của kết quả.
\end{itemchoice}
}
\end{ex}

% ===== Câu 134 =====
% ID: L12C4B22-34-TL
% Mức: VD
\begin{bt}
% ID: L12C4B22-34-TL
% Mức: VD
Đề xuất các bước giải một bài toán thuộc dạng tính năng lượng tỏa ra hoặc thu vào trong phản ứng hạt nhân.
\loigiai{
Cần nêu được: Năng lượng phản ứng $Q=(m_{trước}-m_{sau})c^2$; $Q>0$ là tỏa năng lượng, $Q<0$ là thu năng lượng. Khi vận dụng phải xác định đúng dữ kiện, điều kiện áp dụng, hệ đơn vị và kiểm tra tính hợp lí của kết quả. Nhận định “Phản ứng có $m_{sau}\gt m_{trước}$ luôn tỏa năng lượng.” sai vì trái với kiến thức trên.
}
\end{bt}

% ===== Câu 135 =====
% ID: L12C4B22-34-TLN
% Mức: VD
\begin{ex}
% ID: L12C4B22-34-TLN
% Mức: VD
Trong bốn phát biểu sau về tính năng lượng tỏa ra hoặc thu vào trong phản ứng hạt nhân, có bao nhiêu phát biểu đúng? Chỉ ghi một số nguyên. (1) Năng lượng phản ứng $Q=(m_{trước}-m_{sau})c^2$; $Q>0$ là tỏa năng lượng, $Q<0$ là thu năng lượng. (2) Cần đối chiếu kết quả với điều kiện bài toán. (3) Phản ứng có $m_{sau}\gt m_{trước}$ luôn tỏa năng lượng. (4) Mọi trường hợp đều cho cùng một kết quả.
\shortans{2}
\loigiai{
Đối chiếu từng phát biểu với kiến thức cốt lõi: Năng lượng phản ứng $Q=(m_{trước}-m_{sau})c^2$; $Q>0$ là tỏa năng lượng, $Q<0$ là thu năng lượng. Có 2 phát biểu đúng.
}
\end{ex}

% ===== Câu 136 =====
% ID: L12C4B22-34-TN
% Mức: NB
\dangbt{Hoàn thành phương trình phản ứng hạt nhân}
\begin{ex}
% ID: L12C4B22-34-TN
% Mức: NB
Chọn kết luận chính xác về hoàn thành phương trình phản ứng hạt nhân?
\choice
{Chỉ cần bảo toàn số khối, không cần bảo toàn điện tích.}
{Có thể bỏ qua điều kiện áp dụng và đơn vị trong mọi trường hợp.}
{\True Muốn tìm hạt còn thiếu phải bảo toàn tổng số khối $A$ và tổng điện tích $Z$ ở hai vế.}
{Kết quả của bài toán không phụ thuộc dữ kiện đã cho.}
\loigiai{
Muốn tìm hạt còn thiếu phải bảo toàn tổng số khối $A$ và tổng điện tích $Z$ ở hai vế. Vì vậy phương án $C$ đúng; các phương án còn lại trái với kiến thức hoặc điều kiện áp dụng.
}
\end{ex}

% ===== Câu 137 =====
% ID: L12C4B22-35-DS
% Mức: TH
\dangbt{Bài toán tổng hợp về phản ứng hạt nhân, phân hạch và nhiệt hạch}
\begin{ex}
% ID: L12C4B22-35-DS
% Mức: TH
Xét bài toán tổng hợp về phản ứng hạt nhân, phân hạch và nhiệt hạch (tình huống 1). Hãy xác định tính đúng, sai của bốn phát biểu sau.
\choiceTF
{\True Phân hạch hạt nhân nặng và nhiệt hạch hạt nhân nhẹ đều có thể tỏa năng lượng do sản phẩm bền hơn về năng lượng liên kết riêng.}
{Mọi phản ứng hạt nhân đều tự phát và tỏa cùng một năng lượng.}
{\True Khi giải cần đọc đúng dữ kiện, dùng hệ đơn vị thống nhất và kiểm tra điều kiện áp dụng.}
{Có thể bỏ qua dữ kiện, đơn vị và ý nghĩa vật lí của kết quả.}
\loigiai{
\begin{itemchoice}
\item (Đúng) Phân hạch hạt nhân nặng và nhiệt hạch hạt nhân nhẹ đều có thể tỏa năng lượng do sản phẩm bền hơn về năng lượng liên kết riêng. (Vì): Phân hạch hạt nhân nặng và nhiệt hạch hạt nhân nhẹ đều có thể tỏa năng lượng do sản phẩm bền hơn về năng lượng liên kết riêng. b) Sai: Mọi phản ứng hạt nhân đều tự phát và tỏa cùng một năng lượng. c) Đúng: Khi giải cần đọc đúng dữ kiện, dùng hệ đơn vị thống nhất và kiểm tra điều kiện áp dụng. d) Sai: Có thể bỏ qua dữ kiện, đơn vị và ý nghĩa vật lí của kết quả. Kiến thức cốt lõi: Phân hạch hạt nhân nặng và nhiệt hạch hạt nhân nhẹ đều có thể tỏa năng lượng do sản phẩm bền hơn về năng lượng liên kết riêng
\item (Sai) Mọi phản ứng hạt nhân đều tự phát và tỏa cùng một năng lượng.
\item (Đúng) Khi giải cần đọc đúng dữ kiện, dùng hệ đơn vị thống nhất và kiểm tra điều kiện áp dụng.
\item (Sai) Có thể bỏ qua dữ kiện, đơn vị và ý nghĩa vật lí của kết quả.
\end{itemchoice}
}
\end{ex}

% ===== Câu 138 =====
% ID: L12C4B22-35-TL
% Mức: VD
\dangbt{Tính năng lượng tỏa ra hoặc thu vào trong phản ứng hạt nhân}
\begin{bt}
% ID: L12C4B22-35-TL
% Mức: VD
Nêu một tình huống thực tiễn liên quan đến tính năng lượng tỏa ra hoặc thu vào trong phản ứng hạt nhân và giải thích bằng kiến thức Vật lí.
\loigiai{
Cần nêu được: Năng lượng phản ứng $Q=(m_{trước}-m_{sau})c^2$; $Q>0$ là tỏa năng lượng, $Q<0$ là thu năng lượng. Khi vận dụng phải xác định đúng dữ kiện, điều kiện áp dụng, hệ đơn vị và kiểm tra tính hợp lí của kết quả. Nhận định “Phản ứng có $m_{sau}\gt m_{trước}$ luôn tỏa năng lượng.” sai vì trái với kiến thức trên.
}
\end{bt}

% ===== Câu 139 =====
% ID: L12C4B22-35-TLN
% Mức: VD
\begin{ex}
% ID: L12C4B22-35-TLN
% Mức: VD
Trong bốn phát biểu sau về tính năng lượng tỏa ra hoặc thu vào trong phản ứng hạt nhân, có bao nhiêu phát biểu đúng? Chỉ ghi một số nguyên. (1) Phản ứng có $m_{sau}\gt m_{trước}$ luôn tỏa năng lượng. (2) Cần đọc đúng dữ kiện và giả thiết. (3) Năng lượng phản ứng $Q=(m_{trước}-m_{sau})c^2$; $Q>0$ là tỏa năng lượng, $Q<0$ là thu năng lượng. (4) Có thể thay số trước khi chọn công thức.
\shortans{2}
\loigiai{
Đối chiếu từng phát biểu với kiến thức cốt lõi: Năng lượng phản ứng $Q=(m_{trước}-m_{sau})c^2$; $Q>0$ là tỏa năng lượng, $Q<0$ là thu năng lượng. Có 2 phát biểu đúng.
}
\end{ex}

% ===== Câu 140 =====
% ID: L12C4B22-35-TN
% Mức: TH
\dangbt{Hoàn thành phương trình phản ứng hạt nhân}
\begin{ex}
% ID: L12C4B22-35-TN
% Mức: TH
Nội dung nào mô tả đúng nhất về hoàn thành phương trình phản ứng hạt nhân?
\choice
{Chỉ cần bảo toàn số khối, không cần bảo toàn điện tích.}
{Có thể bỏ qua điều kiện áp dụng và đơn vị trong mọi trường hợp.}
{Kết quả của bài toán không phụ thuộc dữ kiện đã cho.}
{\True Muốn tìm hạt còn thiếu phải bảo toàn tổng số khối $A$ và tổng điện tích $Z$ ở hai vế.}
\loigiai{
Muốn tìm hạt còn thiếu phải bảo toàn tổng số khối $A$ và tổng điện tích $Z$ ở hai vế. Vì vậy phương án $D$ đúng; các phương án còn lại trái với kiến thức hoặc điều kiện áp dụng.
}
\end{ex}

% ===== Câu 141 =====
% ID: L12C4B22-36-DS
% Mức: TH
\dangbt{Bài toán tổng hợp về phản ứng hạt nhân, phân hạch và nhiệt hạch}
\begin{ex}
% ID: L12C4B22-36-DS
% Mức: TH
Xét bài toán tổng hợp về phản ứng hạt nhân, phân hạch và nhiệt hạch (tình huống 2). Hãy xác định tính đúng, sai của bốn phát biểu sau.
\choiceTF
{Mọi phản ứng hạt nhân đều tự phát và tỏa cùng một năng lượng.}
{\True Phân hạch hạt nhân nặng và nhiệt hạch hạt nhân nhẹ đều có thể tỏa năng lượng do sản phẩm bền hơn về năng lượng liên kết riêng.}
{Có thể bỏ qua dữ kiện, đơn vị và ý nghĩa vật lí của kết quả.}
{\True Khi giải cần đọc đúng dữ kiện, dùng hệ đơn vị thống nhất và kiểm tra điều kiện áp dụng.}
\loigiai{
\begin{itemchoice}
\item (Sai) Mọi phản ứng hạt nhân đều tự phát và tỏa cùng một năng lượng. (Vì): Mọi phản ứng hạt nhân đều tự phát và tỏa cùng một năng lượng. b) Đúng: Phân hạch hạt nhân nặng và nhiệt hạch hạt nhân nhẹ đều có thể tỏa năng lượng do sản phẩm bền hơn về năng lượng liên kết riêng. c) Sai: Có thể bỏ qua dữ kiện, đơn vị và ý nghĩa vật lí của kết quả. d) Đúng: Khi giải cần đọc đúng dữ kiện, dùng hệ đơn vị thống nhất và kiểm tra điều kiện áp dụng. Kiến thức cốt lõi: Phân hạch hạt nhân nặng và nhiệt hạch hạt nhân nhẹ đều có thể tỏa năng lượng do sản phẩm bền hơn về năng lượng liên kết riêng
\item (Đúng) Phân hạch hạt nhân nặng và nhiệt hạch hạt nhân nhẹ đều có thể tỏa năng lượng do sản phẩm bền hơn về năng lượng liên kết riêng.
\item (Sai) Có thể bỏ qua dữ kiện, đơn vị và ý nghĩa vật lí của kết quả.
\item (Đúng) Khi giải cần đọc đúng dữ kiện, dùng hệ đơn vị thống nhất và kiểm tra điều kiện áp dụng.
\end{itemchoice}
}
\end{ex}

% ===== Câu 142 =====
% ID: L12C4B22-36-TL
% Mức: VDC
\dangbt{Tính năng lượng tỏa ra hoặc thu vào trong phản ứng hạt nhân}
\begin{bt}
% ID: L12C4B22-36-TL
% Mức: VDC
So sánh nhận định đúng “Năng lượng phản ứng $Q=(m_{trước}-m_{sau})c^2$; $Q>0$ là tỏa năng lượng, $Q<0$ là thu năng lượng.” với nhận định sai “Phản ứng có $m_{sau}\gt m_{trước}$ luôn tỏa năng lượng.”; chỉ rõ căn cứ Vật lí.
\loigiai{
Cần nêu được: Năng lượng phản ứng $Q=(m_{trước}-m_{sau})c^2$; $Q>0$ là tỏa năng lượng, $Q<0$ là thu năng lượng. Khi vận dụng phải xác định đúng dữ kiện, điều kiện áp dụng, hệ đơn vị và kiểm tra tính hợp lí của kết quả. Nhận định “Phản ứng có $m_{sau}\gt m_{trước}$ luôn tỏa năng lượng.” sai vì trái với kiến thức trên.
}
\end{bt}

% ===== Câu 143 =====
% ID: L12C4B22-36-TLN
% Mức: VDC
\begin{ex}
% ID: L12C4B22-36-TLN
% Mức: VDC
Trong bốn phát biểu sau về tính năng lượng tỏa ra hoặc thu vào trong phản ứng hạt nhân, có bao nhiêu phát biểu đúng? Chỉ ghi một số nguyên. (1) Kết quả cần có ý nghĩa vật lí. (2) Phản ứng có $m_{sau}\gt m_{trước}$ luôn tỏa năng lượng. (3) Phải dùng đúng định luật cho tình huống. (4) Năng lượng phản ứng $Q=(m_{trước}-m_{sau})c^2$; $Q>0$ là tỏa năng lượng, $Q<0$ là thu năng lượng.
\shortans{3}
\loigiai{
Đối chiếu từng phát biểu với kiến thức cốt lõi: Năng lượng phản ứng $Q=(m_{trước}-m_{sau})c^2$; $Q>0$ là tỏa năng lượng, $Q<0$ là thu năng lượng. Có 3 phát biểu đúng.
}
\end{ex}

% ===== Câu 144 =====
% ID: L12C4B22-36-TN
% Mức: TH
\dangbt{Hoàn thành phương trình phản ứng hạt nhân}
\begin{ex}
% ID: L12C4B22-36-TN
% Mức: TH
Khi phân tích, phát biểu nào đúng về hoàn thành phương trình phản ứng hạt nhân?
\choice
{\True Muốn tìm hạt còn thiếu phải bảo toàn tổng số khối $A$ và tổng điện tích $Z$ ở hai vế.}
{Chỉ cần bảo toàn số khối, không cần bảo toàn điện tích.}
{Có thể bỏ qua điều kiện áp dụng và đơn vị trong mọi trường hợp.}
{Kết quả của bài toán không phụ thuộc dữ kiện đã cho.}
\loigiai{
Muốn tìm hạt còn thiếu phải bảo toàn tổng số khối $A$ và tổng điện tích $Z$ ở hai vế. Vì vậy phương án $A$ đúng; các phương án còn lại trái với kiến thức hoặc điều kiện áp dụng.
}
\end{ex}

% ===== Câu 145 =====
% ID: L12C4B22-37-DS
% Mức: VD
\dangbt{Bài toán tổng hợp về phản ứng hạt nhân, phân hạch và nhiệt hạch}
\begin{ex}
% ID: L12C4B22-37-DS
% Mức: VD
Xét bài toán tổng hợp về phản ứng hạt nhân, phân hạch và nhiệt hạch (tình huống 3). Hãy xác định tính đúng, sai của bốn phát biểu sau.
\choiceTF
{\True Phân hạch hạt nhân nặng và nhiệt hạch hạt nhân nhẹ đều có thể tỏa năng lượng do sản phẩm bền hơn về năng lượng liên kết riêng.}
{Mọi phản ứng hạt nhân đều tự phát và tỏa cùng một năng lượng.}
{\True Khi giải cần đọc đúng dữ kiện, dùng hệ đơn vị thống nhất và kiểm tra điều kiện áp dụng.}
{Có thể bỏ qua dữ kiện, đơn vị và ý nghĩa vật lí của kết quả.}
\loigiai{
\begin{itemchoice}
\item (Đúng) Phân hạch hạt nhân nặng và nhiệt hạch hạt nhân nhẹ đều có thể tỏa năng lượng do sản phẩm bền hơn về năng lượng liên kết riêng. (Vì): Phân hạch hạt nhân nặng và nhiệt hạch hạt nhân nhẹ đều có thể tỏa năng lượng do sản phẩm bền hơn về năng lượng liên kết riêng. b) Sai: Mọi phản ứng hạt nhân đều tự phát và tỏa cùng một năng lượng. c) Đúng: Khi giải cần đọc đúng dữ kiện, dùng hệ đơn vị thống nhất và kiểm tra điều kiện áp dụng. d) Sai: Có thể bỏ qua dữ kiện, đơn vị và ý nghĩa vật lí của kết quả. Kiến thức cốt lõi: Phân hạch hạt nhân nặng và nhiệt hạch hạt nhân nhẹ đều có thể tỏa năng lượng do sản phẩm bền hơn về năng lượng liên kết riêng
\item (Sai) Mọi phản ứng hạt nhân đều tự phát và tỏa cùng một năng lượng.
\item (Đúng) Khi giải cần đọc đúng dữ kiện, dùng hệ đơn vị thống nhất và kiểm tra điều kiện áp dụng.
\item (Sai) Có thể bỏ qua dữ kiện, đơn vị và ý nghĩa vật lí của kết quả.
\end{itemchoice}
}
\end{ex}

% ===== Câu 146 =====
% ID: L12C4B22-37-TL
% Mức: TH
\begin{bt}
% ID: L12C4B22-37-TL
% Mức: TH
Trình bày kiến thức cốt lõi của bài toán tổng hợp về phản ứng hạt nhân, phân hạch và nhiệt hạch và nêu điều kiện áp dụng.
\loigiai{
Cần nêu được: Phân hạch hạt nhân nặng và nhiệt hạch hạt nhân nhẹ đều có thể tỏa năng lượng do sản phẩm bền hơn về năng lượng liên kết riêng. Khi vận dụng phải xác định đúng dữ kiện, điều kiện áp dụng, hệ đơn vị và kiểm tra tính hợp lí của kết quả. Nhận định “Mọi phản ứng hạt nhân đều tự phát và tỏa cùng một năng lượng.” sai vì trái với kiến thức trên.
}
\end{bt}

% ===== Câu 147 =====
% ID: L12C4B22-37-TLN
% Mức: TH
\begin{ex}
% ID: L12C4B22-37-TLN
% Mức: TH
Trong bốn phát biểu sau về bài toán tổng hợp về phản ứng hạt nhân, phân hạch và nhiệt hạch, có bao nhiêu phát biểu đúng? Chỉ ghi một số nguyên. (1) Phân hạch hạt nhân nặng và nhiệt hạch hạt nhân nhẹ đều có thể tỏa năng lượng do sản phẩm bền hơn về năng lượng liên kết riêng. (2) Mọi phản ứng hạt nhân đều tự phát và tỏa cùng một năng lượng. (3) Cần kiểm tra điều kiện áp dụng trước khi tính toán. (4) Có thể bỏ qua đơn vị của kết quả.
\shortans{2}
\loigiai{
Đối chiếu từng phát biểu với kiến thức cốt lõi: Phân hạch hạt nhân nặng và nhiệt hạch hạt nhân nhẹ đều có thể tỏa năng lượng do sản phẩm bền hơn về năng lượng liên kết riêng. Có 2 phát biểu đúng.
}
\end{ex}

% ===== Câu 148 =====
% ID: L12C4B22-37-TN
% Mức: TH
\dangbt{Hoàn thành phương trình phản ứng hạt nhân}
\begin{ex}
% ID: L12C4B22-37-TN
% Mức: TH
Kết luận nào của học sinh là đúng về hoàn thành phương trình phản ứng hạt nhân?
\choice
{Chỉ cần bảo toàn số khối, không cần bảo toàn điện tích.}
{\True Muốn tìm hạt còn thiếu phải bảo toàn tổng số khối $A$ và tổng điện tích $Z$ ở hai vế.}
{Có thể bỏ qua điều kiện áp dụng và đơn vị trong mọi trường hợp.}
{Kết quả của bài toán không phụ thuộc dữ kiện đã cho.}
\loigiai{
Muốn tìm hạt còn thiếu phải bảo toàn tổng số khối $A$ và tổng điện tích $Z$ ở hai vế. Vì vậy phương án $B$ đúng; các phương án còn lại trái với kiến thức hoặc điều kiện áp dụng.
}
\end{ex}

% ===== Câu 149 =====
% ID: L12C4B22-38-DS
% Mức: VD
\dangbt{Bài toán tổng hợp về phản ứng hạt nhân, phân hạch và nhiệt hạch}
\begin{ex}
% ID: L12C4B22-38-DS
% Mức: VD
Xét bài toán tổng hợp về phản ứng hạt nhân, phân hạch và nhiệt hạch (tình huống 4). Hãy xác định tính đúng, sai của bốn phát biểu sau.
\choiceTF
{Mọi phản ứng hạt nhân đều tự phát và tỏa cùng một năng lượng.}
{\True Phân hạch hạt nhân nặng và nhiệt hạch hạt nhân nhẹ đều có thể tỏa năng lượng do sản phẩm bền hơn về năng lượng liên kết riêng.}
{Có thể bỏ qua dữ kiện, đơn vị và ý nghĩa vật lí của kết quả.}
{\True Khi giải cần đọc đúng dữ kiện, dùng hệ đơn vị thống nhất và kiểm tra điều kiện áp dụng.}
\loigiai{
\begin{itemchoice}
\item (Sai) Mọi phản ứng hạt nhân đều tự phát và tỏa cùng một năng lượng. (Vì): Mọi phản ứng hạt nhân đều tự phát và tỏa cùng một năng lượng. b) Đúng: Phân hạch hạt nhân nặng và nhiệt hạch hạt nhân nhẹ đều có thể tỏa năng lượng do sản phẩm bền hơn về năng lượng liên kết riêng. c) Sai: Có thể bỏ qua dữ kiện, đơn vị và ý nghĩa vật lí của kết quả. d) Đúng: Khi giải cần đọc đúng dữ kiện, dùng hệ đơn vị thống nhất và kiểm tra điều kiện áp dụng. Kiến thức cốt lõi: Phân hạch hạt nhân nặng và nhiệt hạch hạt nhân nhẹ đều có thể tỏa năng lượng do sản phẩm bền hơn về năng lượng liên kết riêng
\item (Đúng) Phân hạch hạt nhân nặng và nhiệt hạch hạt nhân nhẹ đều có thể tỏa năng lượng do sản phẩm bền hơn về năng lượng liên kết riêng.
\item (Sai) Có thể bỏ qua dữ kiện, đơn vị và ý nghĩa vật lí của kết quả.
\item (Đúng) Khi giải cần đọc đúng dữ kiện, dùng hệ đơn vị thống nhất và kiểm tra điều kiện áp dụng.
\end{itemchoice}
}
\end{ex}

% ===== Câu 150 =====
% ID: L12C4B22-38-TL
% Mức: TH
\begin{bt}
% ID: L12C4B22-38-TL
% Mức: TH
Giải thích vì sao nhận định sau đúng: “Phân hạch hạt nhân nặng và nhiệt hạch hạt nhân nhẹ đều có thể tỏa năng lượng do sản phẩm bền hơn về năng lượng liên kết riêng.”
\loigiai{
Cần nêu được: Phân hạch hạt nhân nặng và nhiệt hạch hạt nhân nhẹ đều có thể tỏa năng lượng do sản phẩm bền hơn về năng lượng liên kết riêng. Khi vận dụng phải xác định đúng dữ kiện, điều kiện áp dụng, hệ đơn vị và kiểm tra tính hợp lí của kết quả. Nhận định “Mọi phản ứng hạt nhân đều tự phát và tỏa cùng một năng lượng.” sai vì trái với kiến thức trên.
}
\end{bt}

% ===== Câu 151 =====
% ID: L12C4B22-38-TLN
% Mức: TH
\begin{ex}
% ID: L12C4B22-38-TLN
% Mức: TH
Trong bốn phát biểu sau về bài toán tổng hợp về phản ứng hạt nhân, phân hạch và nhiệt hạch, có bao nhiêu phát biểu đúng? Chỉ ghi một số nguyên. (1) Mọi phản ứng hạt nhân đều tự phát và tỏa cùng một năng lượng. (2) Phân hạch hạt nhân nặng và nhiệt hạch hạt nhân nhẹ đều có thể tỏa năng lượng do sản phẩm bền hơn về năng lượng liên kết riêng. (3) Kết quả phải phù hợp ý nghĩa vật lí. (4) Mọi đại lượng đều có cùng bản chất.
\shortans{2}
\loigiai{
Đối chiếu từng phát biểu với kiến thức cốt lõi: Phân hạch hạt nhân nặng và nhiệt hạch hạt nhân nhẹ đều có thể tỏa năng lượng do sản phẩm bền hơn về năng lượng liên kết riêng. Có 2 phát biểu đúng.
}
\end{ex}

% ===== Câu 152 =====
% ID: L12C4B22-38-TN
% Mức: TH
\dangbt{Hoàn thành phương trình phản ứng hạt nhân}
\begin{ex}
% ID: L12C4B22-38-TN
% Mức: TH
Chọn phát biểu không trái với kiến thức về hoàn thành phương trình phản ứng hạt nhân?
\choice
{Chỉ cần bảo toàn số khối, không cần bảo toàn điện tích.}
{Có thể bỏ qua điều kiện áp dụng và đơn vị trong mọi trường hợp.}
{\True Muốn tìm hạt còn thiếu phải bảo toàn tổng số khối $A$ và tổng điện tích $Z$ ở hai vế.}
{Kết quả của bài toán không phụ thuộc dữ kiện đã cho.}
\loigiai{
Muốn tìm hạt còn thiếu phải bảo toàn tổng số khối $A$ và tổng điện tích $Z$ ở hai vế. Vì vậy phương án $C$ đúng; các phương án còn lại trái với kiến thức hoặc điều kiện áp dụng.
}
\end{ex}

% ===== Câu 153 =====
% ID: L12C4B22-39-DS
% Mức: VD
\dangbt{Bài toán tổng hợp về phản ứng hạt nhân, phân hạch và nhiệt hạch}
\begin{ex}
% ID: L12C4B22-39-DS
% Mức: VD
Xét bài toán tổng hợp về phản ứng hạt nhân, phân hạch và nhiệt hạch (tình huống 5). Hãy xác định tính đúng, sai của bốn phát biểu sau.
\choiceTF
{\True Phân hạch hạt nhân nặng và nhiệt hạch hạt nhân nhẹ đều có thể tỏa năng lượng do sản phẩm bền hơn về năng lượng liên kết riêng.}
{Có thể bỏ qua dữ kiện, đơn vị và ý nghĩa vật lí của kết quả.}
{Mọi phản ứng hạt nhân đều tự phát và tỏa cùng một năng lượng.}
{Có thể bỏ qua dữ kiện, đơn vị và ý nghĩa vật lí của kết quả.}
\loigiai{
\begin{itemchoice}
\item (Đúng) Phân hạch hạt nhân nặng và nhiệt hạch hạt nhân nhẹ đều có thể tỏa năng lượng do sản phẩm bền hơn về năng lượng liên kết riêng. (Vì): Phân hạch hạt nhân nặng và nhiệt hạch hạt nhân nhẹ đều có thể tỏa năng lượng do sản phẩm bền hơn về năng lượng liên kết riêng. b) Sai: Có thể bỏ qua dữ kiện, đơn vị và ý nghĩa vật lí của kết quả. c) Sai: Mọi phản ứng hạt nhân đều tự phát và tỏa cùng một năng lượng. d) Sai: Có thể bỏ qua dữ kiện, đơn vị và ý nghĩa vật lí của kết quả. Kiến thức cốt lõi: Phân hạch hạt nhân nặng và nhiệt hạch hạt nhân nhẹ đều có thể tỏa năng lượng do sản phẩm bền hơn về năng lượng liên kết riêng
\item (Sai) Có thể bỏ qua dữ kiện, đơn vị và ý nghĩa vật lí của kết quả.
\item (Sai) Mọi phản ứng hạt nhân đều tự phát và tỏa cùng một năng lượng.
\item (Sai) Có thể bỏ qua dữ kiện, đơn vị và ý nghĩa vật lí của kết quả.
\end{itemchoice}
}
\end{ex}

% ===== Câu 154 =====
% ID: L12C4B22-39-TL
% Mức: VD
\begin{bt}
% ID: L12C4B22-39-TL
% Mức: VD
Phân tích sai lầm trong nhận định: “Mọi phản ứng hạt nhân đều tự phát và tỏa cùng một năng lượng.”
\loigiai{
Cần nêu được: Phân hạch hạt nhân nặng và nhiệt hạch hạt nhân nhẹ đều có thể tỏa năng lượng do sản phẩm bền hơn về năng lượng liên kết riêng. Khi vận dụng phải xác định đúng dữ kiện, điều kiện áp dụng, hệ đơn vị và kiểm tra tính hợp lí của kết quả. Nhận định “Mọi phản ứng hạt nhân đều tự phát và tỏa cùng một năng lượng.” sai vì trái với kiến thức trên.
}
\end{bt}

% ===== Câu 155 =====
% ID: L12C4B22-39-TLN
% Mức: TH
\begin{ex}
% ID: L12C4B22-39-TLN
% Mức: TH
Trong bốn phát biểu sau về bài toán tổng hợp về phản ứng hạt nhân, phân hạch và nhiệt hạch, có bao nhiêu phát biểu đúng? Chỉ ghi một số nguyên. (1) Cần đổi các đại lượng về hệ đơn vị thống nhất. (2) Phân hạch hạt nhân nặng và nhiệt hạch hạt nhân nhẹ đều có thể tỏa năng lượng do sản phẩm bền hơn về năng lượng liên kết riêng. (3) Mọi phản ứng hạt nhân đều tự phát và tỏa cùng một năng lượng. (4) Không cần kiểm tra dữ kiện.
\shortans{2}
\loigiai{
Đối chiếu từng phát biểu với kiến thức cốt lõi: Phân hạch hạt nhân nặng và nhiệt hạch hạt nhân nhẹ đều có thể tỏa năng lượng do sản phẩm bền hơn về năng lượng liên kết riêng. Có 2 phát biểu đúng.
}
\end{ex}

% ===== Câu 156 =====
% ID: L12C4B22-39-TN
% Mức: VD
\dangbt{Hoàn thành phương trình phản ứng hạt nhân}
\begin{ex}
% ID: L12C4B22-39-TN
% Mức: VD
Cơ sở Vật lí đúng để giải bài là gì về hoàn thành phương trình phản ứng hạt nhân?
\choice
{Chỉ cần bảo toàn số khối, không cần bảo toàn điện tích.}
{Có thể bỏ qua điều kiện áp dụng và đơn vị trong mọi trường hợp.}
{Kết quả của bài toán không phụ thuộc dữ kiện đã cho.}
{\True Muốn tìm hạt còn thiếu phải bảo toàn tổng số khối $A$ và tổng điện tích $Z$ ở hai vế.}
\loigiai{
Muốn tìm hạt còn thiếu phải bảo toàn tổng số khối $A$ và tổng điện tích $Z$ ở hai vế. Vì vậy phương án $D$ đúng; các phương án còn lại trái với kiến thức hoặc điều kiện áp dụng.
}
\end{ex}

% ===== Câu 157 =====
% ID: L12C4B22-40-TL
% Mức: VD
\dangbt{Bài toán tổng hợp về phản ứng hạt nhân, phân hạch và nhiệt hạch}
\begin{bt}
% ID: L12C4B22-40-TL
% Mức: VD
Đề xuất các bước giải một bài toán thuộc dạng bài toán tổng hợp về phản ứng hạt nhân, phân hạch và nhiệt hạch.
\loigiai{
Cần nêu được: Phân hạch hạt nhân nặng và nhiệt hạch hạt nhân nhẹ đều có thể tỏa năng lượng do sản phẩm bền hơn về năng lượng liên kết riêng. Khi vận dụng phải xác định đúng dữ kiện, điều kiện áp dụng, hệ đơn vị và kiểm tra tính hợp lí của kết quả. Nhận định “Mọi phản ứng hạt nhân đều tự phát và tỏa cùng một năng lượng.” sai vì trái với kiến thức trên.
}
\end{bt}

% ===== Câu 158 =====
% ID: L12C4B22-40-TLN
% Mức: VD
\begin{ex}
% ID: L12C4B22-40-TLN
% Mức: VD
Trong bốn phát biểu sau về bài toán tổng hợp về phản ứng hạt nhân, phân hạch và nhiệt hạch, có bao nhiêu phát biểu đúng? Chỉ ghi một số nguyên. (1) Phân hạch hạt nhân nặng và nhiệt hạch hạt nhân nhẹ đều có thể tỏa năng lượng do sản phẩm bền hơn về năng lượng liên kết riêng. (2) Cần đối chiếu kết quả với điều kiện bài toán. (3) Mọi phản ứng hạt nhân đều tự phát và tỏa cùng một năng lượng. (4) Mọi trường hợp đều cho cùng một kết quả.
\shortans{2}
\loigiai{
Đối chiếu từng phát biểu với kiến thức cốt lõi: Phân hạch hạt nhân nặng và nhiệt hạch hạt nhân nhẹ đều có thể tỏa năng lượng do sản phẩm bền hơn về năng lượng liên kết riêng. Có 2 phát biểu đúng.
}
\end{ex}

% ===== Câu 159 =====
% ID: L12C4B22-40-TN
% Mức: VD
\dangbt{Hoàn thành phương trình phản ứng hạt nhân}
\begin{ex}
% ID: L12C4B22-40-TN
% Mức: VD
Trong một bài thực hành, nhận xét nào đúng về hoàn thành phương trình phản ứng hạt nhân?
\choice
{\True Muốn tìm hạt còn thiếu phải bảo toàn tổng số khối $A$ và tổng điện tích $Z$ ở hai vế.}
{Chỉ cần bảo toàn số khối, không cần bảo toàn điện tích.}
{Có thể bỏ qua điều kiện áp dụng và đơn vị trong mọi trường hợp.}
{Kết quả của bài toán không phụ thuộc dữ kiện đã cho.}
\loigiai{
Muốn tìm hạt còn thiếu phải bảo toàn tổng số khối $A$ và tổng điện tích $Z$ ở hai vế. Vì vậy phương án $A$ đúng; các phương án còn lại trái với kiến thức hoặc điều kiện áp dụng.
}
\end{ex}

% ===== Câu 160 =====
% ID: L12C4B22-41-TL
% Mức: VD
\dangbt{Bài toán tổng hợp về phản ứng hạt nhân, phân hạch và nhiệt hạch}
\begin{bt}
% ID: L12C4B22-41-TL
% Mức: VD
Nêu một tình huống thực tiễn liên quan đến bài toán tổng hợp về phản ứng hạt nhân, phân hạch và nhiệt hạch và giải thích bằng kiến thức Vật lí.
\loigiai{
Cần nêu được: Phân hạch hạt nhân nặng và nhiệt hạch hạt nhân nhẹ đều có thể tỏa năng lượng do sản phẩm bền hơn về năng lượng liên kết riêng. Khi vận dụng phải xác định đúng dữ kiện, điều kiện áp dụng, hệ đơn vị và kiểm tra tính hợp lí của kết quả. Nhận định “Mọi phản ứng hạt nhân đều tự phát và tỏa cùng một năng lượng.” sai vì trái với kiến thức trên.
}
\end{bt}

% ===== Câu 161 =====
% ID: L12C4B22-41-TLN
% Mức: VD
\begin{ex}
% ID: L12C4B22-41-TLN
% Mức: VD
Trong bốn phát biểu sau về bài toán tổng hợp về phản ứng hạt nhân, phân hạch và nhiệt hạch, có bao nhiêu phát biểu đúng? Chỉ ghi một số nguyên. (1) Mọi phản ứng hạt nhân đều tự phát và tỏa cùng một năng lượng. (2) Cần đọc đúng dữ kiện và giả thiết. (3) Phân hạch hạt nhân nặng và nhiệt hạch hạt nhân nhẹ đều có thể tỏa năng lượng do sản phẩm bền hơn về năng lượng liên kết riêng. (4) Có thể thay số trước khi chọn công thức.
\shortans{2}
\loigiai{
Đối chiếu từng phát biểu với kiến thức cốt lõi: Phân hạch hạt nhân nặng và nhiệt hạch hạt nhân nhẹ đều có thể tỏa năng lượng do sản phẩm bền hơn về năng lượng liên kết riêng. Có 2 phát biểu đúng.
}
\end{ex}

% ===== Câu 162 =====
% ID: L12C4B22-41-TN
% Mức: VD
\dangbt{Hoàn thành phương trình phản ứng hạt nhân}
\begin{ex}
% ID: L12C4B22-41-TN
% Mức: VD
Kết luận đúng khi vận dụng là về hoàn thành phương trình phản ứng hạt nhân?
\choice
{Chỉ cần bảo toàn số khối, không cần bảo toàn điện tích.}
{\True Muốn tìm hạt còn thiếu phải bảo toàn tổng số khối $A$ và tổng điện tích $Z$ ở hai vế.}
{Có thể bỏ qua điều kiện áp dụng và đơn vị trong mọi trường hợp.}
{Kết quả của bài toán không phụ thuộc dữ kiện đã cho.}
\loigiai{
Muốn tìm hạt còn thiếu phải bảo toàn tổng số khối $A$ và tổng điện tích $Z$ ở hai vế. Vì vậy phương án $B$ đúng; các phương án còn lại trái với kiến thức hoặc điều kiện áp dụng.
}
\end{ex}

% ===== Câu 163 =====
% ID: L12C4B22-42-TL
% Mức: VDC
\dangbt{Bài toán tổng hợp về phản ứng hạt nhân, phân hạch và nhiệt hạch}
\begin{bt}
% ID: L12C4B22-42-TL
% Mức: VDC
So sánh nhận định đúng “Phân hạch hạt nhân nặng và nhiệt hạch hạt nhân nhẹ đều có thể tỏa năng lượng do sản phẩm bền hơn về năng lượng liên kết riêng.” với nhận định sai “Mọi phản ứng hạt nhân đều tự phát và tỏa cùng một năng lượng.”; chỉ rõ căn cứ Vật lí.
\loigiai{
Cần nêu được: Phân hạch hạt nhân nặng và nhiệt hạch hạt nhân nhẹ đều có thể tỏa năng lượng do sản phẩm bền hơn về năng lượng liên kết riêng. Khi vận dụng phải xác định đúng dữ kiện, điều kiện áp dụng, hệ đơn vị và kiểm tra tính hợp lí của kết quả. Nhận định “Mọi phản ứng hạt nhân đều tự phát và tỏa cùng một năng lượng.” sai vì trái với kiến thức trên.
}
\end{bt}

% ===== Câu 164 =====
% ID: L12C4B22-42-TLN
% Mức: VDC
\begin{ex}
% ID: L12C4B22-42-TLN
% Mức: VDC
Trong bốn phát biểu sau về bài toán tổng hợp về phản ứng hạt nhân, phân hạch và nhiệt hạch, có bao nhiêu phát biểu đúng? Chỉ ghi một số nguyên. (1) Kết quả cần có ý nghĩa vật lí. (2) Mọi phản ứng hạt nhân đều tự phát và tỏa cùng một năng lượng. (3) Phải dùng đúng định luật cho tình huống. (4) Phân hạch hạt nhân nặng và nhiệt hạch hạt nhân nhẹ đều có thể tỏa năng lượng do sản phẩm bền hơn về năng lượng liên kết riêng.
\shortans{3}
\loigiai{
Đối chiếu từng phát biểu với kiến thức cốt lõi: Phân hạch hạt nhân nặng và nhiệt hạch hạt nhân nhẹ đều có thể tỏa năng lượng do sản phẩm bền hơn về năng lượng liên kết riêng. Có 3 phát biểu đúng.
}
\end{ex}

% ===== Câu 165 =====
% ID: L12C4B22-42-TN
% Mức: NB
\dangbt{Tính độ hụt khối, năng lượng liên kết và năng lượng liên kết riêng}
\begin{ex}
% ID: L12C4B22-42-TN
% Mức: NB
Phát biểu nào sau đây đúng về tính độ hụt khối, năng lượng liên kết và năng lượng liên kết riêng?
\choice
{\True Năng lượng liên kết $W_{lk}=\Delta m c^2$; năng lượng liên kết riêng bằng $W_{lk}/A$.}
{Năng lượng liên kết riêng bằng năng lượng liên kết nhân với số khối.}
{Có thể bỏ qua điều kiện áp dụng và đơn vị trong mọi trường hợp.}
{Kết quả của bài toán không phụ thuộc dữ kiện đã cho.}
\loigiai{
Năng lượng liên kết $W_{lk}=\Delta m c^2$; năng lượng liên kết riêng bằng $W_{lk}/A$. Vì vậy phương án $A$ đúng; các phương án còn lại trái với kiến thức hoặc điều kiện áp dụng.
}
\end{ex}

% ===== Câu 166 =====
% ID: L12C4B22-43-TN
% Mức: NB
\begin{ex}
% ID: L12C4B22-43-TN
% Mức: NB
Nhận định nào phù hợp với kiến thức Vật lí về tính độ hụt khối, năng lượng liên kết và năng lượng liên kết riêng?
\choice
{Năng lượng liên kết riêng bằng năng lượng liên kết nhân với số khối.}
{\True Năng lượng liên kết $W_{lk}=\Delta m c^2$; năng lượng liên kết riêng bằng $W_{lk}/A$.}
{Có thể bỏ qua điều kiện áp dụng và đơn vị trong mọi trường hợp.}
{Kết quả của bài toán không phụ thuộc dữ kiện đã cho.}
\loigiai{
Năng lượng liên kết $W_{lk}=\Delta m c^2$; năng lượng liên kết riêng bằng $W_{lk}/A$. Vì vậy phương án $B$ đúng; các phương án còn lại trái với kiến thức hoặc điều kiện áp dụng.
}
\end{ex}

% ===== Câu 167 =====
% ID: L12C4B22-44-TN
% Mức: NB
\begin{ex}
% ID: L12C4B22-44-TN
% Mức: NB
Chọn kết luận chính xác về tính độ hụt khối, năng lượng liên kết và năng lượng liên kết riêng?
\choice
{Năng lượng liên kết riêng bằng năng lượng liên kết nhân với số khối.}
{Có thể bỏ qua điều kiện áp dụng và đơn vị trong mọi trường hợp.}
{\True Năng lượng liên kết $W_{lk}=\Delta m c^2$; năng lượng liên kết riêng bằng $W_{lk}/A$.}
{Kết quả của bài toán không phụ thuộc dữ kiện đã cho.}
\loigiai{
Năng lượng liên kết $W_{lk}=\Delta m c^2$; năng lượng liên kết riêng bằng $W_{lk}/A$. Vì vậy phương án $C$ đúng; các phương án còn lại trái với kiến thức hoặc điều kiện áp dụng.
}
\end{ex}

% ===== Câu 168 =====
% ID: L12C4B22-45-TN
% Mức: TH
\begin{ex}
% ID: L12C4B22-45-TN
% Mức: TH
Nội dung nào mô tả đúng nhất về tính độ hụt khối, năng lượng liên kết và năng lượng liên kết riêng?
\choice
{Năng lượng liên kết riêng bằng năng lượng liên kết nhân với số khối.}
{Có thể bỏ qua điều kiện áp dụng và đơn vị trong mọi trường hợp.}
{Kết quả của bài toán không phụ thuộc dữ kiện đã cho.}
{\True Năng lượng liên kết $W_{lk}=\Delta m c^2$; năng lượng liên kết riêng bằng $W_{lk}/A$.}
\loigiai{
Năng lượng liên kết $W_{lk}=\Delta m c^2$; năng lượng liên kết riêng bằng $W_{lk}/A$. Vì vậy phương án $D$ đúng; các phương án còn lại trái với kiến thức hoặc điều kiện áp dụng.
}
\end{ex}

% ===== Câu 169 =====
% ID: L12C4B22-46-TN
% Mức: TH
\begin{ex}
% ID: L12C4B22-46-TN
% Mức: TH
Khi phân tích, phát biểu nào đúng về tính độ hụt khối, năng lượng liên kết và năng lượng liên kết riêng?
\choice
{\True Năng lượng liên kết $W_{lk}=\Delta m c^2$; năng lượng liên kết riêng bằng $W_{lk}/A$.}
{Năng lượng liên kết riêng bằng năng lượng liên kết nhân với số khối.}
{Có thể bỏ qua điều kiện áp dụng và đơn vị trong mọi trường hợp.}
{Kết quả của bài toán không phụ thuộc dữ kiện đã cho.}
\loigiai{
Năng lượng liên kết $W_{lk}=\Delta m c^2$; năng lượng liên kết riêng bằng $W_{lk}/A$. Vì vậy phương án $A$ đúng; các phương án còn lại trái với kiến thức hoặc điều kiện áp dụng.
}
\end{ex}

% ===== Câu 170 =====
% ID: L12C4B22-47-TN
% Mức: TH
\begin{ex}
% ID: L12C4B22-47-TN
% Mức: TH
Kết luận nào của học sinh là đúng về tính độ hụt khối, năng lượng liên kết và năng lượng liên kết riêng?
\choice
{Năng lượng liên kết riêng bằng năng lượng liên kết nhân với số khối.}
{\True Năng lượng liên kết $W_{lk}=\Delta m c^2$; năng lượng liên kết riêng bằng $W_{lk}/A$.}
{Có thể bỏ qua điều kiện áp dụng và đơn vị trong mọi trường hợp.}
{Kết quả của bài toán không phụ thuộc dữ kiện đã cho.}
\loigiai{
Năng lượng liên kết $W_{lk}=\Delta m c^2$; năng lượng liên kết riêng bằng $W_{lk}/A$. Vì vậy phương án $B$ đúng; các phương án còn lại trái với kiến thức hoặc điều kiện áp dụng.
}
\end{ex}

% ===== Câu 171 =====
% ID: L12C4B22-48-TN
% Mức: TH
\begin{ex}
% ID: L12C4B22-48-TN
% Mức: TH
Chọn phát biểu không trái với kiến thức về tính độ hụt khối, năng lượng liên kết và năng lượng liên kết riêng?
\choice
{Năng lượng liên kết riêng bằng năng lượng liên kết nhân với số khối.}
{Có thể bỏ qua điều kiện áp dụng và đơn vị trong mọi trường hợp.}
{\True Năng lượng liên kết $W_{lk}=\Delta m c^2$; năng lượng liên kết riêng bằng $W_{lk}/A$.}
{Kết quả của bài toán không phụ thuộc dữ kiện đã cho.}
\loigiai{
Năng lượng liên kết $W_{lk}=\Delta m c^2$; năng lượng liên kết riêng bằng $W_{lk}/A$. Vì vậy phương án $C$ đúng; các phương án còn lại trái với kiến thức hoặc điều kiện áp dụng.
}
\end{ex}

% ===== Câu 172 =====
% ID: L12C4B22-49-TN
% Mức: VD
\begin{ex}
% ID: L12C4B22-49-TN
% Mức: VD
Cơ sở Vật lí đúng để giải bài là gì về tính độ hụt khối, năng lượng liên kết và năng lượng liên kết riêng?
\choice
{Năng lượng liên kết riêng bằng năng lượng liên kết nhân với số khối.}
{Có thể bỏ qua điều kiện áp dụng và đơn vị trong mọi trường hợp.}
{Kết quả của bài toán không phụ thuộc dữ kiện đã cho.}
{\True Năng lượng liên kết $W_{lk}=\Delta m c^2$; năng lượng liên kết riêng bằng $W_{lk}/A$.}
\loigiai{
Năng lượng liên kết $W_{lk}=\Delta m c^2$; năng lượng liên kết riêng bằng $W_{lk}/A$. Vì vậy phương án $D$ đúng; các phương án còn lại trái với kiến thức hoặc điều kiện áp dụng.
}
\end{ex}

% ===== Câu 173 =====
% ID: L12C4B22-50-TN
% Mức: VD
\begin{ex}
% ID: L12C4B22-50-TN
% Mức: VD
Trong một bài thực hành, nhận xét nào đúng về tính độ hụt khối, năng lượng liên kết và năng lượng liên kết riêng?
\choice
{\True Năng lượng liên kết $W_{lk}=\Delta m c^2$; năng lượng liên kết riêng bằng $W_{lk}/A$.}
{Năng lượng liên kết riêng bằng năng lượng liên kết nhân với số khối.}
{Có thể bỏ qua điều kiện áp dụng và đơn vị trong mọi trường hợp.}
{Kết quả của bài toán không phụ thuộc dữ kiện đã cho.}
\loigiai{
Năng lượng liên kết $W_{lk}=\Delta m c^2$; năng lượng liên kết riêng bằng $W_{lk}/A$. Vì vậy phương án $A$ đúng; các phương án còn lại trái với kiến thức hoặc điều kiện áp dụng.
}
\end{ex}

% ===== Câu 174 =====
% ID: L12C4B22-51-TN
% Mức: VD
\begin{ex}
% ID: L12C4B22-51-TN
% Mức: VD
Kết luận đúng khi vận dụng là về tính độ hụt khối, năng lượng liên kết và năng lượng liên kết riêng?
\choice
{Năng lượng liên kết riêng bằng năng lượng liên kết nhân với số khối.}
{\True Năng lượng liên kết $W_{lk}=\Delta m c^2$; năng lượng liên kết riêng bằng $W_{lk}/A$.}
{Có thể bỏ qua điều kiện áp dụng và đơn vị trong mọi trường hợp.}
{Kết quả của bài toán không phụ thuộc dữ kiện đã cho.}
\loigiai{
Năng lượng liên kết $W_{lk}=\Delta m c^2$; năng lượng liên kết riêng bằng $W_{lk}/A$. Vì vậy phương án $B$ đúng; các phương án còn lại trái với kiến thức hoặc điều kiện áp dụng.
}
\end{ex}

% ===== Câu 175 =====
% ID: L12C4B22-52-TN
% Mức: NB
\dangbt{So sánh độ bền vững của các hạt nhân}
\begin{ex}
% ID: L12C4B22-52-TN
% Mức: NB
Phát biểu nào sau đây đúng về so sánh độ bền vững của các hạt nhân?
\choice
{\True Trong phạm vi so sánh, hạt nhân có năng lượng liên kết riêng lớn hơn thường bền vững hơn.}
{Hạt nhân có số khối lớn hơn luôn bền vững hơn.}
{Có thể bỏ qua điều kiện áp dụng và đơn vị trong mọi trường hợp.}
{Kết quả của bài toán không phụ thuộc dữ kiện đã cho.}
\loigiai{
Trong phạm vi so sánh, hạt nhân có năng lượng liên kết riêng lớn hơn thường bền vững hơn. Vì vậy phương án $A$ đúng; các phương án còn lại trái với kiến thức hoặc điều kiện áp dụng.
}
\end{ex}

% ===== Câu 176 =====
% ID: L12C4B22-53-TN
% Mức: NB
\begin{ex}
% ID: L12C4B22-53-TN
% Mức: NB
Nhận định nào phù hợp với kiến thức Vật lí về so sánh độ bền vững của các hạt nhân?
\choice
{Hạt nhân có số khối lớn hơn luôn bền vững hơn.}
{\True Trong phạm vi so sánh, hạt nhân có năng lượng liên kết riêng lớn hơn thường bền vững hơn.}
{Có thể bỏ qua điều kiện áp dụng và đơn vị trong mọi trường hợp.}
{Kết quả của bài toán không phụ thuộc dữ kiện đã cho.}
\loigiai{
Trong phạm vi so sánh, hạt nhân có năng lượng liên kết riêng lớn hơn thường bền vững hơn. Vì vậy phương án $B$ đúng; các phương án còn lại trái với kiến thức hoặc điều kiện áp dụng.
}
\end{ex}

% ===== Câu 177 =====
% ID: L12C4B22-54-TN
% Mức: NB
\begin{ex}
% ID: L12C4B22-54-TN
% Mức: NB
Chọn kết luận chính xác về so sánh độ bền vững của các hạt nhân?
\choice
{Hạt nhân có số khối lớn hơn luôn bền vững hơn.}
{Có thể bỏ qua điều kiện áp dụng và đơn vị trong mọi trường hợp.}
{\True Trong phạm vi so sánh, hạt nhân có năng lượng liên kết riêng lớn hơn thường bền vững hơn.}
{Kết quả của bài toán không phụ thuộc dữ kiện đã cho.}
\loigiai{
Trong phạm vi so sánh, hạt nhân có năng lượng liên kết riêng lớn hơn thường bền vững hơn. Vì vậy phương án $C$ đúng; các phương án còn lại trái với kiến thức hoặc điều kiện áp dụng.
}
\end{ex}

% ===== Câu 178 =====
% ID: L12C4B22-55-TN
% Mức: TH
\begin{ex}
% ID: L12C4B22-55-TN
% Mức: TH
Nội dung nào mô tả đúng nhất về so sánh độ bền vững của các hạt nhân?
\choice
{Hạt nhân có số khối lớn hơn luôn bền vững hơn.}
{Có thể bỏ qua điều kiện áp dụng và đơn vị trong mọi trường hợp.}
{Kết quả của bài toán không phụ thuộc dữ kiện đã cho.}
{\True Trong phạm vi so sánh, hạt nhân có năng lượng liên kết riêng lớn hơn thường bền vững hơn.}
\loigiai{
Trong phạm vi so sánh, hạt nhân có năng lượng liên kết riêng lớn hơn thường bền vững hơn. Vì vậy phương án $D$ đúng; các phương án còn lại trái với kiến thức hoặc điều kiện áp dụng.
}
\end{ex}

% ===== Câu 179 =====
% ID: L12C4B22-56-TN
% Mức: TH
\begin{ex}
% ID: L12C4B22-56-TN
% Mức: TH
Khi phân tích, phát biểu nào đúng về so sánh độ bền vững của các hạt nhân?
\choice
{\True Trong phạm vi so sánh, hạt nhân có năng lượng liên kết riêng lớn hơn thường bền vững hơn.}
{Hạt nhân có số khối lớn hơn luôn bền vững hơn.}
{Có thể bỏ qua điều kiện áp dụng và đơn vị trong mọi trường hợp.}
{Kết quả của bài toán không phụ thuộc dữ kiện đã cho.}
\loigiai{
Trong phạm vi so sánh, hạt nhân có năng lượng liên kết riêng lớn hơn thường bền vững hơn. Vì vậy phương án $A$ đúng; các phương án còn lại trái với kiến thức hoặc điều kiện áp dụng.
}
\end{ex}

% ===== Câu 180 =====
% ID: L12C4B22-57-TN
% Mức: TH
\begin{ex}
% ID: L12C4B22-57-TN
% Mức: TH
Kết luận nào của học sinh là đúng về so sánh độ bền vững của các hạt nhân?
\choice
{Hạt nhân có số khối lớn hơn luôn bền vững hơn.}
{\True Trong phạm vi so sánh, hạt nhân có năng lượng liên kết riêng lớn hơn thường bền vững hơn.}
{Có thể bỏ qua điều kiện áp dụng và đơn vị trong mọi trường hợp.}
{Kết quả của bài toán không phụ thuộc dữ kiện đã cho.}
\loigiai{
Trong phạm vi so sánh, hạt nhân có năng lượng liên kết riêng lớn hơn thường bền vững hơn. Vì vậy phương án $B$ đúng; các phương án còn lại trái với kiến thức hoặc điều kiện áp dụng.
}
\end{ex}

% ===== Câu 181 =====
% ID: L12C4B22-58-TN
% Mức: TH
\begin{ex}
% ID: L12C4B22-58-TN
% Mức: TH
Chọn phát biểu không trái với kiến thức về so sánh độ bền vững của các hạt nhân?
\choice
{Hạt nhân có số khối lớn hơn luôn bền vững hơn.}
{Có thể bỏ qua điều kiện áp dụng và đơn vị trong mọi trường hợp.}
{\True Trong phạm vi so sánh, hạt nhân có năng lượng liên kết riêng lớn hơn thường bền vững hơn.}
{Kết quả của bài toán không phụ thuộc dữ kiện đã cho.}
\loigiai{
Trong phạm vi so sánh, hạt nhân có năng lượng liên kết riêng lớn hơn thường bền vững hơn. Vì vậy phương án $C$ đúng; các phương án còn lại trái với kiến thức hoặc điều kiện áp dụng.
}
\end{ex}

% ===== Câu 182 =====
% ID: L12C4B22-59-TN
% Mức: VD
\begin{ex}
% ID: L12C4B22-59-TN
% Mức: VD
Cơ sở Vật lí đúng để giải bài là gì về so sánh độ bền vững của các hạt nhân?
\choice
{Hạt nhân có số khối lớn hơn luôn bền vững hơn.}
{Có thể bỏ qua điều kiện áp dụng và đơn vị trong mọi trường hợp.}
{Kết quả của bài toán không phụ thuộc dữ kiện đã cho.}
{\True Trong phạm vi so sánh, hạt nhân có năng lượng liên kết riêng lớn hơn thường bền vững hơn.}
\loigiai{
Trong phạm vi so sánh, hạt nhân có năng lượng liên kết riêng lớn hơn thường bền vững hơn. Vì vậy phương án $D$ đúng; các phương án còn lại trái với kiến thức hoặc điều kiện áp dụng.
}
\end{ex}

% ===== Câu 183 =====
% ID: L12C4B22-60-TN
% Mức: VD
\begin{ex}
% ID: L12C4B22-60-TN
% Mức: VD
Trong một bài thực hành, nhận xét nào đúng về so sánh độ bền vững của các hạt nhân?
\choice
{\True Trong phạm vi so sánh, hạt nhân có năng lượng liên kết riêng lớn hơn thường bền vững hơn.}
{Hạt nhân có số khối lớn hơn luôn bền vững hơn.}
{Có thể bỏ qua điều kiện áp dụng và đơn vị trong mọi trường hợp.}
{Kết quả của bài toán không phụ thuộc dữ kiện đã cho.}
\loigiai{
Trong phạm vi so sánh, hạt nhân có năng lượng liên kết riêng lớn hơn thường bền vững hơn. Vì vậy phương án $A$ đúng; các phương án còn lại trái với kiến thức hoặc điều kiện áp dụng.
}
\end{ex}

% ===== Câu 184 =====
% ID: L12C4B22-61-TN
% Mức: VD
\begin{ex}
% ID: L12C4B22-61-TN
% Mức: VD
Kết luận đúng khi vận dụng là về so sánh độ bền vững của các hạt nhân?
\choice
{Hạt nhân có số khối lớn hơn luôn bền vững hơn.}
{\True Trong phạm vi so sánh, hạt nhân có năng lượng liên kết riêng lớn hơn thường bền vững hơn.}
{Có thể bỏ qua điều kiện áp dụng và đơn vị trong mọi trường hợp.}
{Kết quả của bài toán không phụ thuộc dữ kiện đã cho.}
\loigiai{
Trong phạm vi so sánh, hạt nhân có năng lượng liên kết riêng lớn hơn thường bền vững hơn. Vì vậy phương án $B$ đúng; các phương án còn lại trái với kiến thức hoặc điều kiện áp dụng.
}
\end{ex}

% ===== Câu 185 =====
% ID: L12C4B22-62-TN
% Mức: NB
\dangbt{Tính năng lượng tỏa ra hoặc thu vào trong phản ứng hạt nhân}
\begin{ex}
% ID: L12C4B22-62-TN
% Mức: NB
Phát biểu nào sau đây đúng về tính năng lượng tỏa ra hoặc thu vào trong phản ứng hạt nhân?
\choice
{\True Năng lượng phản ứng $Q=(m_{trước}-m_{sau})c^2$; $Q>0$ là tỏa năng lượng, $Q<0$ là thu năng lượng.}
{Phản ứng có $m_{sau}\gt m_{trước}$ luôn tỏa năng lượng.}
{Có thể bỏ qua điều kiện áp dụng và đơn vị trong mọi trường hợp.}
{Kết quả của bài toán không phụ thuộc dữ kiện đã cho.}
\loigiai{
Năng lượng phản ứng $Q=(m_{trước}-m_{sau})c^2$; $Q>0$ là tỏa năng lượng, $Q<0$ là thu năng lượng. Vì vậy phương án $A$ đúng; các phương án còn lại trái với kiến thức hoặc điều kiện áp dụng.
}
\end{ex}

% ===== Câu 186 =====
% ID: L12C4B22-63-TN
% Mức: NB
\begin{ex}
% ID: L12C4B22-63-TN
% Mức: NB
Nhận định nào phù hợp với kiến thức Vật lí về tính năng lượng tỏa ra hoặc thu vào trong phản ứng hạt nhân?
\choice
{Phản ứng có $m_{sau}\gt m_{trước}$ luôn tỏa năng lượng.}
{\True Năng lượng phản ứng $Q=(m_{trước}-m_{sau})c^2$; $Q>0$ là tỏa năng lượng, $Q<0$ là thu năng lượng.}
{Có thể bỏ qua điều kiện áp dụng và đơn vị trong mọi trường hợp.}
{Kết quả của bài toán không phụ thuộc dữ kiện đã cho.}
\loigiai{
Năng lượng phản ứng $Q=(m_{trước}-m_{sau})c^2$; $Q>0$ là tỏa năng lượng, $Q<0$ là thu năng lượng. Vì vậy phương án $B$ đúng; các phương án còn lại trái với kiến thức hoặc điều kiện áp dụng.
}
\end{ex}

% ===== Câu 187 =====
% ID: L12C4B22-64-TN
% Mức: NB
\begin{ex}
% ID: L12C4B22-64-TN
% Mức: NB
Chọn kết luận chính xác về tính năng lượng tỏa ra hoặc thu vào trong phản ứng hạt nhân?
\choice
{Phản ứng có $m_{sau}\gt m_{trước}$ luôn tỏa năng lượng.}
{Có thể bỏ qua điều kiện áp dụng và đơn vị trong mọi trường hợp.}
{\True Năng lượng phản ứng $Q=(m_{trước}-m_{sau})c^2$; $Q>0$ là tỏa năng lượng, $Q<0$ là thu năng lượng.}
{Kết quả của bài toán không phụ thuộc dữ kiện đã cho.}
\loigiai{
Năng lượng phản ứng $Q=(m_{trước}-m_{sau})c^2$; $Q>0$ là tỏa năng lượng, $Q<0$ là thu năng lượng. Vì vậy phương án $C$ đúng; các phương án còn lại trái với kiến thức hoặc điều kiện áp dụng.
}
\end{ex}

% ===== Câu 188 =====
% ID: L12C4B22-65-TN
% Mức: TH
\begin{ex}
% ID: L12C4B22-65-TN
% Mức: TH
Nội dung nào mô tả đúng nhất về tính năng lượng tỏa ra hoặc thu vào trong phản ứng hạt nhân?
\choice
{Phản ứng có $m_{sau}\gt m_{trước}$ luôn tỏa năng lượng.}
{Có thể bỏ qua điều kiện áp dụng và đơn vị trong mọi trường hợp.}
{Kết quả của bài toán không phụ thuộc dữ kiện đã cho.}
{\True Năng lượng phản ứng $Q=(m_{trước}-m_{sau})c^2$; $Q>0$ là tỏa năng lượng, $Q<0$ là thu năng lượng.}
\loigiai{
Năng lượng phản ứng $Q=(m_{trước}-m_{sau})c^2$; $Q>0$ là tỏa năng lượng, $Q<0$ là thu năng lượng. Vì vậy phương án $D$ đúng; các phương án còn lại trái với kiến thức hoặc điều kiện áp dụng.
}
\end{ex}

% ===== Câu 189 =====
% ID: L12C4B22-66-TN
% Mức: TH
\begin{ex}
% ID: L12C4B22-66-TN
% Mức: TH
Khi phân tích, phát biểu nào đúng về tính năng lượng tỏa ra hoặc thu vào trong phản ứng hạt nhân?
\choice
{\True Năng lượng phản ứng $Q=(m_{trước}-m_{sau})c^2$; $Q>0$ là tỏa năng lượng, $Q<0$ là thu năng lượng.}
{Phản ứng có $m_{sau}\gt m_{trước}$ luôn tỏa năng lượng.}
{Có thể bỏ qua điều kiện áp dụng và đơn vị trong mọi trường hợp.}
{Kết quả của bài toán không phụ thuộc dữ kiện đã cho.}
\loigiai{
Năng lượng phản ứng $Q=(m_{trước}-m_{sau})c^2$; $Q>0$ là tỏa năng lượng, $Q<0$ là thu năng lượng. Vì vậy phương án $A$ đúng; các phương án còn lại trái với kiến thức hoặc điều kiện áp dụng.
}
\end{ex}

% ===== Câu 190 =====
% ID: L12C4B22-67-TN
% Mức: TH
\begin{ex}
% ID: L12C4B22-67-TN
% Mức: TH
Kết luận nào của học sinh là đúng về tính năng lượng tỏa ra hoặc thu vào trong phản ứng hạt nhân?
\choice
{Phản ứng có $m_{sau}\gt m_{trước}$ luôn tỏa năng lượng.}
{\True Năng lượng phản ứng $Q=(m_{trước}-m_{sau})c^2$; $Q>0$ là tỏa năng lượng, $Q<0$ là thu năng lượng.}
{Có thể bỏ qua điều kiện áp dụng và đơn vị trong mọi trường hợp.}
{Kết quả của bài toán không phụ thuộc dữ kiện đã cho.}
\loigiai{
Năng lượng phản ứng $Q=(m_{trước}-m_{sau})c^2$; $Q>0$ là tỏa năng lượng, $Q<0$ là thu năng lượng. Vì vậy phương án $B$ đúng; các phương án còn lại trái với kiến thức hoặc điều kiện áp dụng.
}
\end{ex}

% ===== Câu 191 =====
% ID: L12C4B22-68-TN
% Mức: TH
\begin{ex}
% ID: L12C4B22-68-TN
% Mức: TH
Chọn phát biểu không trái với kiến thức về tính năng lượng tỏa ra hoặc thu vào trong phản ứng hạt nhân?
\choice
{Phản ứng có $m_{sau}\gt m_{trước}$ luôn tỏa năng lượng.}
{Có thể bỏ qua điều kiện áp dụng và đơn vị trong mọi trường hợp.}
{\True Năng lượng phản ứng $Q=(m_{trước}-m_{sau})c^2$; $Q>0$ là tỏa năng lượng, $Q<0$ là thu năng lượng.}
{Kết quả của bài toán không phụ thuộc dữ kiện đã cho.}
\loigiai{
Năng lượng phản ứng $Q=(m_{trước}-m_{sau})c^2$; $Q>0$ là tỏa năng lượng, $Q<0$ là thu năng lượng. Vì vậy phương án $C$ đúng; các phương án còn lại trái với kiến thức hoặc điều kiện áp dụng.
}
\end{ex}

% ===== Câu 192 =====
% ID: L12C4B22-69-TN
% Mức: VD
\begin{ex}
% ID: L12C4B22-69-TN
% Mức: VD
Cơ sở Vật lí đúng để giải bài là gì về tính năng lượng tỏa ra hoặc thu vào trong phản ứng hạt nhân?
\choice
{Phản ứng có $m_{sau}\gt m_{trước}$ luôn tỏa năng lượng.}
{Có thể bỏ qua điều kiện áp dụng và đơn vị trong mọi trường hợp.}
{Kết quả của bài toán không phụ thuộc dữ kiện đã cho.}
{\True Năng lượng phản ứng $Q=(m_{trước}-m_{sau})c^2$; $Q>0$ là tỏa năng lượng, $Q<0$ là thu năng lượng.}
\loigiai{
Năng lượng phản ứng $Q=(m_{trước}-m_{sau})c^2$; $Q>0$ là tỏa năng lượng, $Q<0$ là thu năng lượng. Vì vậy phương án $D$ đúng; các phương án còn lại trái với kiến thức hoặc điều kiện áp dụng.
}
\end{ex}

% ===== Câu 193 =====
% ID: L12C4B22-70-TN
% Mức: VD
\begin{ex}
% ID: L12C4B22-70-TN
% Mức: VD
Trong một bài thực hành, nhận xét nào đúng về tính năng lượng tỏa ra hoặc thu vào trong phản ứng hạt nhân?
\choice
{\True Năng lượng phản ứng $Q=(m_{trước}-m_{sau})c^2$; $Q>0$ là tỏa năng lượng, $Q<0$ là thu năng lượng.}
{Phản ứng có $m_{sau}\gt m_{trước}$ luôn tỏa năng lượng.}
{Có thể bỏ qua điều kiện áp dụng và đơn vị trong mọi trường hợp.}
{Kết quả của bài toán không phụ thuộc dữ kiện đã cho.}
\loigiai{
Năng lượng phản ứng $Q=(m_{trước}-m_{sau})c^2$; $Q>0$ là tỏa năng lượng, $Q<0$ là thu năng lượng. Vì vậy phương án $A$ đúng; các phương án còn lại trái với kiến thức hoặc điều kiện áp dụng.
}
\end{ex}

% ===== Câu 194 =====
% ID: L12C4B22-71-TN
% Mức: VD
\begin{ex}
% ID: L12C4B22-71-TN
% Mức: VD
Kết luận đúng khi vận dụng là về tính năng lượng tỏa ra hoặc thu vào trong phản ứng hạt nhân?
\choice
{Phản ứng có $m_{sau}\gt m_{trước}$ luôn tỏa năng lượng.}
{\True Năng lượng phản ứng $Q=(m_{trước}-m_{sau})c^2$; $Q>0$ là tỏa năng lượng, $Q<0$ là thu năng lượng.}
{Có thể bỏ qua điều kiện áp dụng và đơn vị trong mọi trường hợp.}
{Kết quả của bài toán không phụ thuộc dữ kiện đã cho.}
\loigiai{
Năng lượng phản ứng $Q=(m_{trước}-m_{sau})c^2$; $Q>0$ là tỏa năng lượng, $Q<0$ là thu năng lượng. Vì vậy phương án $B$ đúng; các phương án còn lại trái với kiến thức hoặc điều kiện áp dụng.
}
\end{ex}

% ===== Câu 195 =====
% ID: L12C4B22-72-TN
% Mức: NB
\dangbt{Bài toán tổng hợp về phản ứng hạt nhân, phân hạch và nhiệt hạch}
\begin{ex}
% ID: L12C4B22-72-TN
% Mức: NB
Phát biểu nào sau đây đúng về bài toán tổng hợp về phản ứng hạt nhân, phân hạch và nhiệt hạch?
\choice
{\True Phân hạch hạt nhân nặng và nhiệt hạch hạt nhân nhẹ đều có thể tỏa năng lượng do sản phẩm bền hơn về năng lượng liên kết riêng.}
{Mọi phản ứng hạt nhân đều tự phát và tỏa cùng một năng lượng.}
{Có thể bỏ qua điều kiện áp dụng và đơn vị trong mọi trường hợp.}
{Kết quả của bài toán không phụ thuộc dữ kiện đã cho.}
\loigiai{
Phân hạch hạt nhân nặng và nhiệt hạch hạt nhân nhẹ đều có thể tỏa năng lượng do sản phẩm bền hơn về năng lượng liên kết riêng. Vì vậy phương án $A$ đúng; các phương án còn lại trái với kiến thức hoặc điều kiện áp dụng.
}
\end{ex}

% ===== Câu 196 =====
% ID: L12C4B22-73-TN
% Mức: NB
\begin{ex}
% ID: L12C4B22-73-TN
% Mức: NB
Nhận định nào phù hợp với kiến thức Vật lí về bài toán tổng hợp về phản ứng hạt nhân, phân hạch và nhiệt hạch?
\choice
{Mọi phản ứng hạt nhân đều tự phát và tỏa cùng một năng lượng.}
{\True Phân hạch hạt nhân nặng và nhiệt hạch hạt nhân nhẹ đều có thể tỏa năng lượng do sản phẩm bền hơn về năng lượng liên kết riêng.}
{Có thể bỏ qua điều kiện áp dụng và đơn vị trong mọi trường hợp.}
{Kết quả của bài toán không phụ thuộc dữ kiện đã cho.}
\loigiai{
Phân hạch hạt nhân nặng và nhiệt hạch hạt nhân nhẹ đều có thể tỏa năng lượng do sản phẩm bền hơn về năng lượng liên kết riêng. Vì vậy phương án $B$ đúng; các phương án còn lại trái với kiến thức hoặc điều kiện áp dụng.
}
\end{ex}

% ===== Câu 197 =====
% ID: L12C4B22-74-TN
% Mức: NB
\begin{ex}
% ID: L12C4B22-74-TN
% Mức: NB
Chọn kết luận chính xác về bài toán tổng hợp về phản ứng hạt nhân, phân hạch và nhiệt hạch?
\choice
{Mọi phản ứng hạt nhân đều tự phát và tỏa cùng một năng lượng.}
{Có thể bỏ qua điều kiện áp dụng và đơn vị trong mọi trường hợp.}
{\True Phân hạch hạt nhân nặng và nhiệt hạch hạt nhân nhẹ đều có thể tỏa năng lượng do sản phẩm bền hơn về năng lượng liên kết riêng.}
{Kết quả của bài toán không phụ thuộc dữ kiện đã cho.}
\loigiai{
Phân hạch hạt nhân nặng và nhiệt hạch hạt nhân nhẹ đều có thể tỏa năng lượng do sản phẩm bền hơn về năng lượng liên kết riêng. Vì vậy phương án $C$ đúng; các phương án còn lại trái với kiến thức hoặc điều kiện áp dụng.
}
\end{ex}

% ===== Câu 198 =====
% ID: L12C4B22-75-TN
% Mức: TH
\begin{ex}
% ID: L12C4B22-75-TN
% Mức: TH
Nội dung nào mô tả đúng nhất về bài toán tổng hợp về phản ứng hạt nhân, phân hạch và nhiệt hạch?
\choice
{Mọi phản ứng hạt nhân đều tự phát và tỏa cùng một năng lượng.}
{Có thể bỏ qua điều kiện áp dụng và đơn vị trong mọi trường hợp.}
{Kết quả của bài toán không phụ thuộc dữ kiện đã cho.}
{\True Phân hạch hạt nhân nặng và nhiệt hạch hạt nhân nhẹ đều có thể tỏa năng lượng do sản phẩm bền hơn về năng lượng liên kết riêng.}
\loigiai{
Phân hạch hạt nhân nặng và nhiệt hạch hạt nhân nhẹ đều có thể tỏa năng lượng do sản phẩm bền hơn về năng lượng liên kết riêng. Vì vậy phương án $D$ đúng; các phương án còn lại trái với kiến thức hoặc điều kiện áp dụng.
}
\end{ex}

% ===== Câu 199 =====
% ID: L12C4B22-76-TN
% Mức: TH
\begin{ex}
% ID: L12C4B22-76-TN
% Mức: TH
Khi phân tích, phát biểu nào đúng về bài toán tổng hợp về phản ứng hạt nhân, phân hạch và nhiệt hạch?
\choice
{\True Phân hạch hạt nhân nặng và nhiệt hạch hạt nhân nhẹ đều có thể tỏa năng lượng do sản phẩm bền hơn về năng lượng liên kết riêng.}
{Mọi phản ứng hạt nhân đều tự phát và tỏa cùng một năng lượng.}
{Có thể bỏ qua điều kiện áp dụng và đơn vị trong mọi trường hợp.}
{Kết quả của bài toán không phụ thuộc dữ kiện đã cho.}
\loigiai{
Phân hạch hạt nhân nặng và nhiệt hạch hạt nhân nhẹ đều có thể tỏa năng lượng do sản phẩm bền hơn về năng lượng liên kết riêng. Vì vậy phương án $A$ đúng; các phương án còn lại trái với kiến thức hoặc điều kiện áp dụng.
}
\end{ex}

% ===== Câu 200 =====
% ID: L12C4B22-77-TN
% Mức: TH
\begin{ex}
% ID: L12C4B22-77-TN
% Mức: TH
Kết luận nào của học sinh là đúng về bài toán tổng hợp về phản ứng hạt nhân, phân hạch và nhiệt hạch?
\choice
{Mọi phản ứng hạt nhân đều tự phát và tỏa cùng một năng lượng.}
{\True Phân hạch hạt nhân nặng và nhiệt hạch hạt nhân nhẹ đều có thể tỏa năng lượng do sản phẩm bền hơn về năng lượng liên kết riêng.}
{Có thể bỏ qua điều kiện áp dụng và đơn vị trong mọi trường hợp.}
{Kết quả của bài toán không phụ thuộc dữ kiện đã cho.}
\loigiai{
Phân hạch hạt nhân nặng và nhiệt hạch hạt nhân nhẹ đều có thể tỏa năng lượng do sản phẩm bền hơn về năng lượng liên kết riêng. Vì vậy phương án $B$ đúng; các phương án còn lại trái với kiến thức hoặc điều kiện áp dụng.
}
\end{ex}

% ===== Câu 201 =====
% ID: L12C4B22-78-TN
% Mức: TH
\begin{ex}
% ID: L12C4B22-78-TN
% Mức: TH
Chọn phát biểu không trái với kiến thức về bài toán tổng hợp về phản ứng hạt nhân, phân hạch và nhiệt hạch?
\choice
{Mọi phản ứng hạt nhân đều tự phát và tỏa cùng một năng lượng.}
{Có thể bỏ qua điều kiện áp dụng và đơn vị trong mọi trường hợp.}
{\True Phân hạch hạt nhân nặng và nhiệt hạch hạt nhân nhẹ đều có thể tỏa năng lượng do sản phẩm bền hơn về năng lượng liên kết riêng.}
{Kết quả của bài toán không phụ thuộc dữ kiện đã cho.}
\loigiai{
Phân hạch hạt nhân nặng và nhiệt hạch hạt nhân nhẹ đều có thể tỏa năng lượng do sản phẩm bền hơn về năng lượng liên kết riêng. Vì vậy phương án $C$ đúng; các phương án còn lại trái với kiến thức hoặc điều kiện áp dụng.
}
\end{ex}

% ===== Câu 202 =====
% ID: L12C4B22-79-TN
% Mức: VD
\begin{ex}
% ID: L12C4B22-79-TN
% Mức: VD
Cơ sở Vật lí đúng để giải bài là gì về bài toán tổng hợp về phản ứng hạt nhân, phân hạch và nhiệt hạch?
\choice
{Mọi phản ứng hạt nhân đều tự phát và tỏa cùng một năng lượng.}
{Có thể bỏ qua điều kiện áp dụng và đơn vị trong mọi trường hợp.}
{Kết quả của bài toán không phụ thuộc dữ kiện đã cho.}
{\True Phân hạch hạt nhân nặng và nhiệt hạch hạt nhân nhẹ đều có thể tỏa năng lượng do sản phẩm bền hơn về năng lượng liên kết riêng.}
\loigiai{
Phân hạch hạt nhân nặng và nhiệt hạch hạt nhân nhẹ đều có thể tỏa năng lượng do sản phẩm bền hơn về năng lượng liên kết riêng. Vì vậy phương án $D$ đúng; các phương án còn lại trái với kiến thức hoặc điều kiện áp dụng.
}
\end{ex}

% ===== Câu 203 =====
% ID: L12C4B22-80-TN
% Mức: VD
\begin{ex}
% ID: L12C4B22-80-TN
% Mức: VD
Trong một bài thực hành, nhận xét nào đúng về bài toán tổng hợp về phản ứng hạt nhân, phân hạch và nhiệt hạch?
\choice
{\True Phân hạch hạt nhân nặng và nhiệt hạch hạt nhân nhẹ đều có thể tỏa năng lượng do sản phẩm bền hơn về năng lượng liên kết riêng.}
{Mọi phản ứng hạt nhân đều tự phát và tỏa cùng một năng lượng.}
{Có thể bỏ qua điều kiện áp dụng và đơn vị trong mọi trường hợp.}
{Kết quả của bài toán không phụ thuộc dữ kiện đã cho.}
\loigiai{
Phân hạch hạt nhân nặng và nhiệt hạch hạt nhân nhẹ đều có thể tỏa năng lượng do sản phẩm bền hơn về năng lượng liên kết riêng. Vì vậy phương án $A$ đúng; các phương án còn lại trái với kiến thức hoặc điều kiện áp dụng.
}
\end{ex}

% ===== Câu 204 =====
% ID: L12C4B22-81-TN
% Mức: VD
\begin{ex}
% ID: L12C4B22-81-TN
% Mức: VD
Kết luận đúng khi vận dụng là về bài toán tổng hợp về phản ứng hạt nhân, phân hạch và nhiệt hạch?
\choice
{Mọi phản ứng hạt nhân đều tự phát và tỏa cùng một năng lượng.}
{\True Phân hạch hạt nhân nặng và nhiệt hạch hạt nhân nhẹ đều có thể tỏa năng lượng do sản phẩm bền hơn về năng lượng liên kết riêng.}
{Có thể bỏ qua điều kiện áp dụng và đơn vị trong mọi trường hợp.}
{Kết quả của bài toán không phụ thuộc dữ kiện đã cho.}
\loigiai{
Phân hạch hạt nhân nặng và nhiệt hạch hạt nhân nhẹ đều có thể tỏa năng lượng do sản phẩm bền hơn về năng lượng liên kết riêng. Vì vậy phương án $B$ đúng; các phương án còn lại trái với kiến thức hoặc điều kiện áp dụng.
}
\end{ex}
